% Generated mechanically from frozen input inputs/p08/ja/crystalline.tex.
% Do not edit this generated file; edit the bound locale source instead.

% OK, start here.
%

\setcounter{chapter}{59}
\chapter{結晶コホモロジー}



\phantomsection
\label{crystalline-section-phantom}




\section{序論}
\label{crystalline-section-introduction}

\noindent
本章は、2012 年に Columbia University で Johan de Jong が行った
連続講義に基づく。本章の目的は、結晶コホモロジーを手短に導入することである。
参考文献として \cite{Berthelot} がある。

\medskip\noindent
分割冪環に関する、より初等的で純代数的な議論は準備章へ移した。
この議論は可換代数における Tate 分解を扱う際にも有用だからである。
『分割冪代数』第 \ref{dpa-section-introduction} 節を参照されたい。













\section{分割冪包絡}
\label{crystalline-section-divided-power-envelope}

\noindent
次の補題の構成を分割冪包絡と呼ぶ。これは後に重要な役割を果たす。

\begin{lemma}
\label{crystalline-lemma-divided-power-envelope}
$(A, I, \gamma)$ を分割冪環とする。
$A \to B$ を環準同型とし、$J \subset B$ を $IB \subset J$ を満たす
イデアルとする。このとき、分割冪環の準同型
$$
(A, I, \gamma) \longrightarrow (D, \bar J, \bar \gamma)
$$
で、任意の $(A, I, \gamma)$ 上の分割冪代数 $(C, K, \delta)$ に対して
関手的に
$$
\Hom_{(A, I, \gamma)}((D, \bar J, \bar \gamma), (C, K, \delta)) =
\Hom_{(A, I)}((B, J), (C, K))
$$
を満たすものが存在する。ここで左辺は $(A, I, \gamma)$ 上の分割冪環の
射の集合であり、右辺は $(A, I)$ 上の（環、イデアル）の組の射の集合である。
\end{lemma}

\begin{proof}
分割冪環 $(C, K, \delta)$ の圏を $\mathcal{C}$ と書く。次で定義される関手
$F : \mathcal{C} \longrightarrow \textit{Sets}$ を考える：
$$
F(C, K, \delta) =
\left\{
(\varphi, \psi)
\middle|
\begin{matrix}
\varphi : (A, I, \gamma) \to (C, K, \delta)
\text{ は分割冪環の準同型} \\
\psi : (B, J) \to (C, K)\text{ は }
A\text{-代数準同型であり、}\psi(J) \subset K
\end{matrix}
\right\}
$$
『分割冪代数』補題 \ref{dpa-lemma-a-version-of-brown} がこの関手に適用できる
ことを示せば、補題が従う。$(\varphi, \psi) \in F(C, K, \delta)$ と仮定する。
$C' \subset C$ を、$\varphi(A)$、$\psi(B)$、およびすべての $f \in J$ に対する
$\delta_n(\psi(f))$ で生成される部分環とする。$K' \subset K \cap C'$ を、
$\varphi(I)$ と $f \in J$ に対する $\delta_n(\psi(f))$ で生成される $C'$ の
イデアルとする。このとき $(C', K', \delta|_{K'})$ は分割冪環であり、$C'$ の
濃度は基数 $\kappa = |A| \otimes |B|^{\aleph_0}$ 以下である。さらに、
$\varphi$ は $A \to C' \to C$ と分解し、$\psi$ は $B \to C' \to C$ と分解する。
これにより『分割冪代数』補題 \ref{dpa-lemma-a-version-of-brown} の仮定 (1) が
成り立つ。仮定 (2) は明らかである。実際、分割冪環の圏における極限は忘却関手
$(C, K, \delta) \mapsto (C, K)$ と可換する。『分割冪代数』補題
\ref{dpa-lemma-limits} とその証明を参照されたい。
\end{proof}

\begin{definition}
\label{crystalline-definition-divided-power-envelope}
$(A, I, \gamma)$ を分割冪環とする。$A \to B$ を環準同型とし、
$J \subset B$ を $IB \subset J$ を満たすイデアルとする。補題
\ref{crystalline-lemma-divided-power-envelope} で構成した分割冪代数
$(D, \bar J, \bar\gamma)$ を、{\it $(A, I, \gamma)$ に相対的な $B$ における
$J$ の分割冪包絡}と呼び、$D_B(J)$ または $D_{B, \gamma}(J)$ と書く。
\end{definition}

\noindent
$(A, I, \gamma) \to (C, K, \delta)$ を分割冪環の準同型とする。
$D_{B, \gamma}(J) = (D, \bar J, \bar \gamma)$ の普遍性は
$$
\begin{matrix}
\text{環準同型 }B \to C \\
\text{そのもとで }J\text{ の像を含む }K
\end{matrix}
\longleftrightarrow
\begin{matrix}
\text{分割冪準同型} \\
(D, \bar J, \bar \gamma) \to (C, K, \delta)
\end{matrix}
$$
という対応で表される。この対応は $\text{id}_D$ に対応する写像 $B \to D$ との
前合成で与えられる。普遍性から直ちに従う $(D, \bar J, \bar \gamma)$ の性質を
いくつか挙げる。$A$-代数準同型
\begin{equation}
\label{crystalline-equation-divided-power-envelope}
B \longrightarrow D \longrightarrow B/J
\end{equation}
がある。第 1 の矢印は $J$ を $\bar J$ に写し、$\bar J$ は第 2 の矢印の核である。
$n > 0$ かつ $x$ が $J \to D$ の像に属するときの元 $\bar\gamma_n(x)$ は、
$D$ のイデアルとして $\bar J$ を生成し、$B$-代数として $D$ を生成する。

\begin{lemma}
\label{crystalline-lemma-divided-power-envelop-quotient}
$(A, I, \gamma)$ を分割冪環とする。$\varphi : B' \to B$ を、核が $K$ である
$A$-代数の全射とする。$IB \subset J \subset B$ をイデアルとし、
$J' \subset B'$ を $J$ の逆像とする。
$D_{B', \gamma}(J') = (D', \bar J', \bar\gamma)$.
と書く。このとき $D_{B, \gamma}(J) = (D'/K', \bar J'/K', \bar\gamma)$ である。
ここで $K'$ は、$n \geq 1$ および $k \in K$ に対する元
$\bar\gamma_n(k)$ で生成されるイデアルである。
\end{lemma}

\begin{proof}
$D_{B, \gamma}(J) = (D, \bar J, \bar \gamma)$ と書く。$D'$ の普遍性から、
分割冪代数の準同型 $D' \to D$ を得る。$B' \to B$ と $J' \to J$ は全射なので、
$D' \to D$ も全射である（上の注意を参照）。$n \geq 1$ かつ $k \in K$ に対して
$\bar\gamma_n(k)$ が核に属することは明らかであり、したがって準同型
$D'/K' \to D$ を得る。逆に、$\bar J'/K' \subset D'/K'$ 上には分割冪構造が
存在する。『分割冪代数』補題 \ref{dpa-lemma-kernel} を参照されたい。
よって $D$ の普遍性から逆写像 $D \to D'/K'$ が得られ、主張が従う。
\end{proof}

\noindent
定義 \ref{crystalline-definition-divided-power-envelope} の状況では、$A$ 上の多項式代数
$P$ からの全射 $P \to B$ を選び、$J' \subset P$ を $J$ の逆像とできる。
前の補題は $D_{B, \gamma}(J)$ を $D_{P, \gamma}(J')$ によって記述する。
『分割冪代数』補題 \ref{dpa-lemma-gamma-extends} により、$\gamma$ は $IP$ 上の
分割冪構造 $\gamma'$ に拡張することに注意する。したがって
$D_{P, \gamma}(J') = D_{P, \gamma'}(J')$ は、次の補題で記述する分割冪包絡の
特別な場合の一例である。

\begin{lemma}
\label{crystalline-lemma-describe-divided-power-envelope}
$(B, I, \gamma)$ を分割冪代数とする。$I \subset J \subset B$ をイデアルとし、
$(D, \bar J, \bar \gamma)$ を $\gamma$ に相対的な $J$ の分割冪包絡とする。
$J = I + (f_t)$ を満たす元 $f_t \in J$（$t \in T$）を選ぶ。このとき全射
$$
\Psi :
(B\langle x_t \rangle, IB\langle x_t \rangle + B\langle x_t \rangle_+, \delta)
\longrightarrow
(D, \bar J, \bar \gamma)
$$
で、$x_t$ を $D$ における $f_t$ の像へ写す分割冪環の準同型が存在する。
$\Psi$ の核は、元 $x_t - f_t$ と、$B$ においてある $r_t \in B$、$r_0 \in I$
に対し $\sum r_t f_t = r_0$ となる場合のすべての元
$$
\delta_n\left(\sum r_t x_t - r_0\right)
$$
によって生成される。
\end{lemma}

\begin{proof}
補題の主張において、$B\langle x_t \rangle$ はイデアル
$J' = IB\langle x_t \rangle + B\langle x_t \rangle_{+}$ を備えた分割冪環とみなす。
『分割冪代数』注 \ref{dpa-remark-divided-power-polynomial-algebra} を参照されたい。
$\Psi$ の存在は分割冪多項式環の普遍性から従う。その像は $D$ の分割冪部分環
であり、$D$ の普遍性により $D$ に等しいので、$\Psi$ は全射である。
$x_t - f_t$ が核に属することは明らかである。次のようにおく：
$$
\mathcal{R} = \{(r_0, r_t) \in I \oplus \bigoplus\nolimits_{t \in T} B
\mid \sum r_t f_t = r_0 \text{於 }B\}
$$
$(r_0, r_t) \in \mathcal{R}$ ならば、$\sum r_t x_t - r_0$ が核に属することは
明らかである。$\Psi$ は分割冪環の準同型であり、
$\sum r_tx_t - r_0 \in J'$ なので、$\delta_n(\sum r_t x_t - r_0)$ も核に属する。
$K \subset B\langle x_t \rangle$ を、$x_t - f_t$ と、
$(r_0, r_t) \in \mathcal{R}$ に対する元 $\delta_n(\sum r_t x_t - r_0)$ で
生成されるイデアルとする。$K = \Ker(\Psi)$ を示すには、$\delta$ が
$B\langle x_t \rangle/K$ に拡張することを示せば十分である。実際、そうであれば
$D$ の普遍性により $\Psi$ の逆写像 $D \to B\langle x_t \rangle/K$ が得られる。
したがって、$K \cap J'$ が $\delta_n$ で保たれることを示せばよい。
『分割冪代数』補題 \ref{dpa-lemma-kernel} を参照されたい。
$K' \subset B\langle x_t \rangle$ を、次の元で生成されるイデアルとする：
\begin{enumerate}
\item $m > 0$ かつ $(r_0, r_t) \in \mathcal{R}$ に対する
$\delta_m(\sum r_t x_t - r_0)$；
\item $m > 0$ かつ $t', t \in I$ に対する $x_{t'}^{[m]}(x_t - f_t)$。
\end{enumerate}
$K' = K \cap J'$ と主張する。この主張と『分割冪代数』補題
\ref{dpa-lemma-kernel} (2)(c) の判定法、および上に挙げた元の
$\delta_n$ の計算（読者に委ねる）から、$K \cap J'$ は $n > 0$ に対する
$\delta_n$ で保たれる。主張を示すため、まず $K' \subset K \cap J'$ に注意する。
逆に $h \in K \cap J'$ とすると、$K'$ を法として
$$
h = \sum r_t (x_t - f_t)
$$
と書ける。ただし $r_t \in B$ である。$h \in K \cap J' \subset J'$ なので、
$r_0 = \sum r_t f_t \in I$ である。よって $(r_0, r_t) \in \mathcal{R}$ であり、
$$
h = \sum r_t x_t - r_0
$$
は $K'$ に属する。これで所望の主張を得る。
\end{proof}

\begin{lemma}
\label{crystalline-lemma-divided-power-envelope-add-variables}
$(A, I, \gamma)$ を分割冪環とする。$B$ を $A$-代数とし、
$IB \subset J \subset B$ をイデアルとする。$x_i$ を変数の族とする。このとき
$$
D_{B[x_i], \gamma}(JB[x_i] + (x_i)) = D_{B, \gamma}(J) \langle x_i \rangle
$$
\end{lemma}

\begin{proof}
一つの証明は、$B[x_i]$ における $x_i$ 間の関係はすべて自明であることを用いて、
補題 \ref{crystalline-lemma-describe-divided-power-envelope} から導くものである。
別の証明として、分割冪多項式代数の普遍性と分割冪包絡の普遍性からも従う。
\end{proof}

\noindent
次の補題の条件 (1) と (2) は、$B \to B'$ が
$V(IB') \subset \Spec(B')$ のすべての素イデアルにおいて平坦であれば成り立ち、
またこの平坦性条件ときわめて密接に関係している。『代数』補題
\ref{algebra-lemma-what-does-it-mean} を参照されたい。
とくにこの補題は、分割冪包絡をとる操作が局所化と可換であることを述べている。

\begin{lemma}
\label{crystalline-lemma-flat-base-change-divided-power-envelope}
$(A, I, \gamma)$ を分割冪環とする。
$B \to B'$ を $A$-代数の準同型とする。次を仮定する：
\begin{enumerate}
\item $B/IB \to B'/IB'$ は平坦である；
\item $\text{Tor}_1^B(B', B/IB) = 0$.
\end{enumerate}
このとき、任意のイデアル $IB \subset J \subset B$ に対して標準写像
$$
D_B(J) \otimes_B B' \longrightarrow D_{B'}(JB')
$$
は同型である。
\end{lemma}

\begin{proof}
$D = D_B(J)$ とおき、その分割冪イデアルを $\bar J \subset D$、
その分割冪構造を $\bar\gamma$ と書く。$D$ の普遍性から $B$-代数写像
$D \to D_{B'}(JB')$ が得られ、したがって補題の写像が得られる。
分割冪 $\bar\gamma$ が $D \otimes_B B'$ に延長されることを示せば十分である。
実際、そのとき $D_{B'}(JB')$ の普遍性から、補題の写像の逆となる写像
$D_{B'}(JB') \to D \otimes_B B'$ が得られる。

\medskip\noindent
$P$ が $B$ 上の多項式代数であるような全射 $P \to B'$ を選ぶ。
とくに $B \to P$ は平坦であるから、『代数』補題
\ref{algebra-lemma-flat-base-change} により
$D \to D \otimes_B P$ も平坦である。このとき『分割冪代数』補題
\ref{dpa-lemma-gamma-extends} により $\bar\gamma$ は
$D \otimes_B P$ に延長される。この延長も $\bar\gamma$ と書く。
$\mathfrak a = \Ker(P \to B')$ とおけば、短完全列
$$
0 \to \mathfrak a \to P \to B' \to 0
$$
を得る。したがって $\text{Tor}_1^B(B', B/IB) = 0$ から
$\mathfrak a \cap IP = I\mathfrak a$ が従う。
ここで次の可換図式がある：
$$
\xymatrix{
B/J \otimes_B \mathfrak a \ar[r]_\beta &
B/J \otimes_B P \ar[r] &
B/J \otimes_B B' \\
D \otimes_B \mathfrak a \ar[r]^\alpha \ar[u] &
D \otimes_B P \ar[r] \ar[u] &
D \otimes_B B' \ar[u] \\
\bar J \otimes_B \mathfrak a \ar[r] \ar[u] &
\bar J \otimes_B P \ar[r] \ar[u] &
\bar J \otimes_B B' \ar[u]
}
$$
この図式は上端と右端に 0 を加えても完全である。
『分割冪代数』補題 \ref{dpa-lemma-kernel} により、
イデアル $\bar J \otimes_B P$ 上の分割冪がイデアル
$\Im(\alpha) \cap \bar J \otimes_B P$ を保つことを示せばよい。
完全列
$$
0 \to \mathfrak a/I\mathfrak a \to P/IP \to B'/IB' \to 0
$$
を考える（ここでは上で示した $\mathfrak a \cap IP = I\mathfrak a$ を用いた）。
$B'/IB'$ は $B/IB$ 上平坦なので、『代数』補題
\ref{algebra-lemma-flat-tor-zero} により、これに
$B/J \otimes_{B/IB} -$ を施しても完全性は保たれる。したがって
$$
\Ker(B/J \otimes_{B/IB} \mathfrak a/I\mathfrak a \to
B/J \otimes_{B/IB} P/IP) =
\Ker(\mathfrak a/J\mathfrak a \to P/JP)
$$
$0$ である。よって $\beta$ は単射である。したがって
$\Im(\alpha) \cap \bar J \otimes_B P$ は
$\bar J \otimes \mathfrak a$ の像である。いま
$f \in \bar J$ および $a \in \mathfrak a$ ならば
$\bar\gamma_n(f \otimes a) = \bar\gamma_n(f) \otimes a^n$
であるから、結論は明らかである。
\end{proof}

\noindent
次の補題は \cite[Proposition 2.1.7]{dJ-crystalline} の特別な場合であり、
後者はさらに \cite[Proposition 2.8.2]{Berthelot} の一般化である。

\begin{lemma}
\label{crystalline-lemma-flat-extension-divided-power-envelope}
$(B, I, \gamma) \to (B', I', \gamma')$ を分割冪環の準同型とする。
$I \subset J \subset B$ および $I' \subset J' \subset B'$ を
イデアルとする。次を仮定する：
\begin{enumerate}
\item $B/I \to B'/I'$ は平坦である；
\item $J' = JB' + I'$.
\end{enumerate}
このとき標準写像
$$
D_{B, \gamma}(J) \otimes_B B' \longrightarrow D_{B', \gamma'}(J')
$$
は同型である。
\end{lemma}

\begin{proof}
$D = D_{B, \gamma}(J)$ とおく。$J/I$ を生成する元 $f_t \in J$ を選ぶ。
補題 \ref{crystalline-lemma-describe-divided-power-envelope} の証明と同様に
$\mathcal{R} = \{(r_0, r_t) \in I \oplus \bigoplus\nolimits_{t \in T} B
\mid \sum r_t f_t = r_0 \text{於 }B\}$ とおく。同補題により
$$
D = B\langle x_t \rangle/ K
$$
であり、ここで $K$ は $x_t - f_t$ と、$(r_0, r_t) \in \mathcal{R}$ に対する
$\delta_n(\sum r_t x_t - r_0)$ によって生成される。したがって
\begin{equation}
\label{crystalline-equation-base-change}
D \otimes_B B' = B'\langle x_t \rangle/K'
\end{equation}
を得る。ここで $K'$ は、上に挙げた $K$ の生成元の
$B'\langle x_t \rangle$ における像によって生成される。
$f_t$ の像を $f'_t \in B'$ と書く。仮定 (1) により、元 $f'_t \in J'$ は
$J'/I'$ を生成し、また $x_t - f'_t \in K'$ である。ここで
$$
\mathcal{R}' =
\{(r'_0, r'_t) \in I' \oplus \bigoplus\nolimits_{t \in T} B'
\mid \sum r'_t f'_t = r'_0 \text{於 }B'\}
$$
とおく。証明を終えるには、$(r'_0, r'_t) \in \mathcal{R}'$ に対して
$\delta'_n(\sum r'_t x_t - r'_0) \in K'$ であることを示せばよい。
実際、そうすれば (\ref{crystalline-equation-base-change}) による
$D \otimes_B B'$ の表示は、補題
\ref{crystalline-lemma-describe-divided-power-envelope} で生成元 $f'_t$ から得られる
$D_{B', \gamma'}(J')$ の表示と一致する。
$(r'_0, r'_t) \in \mathcal{R}'$ とする。このとき
$B'/I'$ において $\sum r'_t f'_t = 0$ である。
仮定 (1) により $B/I \to B'/I'$ は平坦なので、
平坦性の方程式判定法（『代数』補題 \ref{algebra-lemma-flat-eq}）を適用すると、
$m > 0$ と $r_{jt} \in B$ および $c_j \in B'$（$j = 1, \ldots, m$）で
$$
r_{j0} = \sum\nolimits_t r_{jt} f_t  \in I \text{ 各 } j = 1, \ldots, m
$$
かつ
$$
i'_t  = r'_t - \sum\nolimits_j c_j r_{jt} \in I' \text{ 各 }t
$$
を満たすものが存在する。これからさらに
$r'_0 = \sum_t i'_t f_t + \sum_j c_j r_{j0}$ が従うことに注意する。
したがって
\begin{align*}
\delta'_n(\sum\nolimits_t r'_t x_t  - r'_0)
& =
\delta'_n(
\sum\nolimits_t i'_t x_t +
\sum\nolimits_{t, j} c_j r_{jt} x_t -
\sum\nolimits_t i'_t f_t -
\sum\nolimits_j c_j r_{j0}) \\
& =
\delta'_n(
\sum\nolimits_t i'_t(x_t - f_t) +
\sum\nolimits_j c_j (\sum\nolimits_t r_{jt} x_t  - r_{j0}))
\end{align*}
$\delta_n(a + b) = \sum_{m = 0, \ldots, n} \delta_m(a) \delta_{n - m}(b)$
であり、$m > 0$ のとき $\delta_m(\sum i'_t(x_t - f_t))$ は
$x_t - f_t \in K'$ によって生成されるイデアルに属するので、
$\delta_n(\sum c_j (\sum r_{jt} x_t  - r_{j0}))$ が $K'$ に属することを
示せば十分である。そのために次を用いる：
$$
\delta_n(\sum\nolimits_j c_j (\sum\nolimits_t r_{jt} x_t  - r_{j0}))
=
\sum c_1^{n_1} \ldots c_m^{n_m}
\delta_{n_1}(\sum r_{1t} x_t  - r_{10}) \ldots
\delta_{n_m}(\sum r_{mt} x_t  - r_{m0})
$$
ここで和は $n_1 + \ldots + n_m = n$ を満たすものについてとる。
これで所望の結論が得られた。
\end{proof}





\section{いくつかの具体的な分割冪肥厚化}
\label{crystalline-section-explicit-thickenings}

\noindent
この節の構成は、結晶サイト上の加群のクリスタルに接続を定義する際に用いられる。

\begin{lemma}
\label{crystalline-lemma-divided-power-first-order-thickening}
$(A, I, \gamma)$ を分割冪環とし、$M$ を $A$-加群とする。
$M$ が平方零イデアルとなるように $A$-代数として $B = A \oplus M$ とおき、
$J = I \oplus M$ とおく。$x \in I$ および $z \in M$ に対して
$$
\delta_n(x + z) = \gamma_n(x) + \gamma_{n - 1}(x)z
$$
と定める。このとき $\delta$ は分割冪構造であり、$A \to B$ は
$(A, I, \gamma)$ から $(B, J, \delta)$ への分割冪環の準同型である。
\end{lemma}

\begin{proof}
『分割冪代数』定義 \ref{dpa-definition-divided-powers} の条件 (1) -- (5) を
確認すればよい。この場合には直接証明するが、計算を避ける方法については
次の補題の証明を参照されたい。条件 (1) と (3) は明らかである。
条件 (2) は次の計算から従う：
\begin{align*}
\delta_n(x + z)\delta_m(x + z)
& =
(\gamma_n(x) + \gamma_{n - 1}(x)z)(\gamma_m(x) + \gamma_{m - 1}(x)z) \\
& = \gamma_n(x)\gamma_m(x) + \gamma_n(x)\gamma_{m - 1}(x)z +
\gamma_{n - 1}(x)\gamma_m(x)z \\
& =
\frac{(n + m)!}{n!m!} \gamma_{n + m}(x) +
\left(\frac{(n + m - 1)!}{n!(m - 1)!} +
\frac{(n + m - 1)!}{(n - 1)!m!}\right)
\gamma_{n + m - 1}(x) z \\
& =
\frac{(n + m)!}{n!m!}\delta_{n + m}(x + z)
\end{align*}
条件 (5) は次の計算から従う：
\begin{align*}
\delta_n(\delta_m(x + z))
& =
\delta_n(\gamma_m(x) + \gamma_{m - 1}(x)z) \\
& =
\gamma_n(\gamma_m(x)) + \gamma_{n - 1}(\gamma_m(x))\gamma_{m - 1}(x)z \\
& =
\frac{(nm)!}{n! (m!)^n} \gamma_{nm}(x) +
\frac{((n - 1)m)!}{(n - 1)! (m!)^{n - 1}}
\gamma_{(n - 1)m}(x) \gamma_{m - 1}(x) z \\
& = \frac{(nm)!}{n! (m!)^n}(\gamma_{nm}(x) + \gamma_{nm - 1}(x) z)
\end{align*}
ここでは初等整数論を用いた。条件 (4) を証明するには
$$
\delta_n(x + x' + z + z')
=
\gamma_n(x + x') + \gamma_{n - 1}(x + x')(z + z')
$$
が
$$
\sum\nolimits_{i = 0}^n
(\gamma_i(x) + \gamma_{i - 1}(x)z)
(\gamma_{n - i}(x') + \gamma_{n - i - 1}(x')z')
$$
に等しいことを示せばよい。$1$、$z$、$z'$ の係数をまとめ、
$\gamma$ に対する条件 (4) を用いれば容易に従う。
\end{proof}

\begin{lemma}
\label{crystalline-lemma-divided-power-second-order-thickening}
$(A, I, \gamma)$ を分割冪環とし、$M$、$N$ を $A$-加群とする。
$q : M \times M \to N$ を $A$-双線形写像とする。
$A$-代数として $B = A \oplus M \oplus N$ とおき、その積を
$$
(x, z, w)\cdot (x', z', w') = (xx', xz' + x'z, xw' + x'w + q(z, z') + q(z', z))
$$
で定め、$J = I \oplus M \oplus N$ とおく。$(x, z, w) \in J$ に対して
$$
\delta_n(x, z, w) = (\gamma_n(x), \gamma_{n - 1}(x)z,
\gamma_{n - 1}(x)w + \gamma_{n - 2}(x)q(z, z))
$$
と定める。このとき $\delta$ は分割冪構造であり、$A \to B$ は
$(A, I, \gamma)$ から $(B, J, \delta)$ への分割冪環の準同型である。
\end{lemma}

\begin{proof}
『分割冪代数』定義 \ref{dpa-definition-divided-powers} の性質 (4) が
成り立つことを示したいとする。$J$ の元 $(x, z, w)$ と $(x', z', w')$ を選ぶ。
分割冪多項式環の普遍性により、写像
$$
A_0 = \mathbf{Z}\langle s, s'\rangle \longrightarrow A,\quad
s \longmapsto x,
s' \longmapsto x'
$$
を選べる。$M_0 = A_0 \oplus A_0$ および
$N_0 = A_0 \oplus A_0 \oplus M_0 \otimes_{A_0} M_0$ とおく。
$q_0 : M_0 \times M_0 \to N_0$ を明らかな写像とする。
$M_0$ の基底ベクトルを $z$ と $z'$ に送る $A_0$-線形写像として
$M_0 \to M$ を定める。また、$N_0$ の最初の二つの基底ベクトルを
$w$ と $w'$ に送り、最後の直和因子上では
$M_0 \otimes_{A_0} M_0 \to M \otimes_A M \xrightarrow{q} N$
を用いる $A_0$-線形写像として $N_0 \to N$ を定める。
すると、$(A_0, M_0, N_0, q_0)$ の場合に恒等式 (4) を証明すれば十分である。
他の恒等式についても同様である。これにより、$A$ が
$\mathbf{Z}$-捩れをもたない環で、各加群が $A$-捩れをもたない場合に帰着される。
この場合、示すべきことは環 $A$ における等式
$$
n! \delta_n(x, z, w) = (x, z, w)^n
$$
だけである。『分割冪代数』補題 \ref{dpa-lemma-silly} を参照されたい。
このため、まず
$$
(x, z, w)^2 = (x^2, 2xz, 2xw + 2q(z, z))
$$
であり、帰納法により
$$
(x, z, w)^n = (x^n, nx^{n - 1}z, nx^{n - 1}w + n(n - 1)x^{n - 2}q(z, z))
$$
である。一方、
$$
n! \delta_n(x, z, w) = (n!\gamma_n(x), n!\gamma_{n - 1}(x)z,
n!\gamma_{n - 1}(x)w + n!\gamma_{n - 2}(x) q(z, z))
$$
であり、両者は一致する。これで証明が完了する。
\end{proof}




\section{両立性}
\label{crystalline-section-compatibility}

\noindent
この節は必読ではない。ここでは、本章の議論が \cite{Berthelot} の議論と
どのように対応するかを説明する。次の技術的概念を考える。

\begin{definition}
\label{crystalline-definition-compatible}
$(A, I, \gamma)$ と $(B, J, \delta)$ を分割冪環とし、
$A \to B$ を環準同型とする。$J + IB$ 上に分割冪構造 $\bar\gamma$ が存在して、
$$
(A, I, \gamma) \to (B, J + IB, \bar \gamma)\quad\text{および}\quad
(B, J, \delta) \to (B, J + IB, \bar \gamma)
$$
がいずれも分割冪環の準同型となるとき、
{\it $\delta$ は $\gamma$ と両立する}という。
\end{definition}

\noindent
$p$ を素数とし、$(A, I, \gamma)$ を分割冪環とする。
$A \to C$ を、$C$ において $p$ が冪零となる環準同型とする。
$\gamma$ が $IC$ に延長されると仮定する（『分割冪代数』補題
\ref{dpa-lemma-gamma-extends} を参照）。この状況で、\cite{Berthelot} により
定義される $\Spec(A)$ 上の $\Spec(C)$ の（大アフィン）結晶サイトは、
次の系のなす圏の反対圏である：
$$
(B, J, \delta, A \to B, C \to B/J)
$$
ここで
\begin{enumerate}
\item $(B, J, \delta)$ は $B$ において $p$ が冪零となる分割冪環である；
\item $\delta$ は $\gamma$ と両立する；
\item 図式
$$
\xymatrix{
B \ar[r] & B/J \\
A \ar[u] \ar[r] & C \ar[u]
}
$$
は可換である。
\end{enumerate}
\cite{Berthelot} では、「$\gamma$ が $C$ に延長され、$\delta$ が
$\gamma$ と両立する」という条件を用いて、$\Spec(C)$ の結晶コホモロジーが
$\Spec(C/IC)$ の結晶コホモロジーと同じになることを保証している。
本章では $IC = 0$ を満たす $C$ のみを扱うことで、この問題を避ける
\footnote{もちろん、その代償はある。}。この場合、上の系
$(B, J, \delta, A \to B, C \to B/J)$ に対して、表示された図式の可換性から
$IB \subset J$ が従い、両立性は
$(A, I, \gamma) \to (B, J, \delta)$ が分割冪環の準同型であるという条件と
同値になる。




\section{アフィン結晶サイト}
\label{crystalline-section-affine-site}

\noindent
この節では結晶サイトの代数的な変種を論じる。以下では、次の状況を
基本的な設定として用いる。

\begin{situation}
\label{crystalline-situation-affine}
ここで $p$ は素数、$(A, I, \gamma)$ は $A$ が
$\mathbf{Z}_{(p)}$-代数であるような分割冪環であり、$A \to C$ は
$IC = 0$ かつ $C$ において $p$ が冪零となる環準同型である。
\end{situation}

\noindent
通常、素数 $p$ は分割冪イデアル $I$ に含まれる。

\begin{definition}
\label{crystalline-definition-affine-thickening}
状況 \ref{crystalline-situation-affine} において、次のように定める。
\begin{enumerate}
\item $(A, I, \gamma)$ 上の $C$ の {\it 分割冪肥厚化}とは、
$B$ において $p$ が冪零となるような分割冪代数の準同型
$(A, I, \gamma) \to (B, J, \delta)$ と、次の図式を可換にする環準同型
$C \to B/J$ の組である：
$$
\xymatrix{
B \ar[r] & B/J \\
& C \ar[u] \\
A \ar[uu] \ar[r] & A/I \ar[u]
}
$$
\item {\it 分割冪肥厚化の準同型}
$$
(B, J, \delta, C \to B/J) \longrightarrow (B', J', \delta', C \to B'/J')
$$
とは、合成 $C \to B/J \to B'/J'$ が与えられた写像 $C \to B'/J'$ と
一致するような分割冪 $A$-代数の準同型 $\varphi : B \to B'$ である。
\item $(A, I, \gamma)$ 上の $C$ の分割冪肥厚化の圏を
$\text{CRIS}(C/A, I, \gamma)$、または単に $\text{CRIS}(C/A)$ と書く。
\item $C \to B/J$ が同型であるような $(B, J, \delta, C \to B/J)$ からなる
充満部分圏を $\text{Cris}(C/A, I, \gamma)$、または単に
$\text{Cris}(C/A)$ と書く。このような対象は、しばしば
$J = \Ker(B \to C)$ を暗黙にして $(B \to C, \delta)$ と書く。
\end{enumerate}
\end{definition}

\noindent
上の分割冪肥厚化 $(B, J, \delta)$ に対してイデアル $J$ は局所冪零であることに
注意する。『分割冪代数』補題 \ref{dpa-lemma-nil} を参照されたい。
標準関手
\begin{equation}
\label{crystalline-equation-forget-affine}
\text{CRIS}(C/A) \longrightarrow C\text{-algebras},\quad
(B, J, \delta) \longmapsto B/J
\end{equation}
がある。この圏は一般には等化子もファイバー積ももたない。
また一般には始対象（$=$ 空余極限）ももたない。

\begin{lemma}
\label{crystalline-lemma-affine-thickenings-colimits}
状況 \ref{crystalline-situation-affine} において、次が成り立つ。
\begin{enumerate}
\item $\text{CRIS}(C/A)$ は有限積をもつ（ただし無限積は一般にもたない）。
\item $\text{CRIS}(C/A)$ はすべての空でない有限余極限をもち、
(\ref{crystalline-equation-forget-affine}) はそれらと可換である。
\item $\text{Cris}(C/A)$ はすべての空でない有限余極限をもち、
$\text{Cris}(C/A) \to \text{CRIS}(C/A)$ はそれらと可換である。
\end{enumerate}
\end{lemma}

\begin{proof}
空積、すなわち $(A, I, \gamma)$ 上の $C$ の分割冪肥厚化の圏における終対象は
零環である。ここで零環を $A$-代数とみなし、零イデアルと、その零イデアル上の
一意な分割冪を入れ、さらに $C$ から零環への一意な準同型を入れる。
$(B_t, J_t, \delta_t)_{t \in T}$ が $\text{CRIS}(C/A)$ の対象の族なら、
『分割冪代数』補題 \ref{dpa-lemma-limits} と同様に積
$(\prod_t B_t, \prod_t J_t, \prod_t \delta_t)$ を作れる。
写像 $C \to \prod B_t/\prod J_t = \prod B_t/J_t$ は明らかである。
ただし、$\prod_t B_t$ において $p$ が冪零であることが保証されるのは
$T$ が有限の場合だけである。

\medskip\noindent
$\text{CRIS}(C/A)$ の二つの対象 $(B, J, \gamma)$ と
$(B', J', \gamma')$ が与えられたとする。分割冪環の圏において余カルテシアン図式
$$
\xymatrix{
(B, J, \delta) \ar[r] & (B'', J'', \delta'') \\
(A, I, \gamma) \ar[r] \ar[u] & (B', J', \delta') \ar[u]
}
$$
を作ることができる。このとき
$$
B''/J'' = B/J \otimes_{A/I} B'/J' \longleftarrow C \otimes_{A/I} C
$$
を得る。『分割冪代数』注意 \ref{dpa-remark-pushout} を参照されたい。
次の図式が押し出しとなるようなイデアル $J'' \subset K \subset B''$ をとる：
$$
\xymatrix{
B''/J'' \ar[r] & B''/K \\
C \otimes_{A/I} C \ar[r] \ar[u] & C \ar[u]
}
$$
すなわち $B''/K \cong B/J \otimes_C B'/J'$ である。
$D_{B''}(K) = (D, \bar K, \bar \delta)$ を、$(B'', J'', \delta'')$ に関する
$B''$ 内の $K$ の分割冪包絡とする。このとき
$(D, \bar K, \bar \delta)$ が $\text{CRIS}(C/A)$ における
$(B, J, \delta)$ と $(B', J', \delta')$ の余積であることは容易に確認できる。

\medskip\noindent
次に余等化子を扱う。$\alpha, \beta : (B, J, \delta) \to
(B', J', \delta')$ を $\text{CRIS}(C/A)$ の射とする。
$B'' = B'/ (\alpha(b) - \beta(b))$ とし、$J'' \subset B''$ を $J'$ の像とする。
$D_{B''}(J'') = (D, \bar J, \bar\delta)$ を、$(B', J', \delta')$ に関する
$B''$ 内の $J''$ の分割冪包絡とする。このとき
$(D, \bar J, \bar \delta)$ が $\text{CRIS}(C/A)$ における
$(B, J, \delta)$ と $(B', J', \delta')$ の余等化子であることは容易に確認できる。

\medskip\noindent
『圏論』補題 \ref{categories-lemma-almost-finite-colimits-exist} により、
$\text{CRIS}(C/A)$ はすべての空でない有限余極限をもつ。
上の構成から、(\ref{crystalline-equation-forget-affine}) がそれらと可換であることが分かる。
$\text{Cris}(C/A)$ は (\ref{crystalline-equation-forget-affine}) の $C$ 上のファイバー圏なので、
これから形式的に (3) が従う。
\end{proof}

\begin{remark}
\label{crystalline-remark-completed-affine-site}
状況 \ref{crystalline-situation-affine} において、次の条件を満たす組
$(B \to C, \delta)$ を対象とする圏を $\text{Cris}^\wedge(C/A)$ と書く：
\begin{enumerate}
\item $B$ は $p$ 進完備な $A$-代数である；
\item $B \to C$ は $A$-代数の全射である；
\item $\delta$ は $\Ker(B \to C)$ 上の分割冪構造である；
\item $A \to B$ は分割冪環の準同型である。
\end{enumerate}
射は定義 \ref{crystalline-definition-affine-thickening} と同様に定める。
このとき $\text{Cris}(C/A) \subset \text{Cris}^\wedge(C/A)$ は、
$B$ において $p$ が冪零であるような $B$ からなる充満部分圏である。
逆に、$\text{Cris}^\wedge(C/A)$ の任意の対象 $(B \to C, \delta)$ は極限
$$
(B \to C, \delta) = \lim_e (B/p^eB \to C, \delta)
$$
に等しい。ここで $e \gg 0$ に対して対象 $(B/p^eB \to C, \delta)$ は
$\text{Cris}(C/A)$ に属する。『分割冪代数』補題
\ref{dpa-lemma-extend-to-completion} を参照されたい。
とくに、$\text{Cris}^\wedge(C/A)$ は $\text{Cris}(C/A)$ の pro 対象の圏の
充満部分圏である。『圏論』注意 \ref{categories-remark-pro-category} を参照されたい。
\end{remark}

\begin{lemma}
\label{crystalline-lemma-list-properties}
状況 \ref{crystalline-situation-affine} において、$P \to C$ を核 $J$ をもつ
$A$-代数の全射とし、$D_{P, \gamma}(J) = (D, \bar J, \bar\gamma)$ と書く。
$(D^\wedge, J^\wedge, \bar\gamma^\wedge)$ を $D$ の $p$ 進完備化とする。
『分割冪代数』補題 \ref{dpa-lemma-extend-to-completion} を参照されたい。
各 $e \geq 1$ に対して $P_e = P/p^eP$ とおき、$J_e \subset P_e$ を
$J$ の像とし、$D_{P_e, \gamma}(J_e) = (D_e, \bar J_e, \bar\gamma)$ と書く。
このとき、十分大きいすべての $e$ に対して次が成り立つ。
\begin{enumerate}
\item $p^eD \subset \bar J$ と $p^eD^\wedge \subset \bar J^\wedge$ は
分割冪によって保たれる。
\item 分割冪環として $D^\wedge/p^eD^\wedge = D/p^eD = D_e$ である。
\item $(D_e, \bar J_e, \bar\gamma)$ は $\text{Cris}(C/A)$ の対象である。
\item $(D^\wedge, \bar J^\wedge, \bar\gamma^\wedge)$ は
$\lim_e (D_e, \bar J_e, \bar\gamma)$ に等しい。
\item $(D^\wedge, \bar J^\wedge, \bar\gamma^\wedge)$ は
$\text{Cris}^\wedge(C/A)$ の対象である。
\end{enumerate}
\end{lemma}

\begin{proof}
(1) は『分割冪代数』補題 \ref{dpa-lemma-extend-to-completion} から従う。
$p$ 進完備化の一般的性質により $D/p^eD = D^\wedge/p^eD^\wedge$ である。
$D/p^eD$ は分割冪環であり、$P \to D/p^eD$ は $P_e$ を経由するので、
$D_e$ の普遍性から写像 $D_e \to D/p^eD$ が得られる。逆に、$D$ の普遍性から
$D \to D_e$ が得られ、これは $D/p^eD$ を経由する。これらの写像が互いに逆で
あることの確認は省略する。これで (2) が示された。$e$ が十分大きければ
$p^eC = 0$ なので (3) が成り立つ。(4) は『分割冪代数』補題
\ref{dpa-lemma-extend-to-completion} から従い、(5) は定義から明らかである。
\end{proof}

\begin{lemma}
\label{crystalline-lemma-set-generators}
状況 \ref{crystalline-situation-affine} において、$P$ を $A$ 上の多項式代数とし、
$P \to C$ を核 $J$ をもつ $A$-代数の全射とする。
$(D_e, \bar J_e, \bar\gamma)$ を補題 \ref{crystalline-lemma-list-properties} のものとする。
$\text{CRIS}(C/A)$ の任意の対象 $(B, J_B, \delta)$ に対し、ある $e$ と
$\text{CRIS}(C/A)$ の射 $D_e \to B$ が存在する。
\end{lemma}

\begin{proof}
写像 $C \to B/J_B$ を持ち上げる $A$-代数準同型 $P \to B$ を選べる。
$\text{CRIS}(C/A)$ の定義により、ある $e$ に対して $p^eB = 0$ である。
したがって $P \to B$ は $P \to P_e \to B$ と分解する。
分割冪包絡の普遍性により、$P_e \to B$ は $D_e$ を経由する。
\end{proof}

\begin{lemma}
\label{crystalline-lemma-generator-completion}
状況 \ref{crystalline-situation-affine} において、$P$ を $A$ 上の多項式代数とし、
$P \to C$ を核 $J$ をもつ $A$-代数の全射とする。
$(D, \bar J, \bar\gamma)$ を $D_{P, \gamma}(J)$ の $p$ 進完備化とする。
$\text{Cris}^\wedge(C/A)$ の任意の対象 $(B \to C, \delta)$ に対して、
$\text{Cris}^\wedge(C/A)$ の射 $D \to B$ が存在する。
\end{lemma}

\begin{proof}
$C$ への写像と両立する $A$-代数準同型 $P \to B$ を選べる。
$\text{Cris}^\wedge(C/A)$ の定義により、$P \to B$ は
$P \to D_{P, \gamma}(J) \to B$ と分解する。$B$ は $p$ 進完備なので、
この写像を $D$ を経由して分解できる。
\end{proof}



\section{微分加群}
\label{crystalline-section-differentials}

\noindent
この節では分割冪環に対する微分加群の理論を展開する。

\begin{definition}
\label{crystalline-definition-derivation}
$A$ を環、$(B, J, \delta)$ を分割冪環、$A \to B$ を環準同型、
$M$ を $B$-加群とする。$M$ に値をとる {\it 分割冪 $A$-導分}とは、
加法的で、$A$ の元を零に写し、Leibniz 則
$\theta(bb') = b\theta(b') + b'\theta(b)$ を満たし、さらにすべての
$n \geq 1$ と $x \in J$ に対して
$$
\theta(\delta_n(x)) = \delta_{n - 1}(x)\theta(x)
$$
を満たす写像 $\theta : B \to M$ である。
\end{definition}

\noindent
この定義の状況では、通常の導分の場合と同様に、
{\it 普遍分割冪 $A$-導分}
$$
\text{d}_{B/A, \delta} : B \to \Omega_{B/A, \delta}
$$
が存在し、任意の分割冪 $A$-導分 $\theta : B \to M$ は、一意な $B$-線形写像
$\xi : \Omega_{B/A, \delta} \to M$ により
$\theta = \xi \circ d_{B/A, \delta}$ と書ける。
$(A, I, \gamma) \to (B, J, \delta)$ が分割冪環の準同型ならば、
$A$ 上の分割冪を忘れて、$A$ 上の $B$ の分割冪導分を考えることができる。
普遍（分割冪）微分加群の基本的性質をいくつか挙げる。

\begin{lemma}
\label{crystalline-lemma-omega}
$A$ を環、$(B, J, \delta)$ を分割冪環、$A \to B$ を環準同型とする。
\begin{enumerate}
\item $\delta'$ を $\delta$ の $B[x]$ への延長とし、$B[x]$ に分割冪イデアル
$(JB[x], \delta')$ を入れる。このとき
$$
\Omega_{B[x]/A, \delta'} =
\Omega_{B/A, \delta} \otimes_B B[x] \oplus B[x]\text{d}x.
$$
\item $B\langle x \rangle$ に分割冪イデアル
$(JB\langle x \rangle + B\langle x \rangle_{+}, \delta')$ を入れる。このとき
$$
\Omega_{B\langle x\rangle/A, \delta'} =
\Omega_{B/A, \delta} \otimes_B B\langle x \rangle \oplus
B\langle x\rangle \text{d}x.
$$
\item $K \subset J$ をすべての $n > 0$ に対して $\delta_n$ で保たれる
イデアルとする。$B' = B/K$ とおき、$J/K$ 上に誘導される分割冪を
$\delta'$ と書く。このとき $\Omega_{B'/A, \delta'}$ は
$\Omega_{B/A, \delta} \otimes_B B'$ を、$k \in K$ に対する
$\text{d}k$ で生成される $B'$-部分加群で割った商である。
\end{enumerate}
\end{lemma}

\begin{proof}
これらは、$\Omega_{B/A, \delta}$ を元 $\text{d}b$ 上の自由 $B$-加群を
次の関係で割ったものとして構成することから直接従う：
\begin{enumerate}
\item $\text{d}(b + b') = \text{d}b + \text{d}b'$, $b, b' \in B$,
\item $\text{d}a = 0$, $a \in A$,
\item $\text{d}(bb') = b \text{d}b' + b' \text{d}b$, $b, b' \in B$,
\item $\text{d}\delta_n(f) = \delta_{n - 1}(f)\text{d}f$, $f \in J$, $n > 1$.
\end{enumerate}
最後の関係は、分割冪多項式代数と通常の多項式代数について「同じ」答えが
得られる理由を説明している。前者では $x$ が分割冪イデアルの元であるため、
$\text{d}x^{[n]} = x^{[n - 1]}\text{d}x$ となる。
\end{proof}

\noindent
$(A, I, \gamma)$ を分割冪環とする。この設定では、$I$ の通常の冪に対応する
適切な概念は次の分割冪で与えられる：
$$
I^{[n]} = \text{次で生成されるイデアル：}
\gamma_{e_1}(x_1) \ldots \gamma_{e_t}(x_t)
\text{ ただし }\sum e_j \geq n\text{ かつ }x_j \in I.
$$
もちろん $I^n \subset I^{[n]}$ である。また $I^{[1]} = I$ である。
$I^{[0]} = A$ とおくこともある。

\begin{lemma}
\label{crystalline-lemma-diagonal-and-differentials}
$(A, I, \gamma) \to (B, J, \delta)$ を分割冪環の準同型とする。
$(B(1), J(1), \delta(1))$ を $(A, I, \gamma)$ 上での
$(B, J, \delta)$ とそれ自身との余積、すなわち次の図式が
余カルテシアンとなるものとする：
$$
\xymatrix{
(B, J, \delta) \ar[r] & (B(1), J(1), \delta(1)) \\
(A, I, \gamma) \ar[r] \ar[u] & (B, J, \delta) \ar[u]
}
$$
$K = \Ker(B(1) \to B)$ と書く。このとき $K \cap J(1) \subset J(1)$ は
分割冪構造で保たれ、標準的に
$$
\Omega_{B/A, \delta} = K/ \left(K^2 + (K \cap J(1))^{[2]}\right)
$$
である。
\end{lemma}

\begin{proof}
$B(1) \to B$ は分割冪環の準同型なので、$K \cap J(1) \subset J(1)$ が
分割冪構造で保たれることが従う。

\medskip\noindent
$K/K^2$ が標準的な $B$-加群構造をもつことを思い出す。
二つの余射を $s_0, s_1 : B \to B(1)$ と書き、
$b \mapsto s_1(b) - s_0(b)$ で与えられる写像
$\text{d} : B \to K/(K^2 +(K \cap J(1))^{[2]})$ を考える。
$\text{d}$ が加法的で、$A$ を零に写し、Leibniz 則を満たすことは明らかである。
$\text{d}$ が分割冪 $A$-導分であることを示す。$x \in J$ とし、
$y = s_1(x)$ および $z = s_0(x)$ とおく。$J(1)$ 上の分割冪構造
$\delta(1)$ も $\delta$ と書く。$n \geq 1$ に対して、
$K^2 +(K \cap J(1))^{[2]}$ を法として
$\delta_n(y) - \delta_n(z) = \delta_{n - 1}(y)(y - z)$ を示せばよい。
$n = 1$ では等式は成り立つ。$n > 1$ とする。
$i > 1$ に対して $\delta_i(y - z)$ は $(K \cap J(1))^{[2]}$ に属する。
$K^2 + (K \cap J(1))^{[2]}$ を法として計算すると
$$
\delta_n(z) = \delta_n(z - y + y) =
\sum\nolimits_{i = 0}^n \delta_i(z - y)\delta_{n - i}(y) =
\delta_{n - 1}(y) \delta_1(z - y) + \delta_n(y)
$$
となる。これで所望の等式が示された。

\medskip\noindent
$M$ を $B$-加群とし、$\theta : B \to M$ を分割冪 $A$-導分とする。
$M$ が平方零イデアルとなるように $D = B \oplus M$ とおく。
$n > 1$ に対して $\delta_n(x + m) = \delta_n(x) + \delta_{n - 1}(x)m$
と定めることにより、$J \oplus M \subset D$ 上に分割冪構造を入れる。
補題 \ref{crystalline-lemma-divided-power-first-order-thickening} を参照されたい。
分割冪代数の準同型 $B \to D$ が二つある。一つは包含であり、もう一つは
$b \mapsto b + \theta(b)$ で与えられる。したがって $(A, I, \gamma)$ 上の
分割冪代数の標準準同型 $B(1) \to D$ を得る。これは写像 $K \to M$ を誘導し、
$M$ が平方零イデアルなので $K^2$ を零に写し、$M^{[2]} = 0$ なので
$(K \cap J(1))^{[2]}$ も零に写す。構成により合成
$B \to K/K^2 + (K \cap J(1))^{[2]} \to M$ は $\theta$ に等しい。
したがって $\text{d}$ は普遍分割冪 $A$-導分であり、結論が従う。
\end{proof}

\begin{remark}
\label{crystalline-remark-filtration-differentials}
$A \to B$ を環準同型とし、$(J, \delta)$ を $B$ 上の分割冪構造とする。
普遍加群 $\Omega_{B/A, \delta}$ には少し余分な構造がある。すなわち、
$\text{d}_{B/A, \delta}(J)$ で生成される $\Omega_{B/A, \delta}$ の
$B$-部分加群 $N$ である。補題 \ref{crystalline-lemma-diagonal-and-differentials} の
同型のもとでは、これは $\Omega_{B/A, \delta}$ における
$K \cap J(1)$ の像に対応する。補題の証明と同様に、イデアル
$\bar J = J \oplus N$ と分割冪 $\bar \delta$ を備えた $A$-代数
$D = B \oplus \Omega^1_{B/A, \delta}$ を考える。このとき
$(D, \bar J, \bar \delta)$ は分割冪環であり、$b \mapsto b$ および
$b \mapsto b + \text{d}_{B/A, \delta}(b)$ で与えられる二つの写像
$B \to D$ は $A$ 上の分割冪環の準同型である。さらに $N$ は、
この性質が成り立つような $\Omega_{B/A, \delta}$ の最小の部分加群である。
\end{remark}

\begin{lemma}
\label{crystalline-lemma-diagonal-and-differentials-affine-site}
状況 \ref{crystalline-situation-affine} において、$(B, J, \delta)$ を
$\text{CRIS}(C/A)$ の対象とする。$(B(1), J(1), \delta(1))$ を
$\text{CRIS}(C/A)$ における $(B, J, \delta)$ とそれ自身との余積とし、
$K = \Ker(B(1) \to B)$ と書く。このとき $K \cap J(1) \subset J(1)$ は
分割冪構造で保たれ、標準的に
$$
\Omega_{B/A, \delta} = K/ \left(K^2 + (K \cap J(1))^{[2]}\right)
$$
である。
\end{lemma}

\begin{proof}
補題 \ref{crystalline-lemma-diagonal-and-differentials} の証明と逐語的に同じである。
確認すべき唯一の点は、分割冪環 $D = B \oplus M$ が
$\text{CRIS}(C/A)$ の対象であり、二つの写像 $B \to D$ が
$\text{CRIS}(C/A)$ の射であることである。$D/(J \oplus M) = B/J$ なので、
$C \to B/J$ を用いて $D$ を $\text{CRIS}(C/A)$ の対象とみなせる。
射についての主張は構成から明らかである。
\end{proof}

\begin{lemma}
\label{crystalline-lemma-module-differentials-divided-power-envelope}
$(A, I, \gamma)$ を分割冪環、$A \to B$ を環準同型、
$IB \subset J \subset B$ をイデアルとする。
$D_{B, \gamma}(J) = (D, \bar J, \bar \gamma)$ を分割冪包絡とする。このとき
$$
\Omega_{D/A, \bar\gamma} = \Omega_{B/A} \otimes_B D
$$
である。
\end{lemma}

\begin{proof}[First proof]
$M$ を $D$-加群とする。$A$-導分 $\vartheta : B \to M$ を与えることと、
分割冪 $A$-導分 $\theta : D \to M$ を与えることは同じであると主張する。
この主張と Yoneda の補題から結論が従う。

\medskip\noindent
$D$ の平方零肥厚化 $D \oplus M$ を考える。$M$ 上の高次分割冪演算を
零と定めれば、$\bar J \oplus M$ 上に分割冪構造 $\delta$ が入る。
言い換えると、任意の $x \in \bar J$ と $m \in M$ に対して
$\delta_n(x + m) = \bar\gamma_n(x) + \bar\gamma_{n - 1}(x)m$ と定める。
補題 \ref{crystalline-lemma-divided-power-first-order-thickening} を参照されたい。
第一成分が写像 $B \to D$、第二成分が $\vartheta$ である $A$-代数写像
$B \to D \oplus M$ を考える。普遍性により、対応する分割冪代数の準同型
$D \to D \oplus M$ が得られ、その第二成分は $\vartheta$ に対応する
分割冪 $A$-導分 $\theta$ である。
\end{proof}

\begin{proof}[Second proof]
まず $B$ が $A$ 上平坦な場合を証明する。この場合、『分割冪代数』補題
\ref{dpa-lemma-gamma-extends} により、$\gamma$ は $IB$ 上の分割冪構造
$\gamma'$ に延長される。したがって $D = D_{B, \gamma'}(J)$ は、
$D' =  B\langle x_t \rangle$ および
$J' = IB\langle x_t \rangle + B\langle x_t \rangle_{+}$ とおいた分割冪環
$(D', J', \delta)$ を、元 $x_t - f_t$ と
$\delta_n(\sum r_t x_t - r_0)$ で割った商である。記法と説明については
補題 \ref{crystalline-lemma-describe-divided-power-envelope} を参照されたい。
普遍導分を $\text{d} : D' \to \Omega_{D'/A, \delta}$ と書く。このとき
$$
\Omega_{D'/A, \delta} =
\Omega_{B/A} \otimes_B D' \oplus \bigoplus D' \text{d}x_t,
$$
である。補題 \ref{crystalline-lemma-omega} を参照されたい。したがって
$\Omega_{D/A, \bar\gamma}$ は $\Omega_{D'/A, \delta} \otimes_{D'} D$ を、
上に挙げた $D' \to D$ の核の生成元に $\text{d}$ を施して得られる元で
生成される部分加群で割った商である。再び補題 \ref{crystalline-lemma-omega} を用いた。
$\text{d}(x_t - f_t) = - \text{d}f_t + \text{d}x_t$ なので、商では
$\text{d}x_t = \text{d}f_t$ である。とくに
$\Omega_{B/A} \otimes_B D \to \Omega_{D/A, \gamma}$ は全射であり、
その核は元 $\delta_n(\sum r_t x_t - r_0)$ に $\text{d}$ を施した像で与えられる。
一方、$r_t \in B$ および $r_0 \in IB$ に対して $B$ における関係
$\sum r_tf_t - r_0 = 0$ が与えられると
\begin{align*}
\text{d}\delta_n(\sum r_t x_t - r_0)
& =
\delta_{n - 1}(\sum r_t x_t - r_0)\text{d}(\sum r_t x_t - r_0)
\\
& =
\delta_{n - 1}(\sum r_t x_t - r_0)
\left(
\sum r_t\text{d}(x_t - f_t) + \sum (x_t - f_t)\text{d}r_t
\right)
\end{align*}
である。ここで $B$ において $\sum r_tf_t - r_0 = 0$ を用いた。
したがって、これはすでに $\Omega_{B/A} \otimes_A D$ において零であり、
$B$ が $A$ 上平坦な場合が示された。

\medskip\noindent
一般の場合、$B$ を多項式環の商 $P \to B$ として表し、$J' \subset P$ を
$J$ の逆像とする。補題 \ref{crystalline-lemma-divided-power-envelop-quotient} の記法により
$D = D'/K'$ である。証明の第一段落で扱った場合から
$\Omega_{D'/A, \bar\gamma'} = \Omega_{P/A} \otimes_P D'$ を得る。
このとき $\Omega_{D/A, \bar \gamma}$ は
$\Omega_{P/A} \otimes_P D$ を、$P \to B$ の核の元 $k$ に対する
$\text{d}\bar\gamma_n'(k)$ で生成される部分加群で割った商である。
補題 \ref{crystalline-lemma-omega} と、補題
\ref{crystalline-lemma-divided-power-envelop-quotient} における $K'$ の記述を用いた。
$\text{d}\bar\gamma_n'(k) = \bar\gamma'_{n - 1}(k)\text{d}k$ なので、
再び $k \in \Ker(P \to B)$ に対する $\text{d}k$ で生成される部分加群で
割れば十分である。また $\Omega_{B/A}$ は
$\Omega_{P/A} \otimes_A B$ をこれらの元で割った商なので、結論が従う。
『代数』補題 \ref{algebra-lemma-differential-seq} を参照されたい。
\end{proof}

\begin{remark}
\label{crystalline-remark-divided-powers-de-rham-complex}
$A \to B$ を環準同型とし、$(J, \delta)$ を $B$ 上の分割冪構造とする。
$\Omega_{B/A, \delta}^i = \wedge^i_B \Omega_{B/A, \delta}$ とおく。
ここで $\Omega_{B/A, \delta}$ は普遍分割冪 $A$-導分
$\text{d} = \text{d}_{B/A} : B \to \Omega_{B/A, \delta}$ の終域である。
$\Omega_{B/A, \delta}$ は $\Omega_{B/A}$ を、$x \in J$ に対する元
$\text{d}\delta_n(x) - \delta_{n - 1}(x)\text{d}x$ で生成される
$B$-部分加群で割った商である。『代数』補題
\ref{algebra-lemma-de-rham-complex} を適用できると主張する。
そのためには、$\Omega_B$ の元
$\text{d}\delta_n(x) - \delta_{n - 1}(x)\text{d}x$ が
$\Omega^2_{B/A, \delta}$ において零に写ることを確認すれば十分である。実際、
$$
\text{d}(\delta_{n - 1}(x)) \wedge \text{d}x
= \delta_{n - 2}(x) \text{d}x \wedge \text{d}x = 0
$$
が $\Omega^2_{B/A, \delta}$ において成り立つ。したがって
{\it 分割冪 de Rham 複体}
$$
\Omega^0_{B/A, \delta} \to \Omega^1_{B/A, \delta} \to
\Omega^2_{B/A, \delta} \to \ldots
$$
を得る。これは後で重要な役割を果たす。
\end{remark}

\begin{remark}
\label{crystalline-remark-connection}
$A \to B$ を環準同型とし、$\Omega_{B/A} \to \Omega$ を
『代数』補題 \ref{algebra-lemma-de-rham-complex} の仮定を満たす商とする。
$M$ を $B$-加群とする。{\it 接続}とは加法写像
$$
\nabla : M \longrightarrow M \otimes_B \Omega
$$
であって、$b \in B$ と $m \in M$ に対して
$\nabla(bm) = b \nabla(m) + m \otimes \text{d}b$ を満たすものをいう。
この状況では、写像
$$
\nabla : M \otimes_B \Omega^i \longrightarrow M \otimes_B \Omega^{i + 1}
$$
を $\nabla(m \otimes \omega) = \nabla(m) \wedge \omega +
m \otimes \text{d}\omega$ によって定義できる。実際、$b \in B$ ならば
\begin{align*}
\nabla(bm \otimes \omega) - \nabla(m \otimes b\omega)
& =
\nabla(bm) \wedge \omega + bm \otimes \text{d}\omega
- \nabla(m) \wedge b\omega - m \otimes \text{d}(b\omega) \\
& =
b\nabla(m) \wedge \omega + m \otimes \text{d}b \wedge \omega
+ bm \otimes \text{d}\omega \\
& \ \ \ \ \ \ - b\nabla(m) \wedge \omega - bm \otimes \text{d}(\omega)
- m \otimes \text{d}b \wedge \omega = 0
\end{align*}
である。通常どおり、合成
$$
M \xrightarrow{\nabla} M \otimes_B \Omega^1
\xrightarrow{\nabla} M \otimes_B \Omega^2
$$
が零であるとき、かつそのときに限り、接続は {\it 可積分}であるという。
この場合、複体
$$
M \xrightarrow{\nabla} M \otimes_B \Omega^1
\xrightarrow{\nabla} M \otimes_B \Omega^2
\xrightarrow{\nabla} M \otimes_B \Omega^3
\xrightarrow{\nabla} M \otimes_B \Omega^4 \to \ldots
$$
を得る。これを接続の de Rham 複体という。
\end{remark}

\begin{remark}
\label{crystalline-remark-base-change-connection}
環の可換図式
$$
\xymatrix{
B \ar[r]_\varphi & B' \\
A \ar[u] \ar[r] & A' \ar[u]
}
$$
を考える。$\Omega_{B/A} \to \Omega$ と $\Omega_{B'/A'} \to \Omega'$ を、
『代数』補題 \ref{algebra-lemma-de-rham-complex} の仮定を満たす商とする。
次の図式を可換にする写像 $\varphi : \Omega \to \Omega'$ があると仮定する：
$$
\xymatrix{
\Omega_{B/A} \ar[r] \ar[d] &
\Omega_{B'/A'} \ar[d] \\
\Omega \ar[r]^{\varphi} &
\Omega'
}
$$
ここで上の水平矢印は、$\varphi : B \to B'$ により誘導される標準写像
$\Omega_{B/A} \to \Omega_{B'/A'}$ である。この状況で、$M$ が $B$-加群、
$\nabla : M \to M \otimes_B \Omega$ が接続であるような任意の組
$(M, \nabla)$ に対し、{\it 基底変換} $(M \otimes_B B', \nabla')$ を得る。
ここで
$$
\nabla' :
M \otimes_B B'
\longrightarrow
(M \otimes_B B') \otimes_{B'} \Omega' = M \otimes_B \Omega'
$$
は、$\nabla(m) = \sum m_i \otimes \text{d}b_i$ のとき
$$
\nabla'(m \otimes b') =
\sum m_i \otimes b'\text{d}\varphi(b_i) + m \otimes \text{d}b' 
$$
により定義される。$\nabla$ が可積分ならば $\nabla'$ も可積分であり、
この場合、de Rham 複体の標準写像（注意 \ref{crystalline-remark-connection}）
\begin{equation}
\label{crystalline-equation-base-change-map-complexes}
M \otimes_B \Omega^\bullet
\longrightarrow
(M \otimes_B B') \otimes_{B'} (\Omega')^\bullet =
M \otimes_B (\Omega')^\bullet
\end{equation}
があり、これは $m \otimes \eta$ を $m \otimes \varphi(\eta)$ に写す。
\end{remark}

\begin{lemma}
\label{crystalline-lemma-differentials-completion}
$A \to B$ を環準同型、$(J, \delta)$ を $B$ 上の分割冪構造、
$p$ を素数とする。$A$ が $\mathbf{Z}_{(p)}$-代数であり、$B/J$ において
$p$ が冪零であると仮定する。このとき
$$
\lim_e \Omega_{B_e/A, \bar\delta} =
\lim_e \Omega_{B/A, \delta}/p^e\Omega_{B/A, \delta} =
\lim_e \Omega_{B^\wedge/A, \delta^\wedge}/p^e \Omega_{B^\wedge/A, \delta^\wedge}
$$
である。記法と説明については証明を参照されたい。
\end{lemma}

\begin{proof}
『分割冪代数』補題 \ref{dpa-lemma-extend-to-completion} により、
十分大きいすべての $e$ に対して $\delta$ は $B_e = B/p^eB$ に延長される。
したがって最初の極限は意味をもつ。同補題は完備化
$B^\wedge = \lim_e B_e$ 上の分割冪構造 $\delta^\wedge$ も与えるので、
最後の極限も意味をもつ。補題 \ref{crystalline-lemma-omega} と、常に
$\text{d}p^e = 0$ であることから、全射
$\Omega_{B/A, \delta} \to \Omega_{B_e/A, \bar\delta}$ の核は
$p^e\Omega_{B/A, \delta}$ である。同様に、
$\Omega_{B^\wedge/A, \delta^\wedge} \to \Omega_{B_e/A, \bar\delta}$ の
核についても同じである。これで補題は明らかである。
\end{proof}



\section{分割冪スキーム}
\label{crystalline-section-divided-power-schemes}

\noindent
これまでの概念を大域化する方法について、いくつか注意する。

\begin{definition}
\label{crystalline-definition-divided-power-structure}
$\mathcal{C}$ をサイト、$\mathcal{O}$ を $\mathcal{C}$ 上の環の層、
$\mathcal{I} \subset \mathcal{O}$ をイデアルの層とする。
$\mathcal{I}$ 上の {\it 分割冪構造 $\gamma$}とは、写像の列
$\gamma_n : \mathcal{I} \to \mathcal{I}$（$n \geq 1$）であって、
$\mathcal{C}$ の任意の対象 $U$ に対して三つ組
$$
(\mathcal{O}(U), \mathcal{I}(U), \gamma)
$$
が分割冪環となるものをいう。
\end{definition}

\noindent
もちろん、これはとくに位相空間上の環の層に適用できる。しかし少し一般的に
しておくのがよい。結晶サイトの構造層が置かれるのは、まさにサイトだからである。
本章では、上の定義における三つ組 $(\mathcal{C}, \mathcal{I}, \gamma)$ を
{\it 分割冪トポス}と呼ぶことがある。別の三つ組
$(\mathcal{C}', \mathcal{I}', \gamma')$ と環付きトポスの射
$(f, f^\sharp) : (\Sh(\mathcal{C}), \mathcal{O}) \to
(\Sh(\mathcal{C}'), \mathcal{O}')$
が与えられたとする。$f^\sharp(f^{-1}\mathcal{I}') \subset \mathcal{I}$ であり、
すべての $n \geq 1$ に対して図式
$$
\xymatrix{
f^{-1}\mathcal{I}' \ar[d]_{f^{-1}\gamma'_n} \ar[r]_{f^\sharp} &
\mathcal{I} \ar[d]^{\gamma_n} \\
f^{-1}\mathcal{I}' \ar[r]^{f^\sharp} & \mathcal{I}
}
$$
が可換であるとき、$(f, f^\sharp)$ は {\it 分割冪トポスの射}を誘導するという。
$f$ が関手 $u : \mathcal{C}' \to \mathcal{C}$ により誘導されるサイトの射から
来る場合、これは単に、すべての $U' \in \Ob(\mathcal{C}')$ に対して
$$
(\mathcal{O}'(U'), \mathcal{I}'(U'), \gamma')
\longrightarrow
(\mathcal{O}(u(U')), \mathcal{I}(u(U')), \gamma)
$$
が分割冪環の準同型であることを意味する。

\medskip\noindent
スキームの場合には、分割冪イデアルが {\bf 準連接}であることを要求する。
この点を除けば定義はトポスの場合とまったく同じである。以下に定義する。

\begin{definition}
\label{crystalline-definition-divided-power-scheme}
{\it 分割冪スキーム}とは三つ組 $(S, \mathcal{I}, \gamma)$ であって、
$S$ がスキーム、$\mathcal{I}$ が準連接なイデアルの層、$\gamma$ が
$\mathcal{I}$ 上の分割冪構造であるものをいう。
{\it 分割冪スキームの射}
$(S, \mathcal{I}, \gamma) \to (S', \mathcal{I}', \gamma')$ とは、
$f^{-1}\mathcal{I}'\mathcal{O}_S \subset \mathcal{I}$ を満たし、さらに
すべての開集合 $U' \subset S'$ に対して
$$
(\mathcal{O}_{S'}(U'), \mathcal{I}'(U'), \gamma')
\longrightarrow
(\mathcal{O}_S(f^{-1}U'), \mathcal{I}(f^{-1}U'), \gamma)
$$
が分割冪環の準同型となるスキームの射 $f : S \to S'$ である。
\end{definition}

\noindent
準連接なイデアルの層と閉埋め込みの間に一対一対応があることを思い出す。
『射』節 \ref{morphisms-section-closed-immersions} を参照されたい。
したがって分割冪スキーム $(T, \mathcal{J}, \gamma)$ が与えられると、
$\mathcal{J}$ で定義される標準閉埋め込み $U \to T$ を得る。
逆に、閉埋め込み $U \to T$ と、$U \to T$ に付随するイデアルの層 $\mathcal{J}$ 上の
分割冪構造 $\gamma$ が与えられると、分割冪スキーム
$(T, \mathcal{J}, \gamma)$ を得る。多くの状況では、射 $U \to T$ が
肥厚化であるような三つ組 $(U, T, \gamma)$ だけを考えたい。
『射の詳細』定義 \ref{more-morphisms-definition-thickening} を参照されたい。

\begin{definition}
\label{crystalline-definition-divided-power-thickening}
上の三つ組 $(U, T, \gamma)$ で $U \to T$ が肥厚化であるものを
{\it 分割冪肥厚化}という。
\end{definition}

\noindent
三つのうち一つが分割冪肥厚化である場合、分割冪スキームのファイバー積は
存在する。正確な主張は次のとおりである。

\begin{lemma}
\label{crystalline-lemma-fibre-product}
$f : (T, \mathcal{J}, \delta) \to (S, \mathcal{I}, \gamma)$ と
$f' : (T', \mathcal{J}', \delta') \to (S, \mathcal{I}, \gamma)$ を
分割冪スキームの射とする。分割冪スキーム $(T'', \mathcal{J}'', \delta'')$ と、
分割冪スキームの圏におけるカルテシアン図式
$$
\xymatrix{
T \ar[d]_f & T'' \ar[d] \ar[l] \\
S & T' \ar[l]_{f'}
}
$$
が存在する。射 $T'' \to T \times_S T'$ は閉埋め込みであり、射
$T''_0 \to T_0 \times_{S_0} T'_0$ は同型である。
\end{lemma}

\begin{proof}
概略を述べる。最後の二つの主張は、$\Spec(-)$ を介して、
『分割冪代数』注意 \ref{dpa-remark-pushout} で見た分割冪代数の押し出しに関する
事実と両立する。そこで、$T''$ を $T \times_S T'$ の閉部分スキームとして
次のように構成する。$f(V), f'(V') \subset U$ を満たす任意のアフィン開集合
$U \subset S$、$V \subset T$、$V' \subset T'$ に対し、
『分割冪代数』注意 \ref{dpa-remark-pushout} の構成で定まる
$V \times_U V'$ の閉部分スキームを考える。スキーム $V \times_U V'$ は
$T \times_S T'$ の開被覆をなすので、次の手順をとれる：(1) これらの
閉部分スキームが貼り合わさることを示す；(2) 得られた分割冪構造が
貼り合わさることを示す；(3) 貼り合わせた結果が分割冪スキームの圏における
ファイバー積であることを示す。(1) を示すことは、分割冪環の圏における
押し出しの構成が局所化と可換であることを、適切に定式化して示すことに等しい。
これは『分割冪代数』補題 \ref{dpa-lemma-gamma-extends} から従う。同補題により、
平坦性を用いて分割冪が局所化へ延長されるからである。
\end{proof}

\noindent
次のことを観察する。$(U, T, \gamma)$ を分割冪スキームとし、$T$ が
$\mathbf{Z}_{(p)}$ 上のスキームで、$U$ 上 $p$ が局所冪零であると仮定する。
このとき
\begin{enumerate}
\item $p$ が $T \Leftrightarrow U \to T$ 上局所冪零であることは肥厚化である
（『分割冪代数』補題 \ref{dpa-lemma-nil} を参照）。
\item $e \gg 0$ に対して、$p^e\mathcal{O}_T$ は $T$ 上局所的に
$\gamma$ で保たれる（『分割冪代数』補題
\ref{dpa-lemma-extend-to-completion} を参照）。
\end{enumerate}
これは、次の仮定のもとで分割冪肥厚化に関する良い結果が得られることを示唆する。

\begin{situation}
\label{crystalline-situation-global}
ここで $p$ は素数、$(S, \mathcal{I}, \gamma)$ は
$\mathbf{Z}_{(p)}$ 上の分割冪スキームである。
$S_0 = V(\mathcal{I}) \subset S$ とおく。さらに $X \to S_0$ は、
$X$ 上 $p$ が局所冪零となるスキームの射である。
\end{situation}

\noindent
この状況で大結晶サイトと小結晶サイトを定義する。









\section{大結晶サイト}
\label{crystalline-section-big-site}

\noindent
まず大サイトを定義する。分割冪スキーム $(S, \mathcal{I}, \gamma)$ が
与えられたとする。$(T, \mathcal{J}, \delta)$ の $T$ に分割冪スキームの射
$T \to S$ が備わっているとき、これを $(S, \mathcal{I}, \gamma)$ 上の
分割冪スキームという。同様に、分割冪肥厚化 $(U, T, \delta)$ の $T$ に
分割冪スキームの射 $T \to S$ が備わっているとき、これを
$(S, \mathcal{I}, \gamma)$ 上の分割冪肥厚化という。

\begin{definition}
\label{crystalline-definition-divided-power-thickening-X}
状況 \ref{crystalline-situation-global} において、次のように定める。
\begin{enumerate}
\item {\it $(S, \mathcal{I}, \gamma)$ に相対的な $X$ の分割冪肥厚化}とは、
$(S, \mathcal{I}, \gamma)$ 上の分割冪肥厚化 $(U, T, \delta)$ と
$S$-射 $U \to X$ の組である。
\item {\it $(S, \mathcal{I}, \gamma)$ に相対的な $X$ の分割冪肥厚化の射}は
明らかな仕方で定義する。
\end{enumerate}
$(S, \mathcal{I}, \gamma)$ に相対的な $X$ の分割冪肥厚化の圏を
$\text{CRIS}(X/S, \mathcal{I}, \gamma)$、または単に
$\text{CRIS}(X/S)$ と書く。
\end{definition}

\noindent
$\text{CRIS}(X/S)$ の任意の $(U, T, \delta)$ に対して、$T$ 上 $p$ は
局所冪零である。状況 \ref{crystalline-situation-global} の直前の議論を参照されたい。
$(U, T, \delta)$ に付随するすべてのデータは、可換図式
$$
\xymatrix{
T \ar[dd] & U \ar[l] \ar[d] \\
& X \ar[d] \\
S & S_0 \ar[l]
}
$$
で表すと分かりやすい。ここで $S_0 = V(\mathcal{I}) \subset S$ である。
$\text{CRIS}(X/S)$ の射も同様に、大きな可換図式として表せる。
とくに標準忘却関手
\begin{equation}
\label{crystalline-equation-forget}
\text{CRIS}(X/S) \longrightarrow \Sch/X,\quad
(U, T, \delta) \longmapsto U
\end{equation}
と、その片側逆（かつ左随伴）
\begin{equation}
\label{crystalline-equation-endow-trivial}
\Sch/X \longrightarrow \text{CRIS}(X/S),\quad
U \longmapsto (U, U, \emptyset)
\end{equation}
がある。後者はときに有用である。

\begin{lemma}
\label{crystalline-lemma-divided-power-thickening-fibre-products}
状況 \ref{crystalline-situation-global} において、圏 $\text{CRIS}(X/S)$ は
すべての空でない有限極限、とくに二対象の積とファイバー積をもつ。
関手 (\ref{crystalline-equation-forget}) は極限と可換である。
\end{lemma}

\begin{proof}
省略する。ヒント：アフィンの場合については補題
\ref{crystalline-lemma-affine-thickenings-colimits} を参照されたい。
『分割冪代数』注意 \ref{dpa-remark-pushout} も参照されたい。
\end{proof}

\begin{lemma}
\label{crystalline-lemma-divided-power-thickening-base-change-flat}
状況 \ref{crystalline-situation-global} において、
$$
\xymatrix{
(U_3, T_3, \delta_3) \ar[d] \ar[r] & (U_2, T_2, \delta_2) \ar[d] \\
(U_1, T_1, \delta_1) \ar[r] & (U, T, \delta)
}
$$
を $(S, \mathcal{I}, \gamma)$ に相対的な $X$ の分割冪肥厚化の圏における
ファイバー平方とする。$T_2 \to T$ が平坦で $U_2 = T_2 \times_T U$ ならば、
スキームとして $T_3 = T_1 \times_T T_2$ である。
\end{lemma}

\begin{proof}
分割冪構造は平坦な環準同型に沿って一意に延長されるので、主張が成り立つ。
『分割冪代数』補題 \ref{dpa-lemma-gamma-extends} を参照されたい。
\end{proof}

\noindent
上の補題は、分割冪肥厚化の平坦射の基底変換が再び平坦射であり、実際には
その射の「通常の」基底変換であることを意味する。したがって次の定義は意味をもつ。

\begin{definition}
\label{crystalline-definition-big-crystalline-site}
状況 \ref{crystalline-situation-global} において、次のように定める。
\begin{enumerate}
\item $X/S$ の分割冪肥厚化の射の族
$\{(U_i, T_i, \delta_i) \to (U, T, \delta)\}$ が
{\it Zariski、\'etale、smooth、syntomic、または fppf 被覆}であるとは、
次の二条件を満たすことをいう：
\begin{enumerate}
\item すべての $i$ に対して $U_i = U \times_T T_i$ である；
\item $\{T_i \to T\}$ は Zariski、\'etale、smooth、syntomic、または
fppf 被覆である。
\end{enumerate}
\item $(S, \mathcal{I}, \gamma)$ 上の $X$ の {\it 大結晶サイト}とは、
圏 $\text{CRIS}(X/S)$ に Zariski 位相を入れたものである。
\item $\text{CRIS}(X/S)$ 上の層のトポスを $(X/S)_{\text{CRIS}}$、
または $(X/S, \mathcal{I}, \gamma)_{\text{CRIS}}$ と書くことがある
\footnote{これは、サイト $\mathcal{C}$ に付随するトポスを
$\Sh(\mathcal{C})$ と書く本章の慣習と衝突する。}。
\end{enumerate}
\end{definition}

\noindent
これらのトポスには、いくつか明らかな関手性がある。

\begin{remark}[関手性]
\label{crystalline-remark-functoriality-big-cris}
$p$ を素数とし、$(S, \mathcal{I}, \gamma) \to
(S', \mathcal{I}', \gamma')$ を $\mathbf{Z}_{(p)}$ 上の分割冪スキームの射とする。
$S_0 = V(\mathcal{I})$ および $S'_0 = V(\mathcal{I}')$ とおく。
$$
\xymatrix{
X \ar[r]_f \ar[d] & Y \ar[d] \\
S_0 \ar[r] & S'_0
}
$$
をスキームの射の可換図式とし、$X$ と $Y$ 上で $p$ が局所冪零であると
仮定する。このとき、$(U, T, \delta)$ を、$U$ から $Y$ への $S'$-射として
$U \to X \to Y$ を備えた同じ $(U, T, \delta)$ に対応させることにより、
連続かつ余連続な関手
$$
\text{CRIS}(X/S) \longrightarrow \text{CRIS}(Y/S')
$$
を得る。したがってトポスの射
$$
f_{\text{CRIS}} : (X/S)_{\text{CRIS}} \longrightarrow (Y/S')_{\text{CRIS}}
$$
を得る。『サイト』節 \ref{sites-section-cocontinuous-morphism-topoi} を
参照されたい。
\end{remark}

\begin{remark}[Zariski サイトとの比較]
\label{crystalline-remark-compare-big-zariski}
状況 \ref{crystalline-situation-global} において、関手 (\ref{crystalline-equation-forget}) は
余連続であり（詳細は省略する）、積およびファイバー積と可換である
（補題 \ref{crystalline-lemma-divided-power-thickening-fibre-products}）。
したがって、$X/S$ の大結晶トポスから $X$ の大 Zariski トポスへのトポスの射
$$
U_{X/S} : (X/S)_{\text{CRIS}} \longrightarrow \Sh((\Sch/X)_{Zar})
$$
を得る。『サイト』節 \ref{sites-section-cocontinuous-morphism-topoi} を
参照されたい。
\end{remark}

\begin{remark}[構造射]
\label{crystalline-remark-big-structure-morphism}
状況 \ref{crystalline-situation-global} において、閉部分スキーム
$S_0 = V(\mathcal{I}) \subset S$ を考える。$S_0$ 上で $p$ が局所冪零であると
仮定すれば（実際上は常にそうである）、定義
\ref{crystalline-definition-divided-power-thickening-X} で $X$ を $S_0$ に置き換えた状況を得る。
したがってサイト $\text{CRIS}(S_0/S)$ を得る。$f : X \to S_0$ を
$S$ 上の $X$ の構造射とすると、環付きトポスの射の可換図式
$$
\xymatrix{
(X/S)_{\text{CRIS}}
\ar[r]_{f_{\text{CRIS}}} \ar[d]_{U_{X/S}} &
(S_0/S)_{\text{CRIS}} \ar[d]^{U_{S_0/S}} \\
\Sh((\Sch/X)_{Zar}) \ar[r]^{f_{big}} & \Sh((\Sch/S_0)_{Zar}) \ar[rd] \\
& & \Sh((\Sch/S)_{Zar})
}
$$
を注意 \ref{crystalline-remark-functoriality-big-cris} により得る。合成
$(X/S)_{\text{CRIS}} \to \Sh((\Sch/S)_{Zar})$ を大結晶サイトの構造射とみなす。
$S_0$ 上で $p$ が局所冪零でなくても、構造射
$$
(X/S)_{\text{CRIS}} \longrightarrow \Sh((\Sch/S)_{Zar})
$$
は、上の図式の下側の経路をとることにより定義される。すなわちこれは、規則
$(U, T, \delta)/S \mapsto U/S$ で与えられる余連続関手
$\text{CRIS}(X/S) \to (\Sch/S)_{Zar}$ に対応するトポスの射である。
『サイト』節 \ref{sites-section-cocontinuous-morphism-topoi} を参照されたい。
\end{remark}

\begin{remark}[両立性]
\label{crystalline-remark-compatibilities-big-cris}
上で定義した射は多数の両立性を満たす。たとえば注意
\ref{crystalline-remark-functoriality-big-cris} の状況では、環付きトポスの可換図式
$$
\xymatrix{
(X/S)_{\text{CRIS}} \ar[d] \ar[r] & (Y/S')_{\text{CRIS}} \ar[d] \\
\Sh((\Sch/S)_{Zar}) \ar[r] & \Sh((\Sch/S')_{Zar})
}
$$
を得る。ここで垂直矢印は構造射である。
\end{remark}




\section{結晶サイト}
\label{crystalline-section-site}

\noindent
(\ref{crystalline-equation-forget}) は積およびファイバー積と可換なので、
$U \to X$ が開埋め込みであるような $(U, T, \delta)$ をとると、
ファイバー積、より一般には空でない有限極限で保たれる充満部分圏が定まる。
したがって次の定義は意味をもつ。

\begin{definition}
\label{crystalline-definition-crystalline-site}
状況 \ref{crystalline-situation-global} において、次のように定める。
\begin{enumerate}
\item $(S, \mathcal{I}, \gamma)$ 上の $X$ の（小）{\it 結晶サイト}を
$\text{Cris}(X/S, \mathcal{I}, \gamma)$、または単に
$\text{Cris}(X/S)$ と書く。これは $\text{CRIS}(X/S)$ の充満部分圏であり、
$\text{CRIS}(X/S)$ の対象 $(U, T, \delta)$ のうち $U \to X$ が
開埋め込みであるものからなる。これに Zariski 位相を入れる。
\item $\text{Cris}(X/S)$ 上の層のトポスを $(X/S)_{\text{cris}}$、
または $(X/S, \mathcal{I}, \gamma)_{\text{cris}}$ と書くことがある
\footnote{これは、サイト $\mathcal{C}$ に付随するトポスを
$\Sh(\mathcal{C})$ と書く本章の慣習と衝突する。}。
\end{enumerate}
\end{definition}

\noindent
$\text{Cris}(X/S)$ の任意の $(U, T, \delta)$ に対し、射 $U \to X$ は
$X$ の小 Zariski サイト $X_{Zar}$ の対象を定める。したがって標準忘却関手
\begin{equation}
\label{crystalline-equation-forget-small}
\text{Cris}(X/S) \longrightarrow X_{Zar},\quad
(U, T, \delta) \longmapsto U
\end{equation}
と左随伴
\begin{equation}
\label{crystalline-equation-endow-trivial-small}
X_{Zar} \longrightarrow \text{Cris}(X/S),\quad
U \longmapsto (U, U, \emptyset)
\end{equation}
がある。後者はときに有用である。

\medskip\noindent
スキームの小 Zariski サイトと大 Zariski サイトを比較できるのと同様に、
小結晶サイトと大結晶サイトを比較できる。『位相』補題
\ref{topologies-lemma-at-the-bottom} を参照されたい。

\begin{lemma}
\label{crystalline-lemma-compare-big-small}
仮定を定義 \ref{crystalline-definition-divided-power-thickening-X} と同じとする。包含関手
$$
\text{Cris}(X/S) \to \text{CRIS}(X/S)
$$
は空でない有限極限と可換であり、充満忠実、連続、かつ余連続である。
トポスの射
$$
(X/S)_{\text{cris}} \xrightarrow{i} (X/S)_{\text{CRIS}}
\xrightarrow{\pi} (X/S)_{\text{cris}}
$$
があり、その合成は恒等射で、最初の射は包含関手により誘導される。
さらに $\pi_* = i^{-1}$ である。
\end{lemma}

\begin{proof}
最初の主張については補題
\ref{crystalline-lemma-divided-power-thickening-fibre-products} を参照されたい。
これによりトポスの射 $i : (X/S)_{\text{cris}} \to
(X/S)_{\text{CRIS}}$ と、$i^{-1}i_! = i^{-1}i_* = \text{id}$ を満たす
左随伴 $i_!$ を得る。『サイト』補題 \ref{sites-lemma-when-shriek}、
\ref{sites-lemma-preserve-equalizers}、\ref{sites-lemma-back-and-forth} を
参照されたい。$i_!$ は完全であると主張する。これが示されれば、
$\pi^{-1} = i_!$ および $\pi_* = i^{-1}$ により $\pi$ を定義でき、
すべて明らかになる。主張を示すため、$i_!$ が右完全でファイバー積を保つことは
すでに分かっている（上の参照先を見よ）。したがって、$*$ を集合の層の圏の
終対象とするとき $i_! * = *$ を示せば十分である。そのためには、
$\text{Cris}(X/S)$ の対象の集合 $(U_i, T_i, \delta_i)$（$i \in I$）で
$$
\coprod\nolimits_{i \in I} h_{(U_i, T_i, \delta_i)} \to *
$$
が $(X/S)_{\text{CRIS}}$ において全射となるものを構成すればよい
（詳細は省略する。ヒント：$\text{Cris}(X/S)$ が積をもち、関手
$\text{Cris}(X/S) \to \text{CRIS}(X/S)$ が積と可換であることを用いる）。
アフィンの場合、これは補題 \ref{crystalline-lemma-set-generators} から従う。
一般の場合の証明は省略する。
\end{proof}

\begin{remark}[関手性]
\label{crystalline-remark-functoriality-cris}
$p$ を素数とし、$(S, \mathcal{I}, \gamma) \to
(S', \mathcal{I}', \gamma')$ を $\mathbf{Z}_{(p)}$ 上の分割冪スキームの射とする。
$$
\xymatrix{
X \ar[r]_f \ar[d] & Y \ar[d] \\
S_0 \ar[r] & S'_0
}
$$
をスキームの射の可換図式とし、$X$ と $Y$ 上で $p$ が局所冪零であると仮定する。
『位相』補題 \ref{topologies-lemma-morphism-big-small} との類似により
$$
f_{\text{cris}} : (X/S)_{\text{cris}} \longrightarrow (Y/S')_{\text{cris}}
$$
を式 $f_{\text{cris}} = \pi_Y \circ f_{\text{CRIS}} \circ i_X$ で定義する。
ここで $i_X$ と $\pi_Y$ は $X$ と $Y$ に対する補題
\ref{crystalline-lemma-compare-big-small} の射であり、$f_{\text{CRIS}}$ は注意
\ref{crystalline-remark-functoriality-big-cris} の射である。
\end{remark}

\begin{remark}[Zariski サイトとの比較]
\label{crystalline-remark-compare-zariski}
状況 \ref{crystalline-situation-global} において、関手
(\ref{crystalline-equation-forget-small}) は連続かつ余連続で、積およびファイバー積と
可換である。したがってトポスの射
$$
u_{X/S} : (X/S)_{\text{cris}} \longrightarrow \Sh(X_{Zar})
$$
を得る。これは $X/S$ の小結晶トポスと $X$ の小 Zariski トポスを結ぶ。
『サイト』節 \ref{sites-section-cocontinuous-morphism-topoi} を参照されたい。
\end{remark}

\begin{lemma}
\label{crystalline-lemma-localize}
状況 \ref{crystalline-situation-global} において、$X' \subset X$ と $S' \subset S$ を、
$X'$ が $S'$ に写るような開部分スキームとする。このとき充満忠実関手
$\text{Cris}(X'/S') \to \text{Cris}(X/S)$ があり、これは可換図式
$$
\xymatrix{
(X'/S')_{\text{cris}} \ar[r] \ar[d]_{u_{X'/S'}} &
(X/S)_{\text{cris}} \ar[d]^{u_{X/S}} \\
\Sh(X'_{Zar}) \ar[r] & \Sh(X_{Zar})
}
$$
に入るトポスの射を誘導する。さらにこの図式は、『サイト』補題
\ref{sites-lemma-localize-morphism-topoi} におけるトポスの射の局所化の例である。
\end{lemma}

\begin{proof}
充満忠実関手は、$\text{Cris}(X'/S')$ の対象を、$U \to X$ が
$X' \subset X$ を経由する $X$ の分割冪肥厚化 $(U, T, \delta)$ とみなすことで
得られる（このとき $T \to S$ は自動的に $S'$ を経由する）。
この関手は明らかに余連続なので、表示したトポスの射を得る。
$X'$ を $X_{Zar}$ の対象とみなしたときの表現可能層を
$h_{X'} \in \Sh(X_{Zar})$ とする。$\Sh(X'_{Zar})$ は局所化
$\Sh(X_{Zar})/h_{X'}$ である。一方、圏
$\text{Cris}(X/S)/u_{X/S}^{-1}h_{X'}$ は上の関手により標準的に
$\text{Cris}(X'/S')$ と同一視される。『サイト』補題
\ref{sites-lemma-localize-topos-site} を参照されたい。これで証明が完了する。
\end{proof}

\begin{remark}[構造射]
\label{crystalline-remark-structure-morphism}
状況 \ref{crystalline-situation-global} において、閉部分スキーム
$S_0 = V(\mathcal{I}) \subset S$ を考える。$S_0$ 上で $p$ が局所冪零であると
仮定すれば（実際上は常にそうである）、定義
\ref{crystalline-definition-divided-power-thickening-X} で $X$ を $S_0$ に置き換えた状況を得る。
したがってサイト $\text{Cris}(S_0/S)$ を得る。$f : X \to S_0$ を
$S$ 上の $X$ の構造射とすると、トポスの可換図式
$$
\xymatrix{
(X/S)_{\text{cris}}
\ar[r]_{f_{\text{cris}}} \ar[d]_{u_{X/S}} &
(S_0/S)_{\text{cris}} \ar[d]^{u_{S_0/S}} \\
\Sh(X_{Zar}) \ar[r]^{f_{small}} & \Sh(S_{0, Zar}) \ar[rd] \\
& & \Sh(S_{Zar})
}
$$
を得る。注意 \ref{crystalline-remark-functoriality-cris} を参照されたい。合成
$(X/S)_{\text{cris}} \to \Sh(S_{Zar})$ を結晶サイトの構造射とみなす。
$S_0$ 上で $p$ が局所冪零でなくても、構造射
$$
\tau_{X/S} : (X/S)_{\text{cris}} \longrightarrow \Sh(S_{Zar})
$$
は上の図式の下側の経路をとることにより定義される。
\end{remark}

\begin{remark}[両立性]
\label{crystalline-remark-compatibilities}
上で定義した射は多数の両立性を満たす。たとえば注意
\ref{crystalline-remark-functoriality-cris} の状況では、環付きトポスの可換図式
$$
\xymatrix{
(X/S)_{\text{cris}} \ar[d] \ar[r] & (Y/S')_{\text{cris}} \ar[d] \\
\Sh((\Sch/S)_{Zar}) \ar[r] & \Sh((\Sch/S')_{Zar})
}
$$
を得る。ここで垂直矢印は構造射である。
\end{remark}



\section{結晶サイト上の層}
\label{crystalline-section-sheaves}

\noindent
記法と仮定を状況 \ref{crystalline-situation-global} と同じとする。この節では $X/S$ の
小結晶サイトと大結晶サイトを同時に扱うため、
$$
\mathcal{C} = \text{CRIS}(X/S)
\quad\text{または}\quad
\mathcal{C} = \text{Cris}(X/S).
$$
とおく。$\mathcal{C}$ 上の層 $\mathcal{F}$ は、$\mathcal{C}$ の各対象
$(U, T, \delta)$ に対して {\it 制限} $\mathcal{F}_T$ を与える。
すなわち $\mathcal{F}_T$ は、規則
$$
\mathcal{F}_T(W) = \mathcal{F}(U \cap W, W, \delta|_W)
$$
により定義されるスキーム $T$ 上の Zariski 層である。ここで $W \subset T$ は
開集合である。さらに、$f : T \to T'$ が $\mathcal{C}$ の対象
$(U, T, \delta)$ と $(U', T', \delta')$ の間の射ならば、標準的な
{\it 比較}写像
\begin{equation}
\label{crystalline-equation-comparison}
c_f : f^{-1}\mathcal{F}_{T'} \longrightarrow \mathcal{F}_T.
\end{equation}
がある。実際、$W' \subset T'$ が開ならば $f$ は $\mathcal{C}$ の射
$$
f|_{f^{-1}W'} :
(U \cap f^{-1}(W'), f^{-1}W', \delta|_{f^{-1}W'})
\longrightarrow
(U' \cap W', W', \delta|_{W'})
$$
を誘導する。したがって $\mathcal{F}$ の制限写像
$(f|_{f^{-1}W'})^*$ を用いて写像
$\mathcal{F}_{T'}(W') \to \mathcal{F}_T(f^{-1}W')$ を定義できる。
これらの写像はさらなる制限と明らかに両立するので、$\mathcal{F}_{T'}$ から
$\mathcal{F}_T$ への $f$-写像を定める。『層』節
\ref{sheaves-section-presheaves-functorial}、とくに『層』定義
\ref{sheaves-definition-f-map} を参照されたい。こうして
(\ref{crystalline-equation-comparison}) の写像 $c_f$ が得られる。$f$ が開埋め込みならば
$c_f$ は同型であることに注意する。この場合 $\mathcal{F}_T$ は単に
$\mathcal{F}_{T'}$ の $T$ への制限だからである。

\medskip\noindent
逆に、$\mathcal{C}$ の各対象 $(U, T, \delta)$ に対する Zariski 層
$\mathcal{F}_T$ と上の比較写像 $c_f$ が与えられ、(a) 開埋め込みに対して
同型であり、(b) 適切なコサイクル条件を満たすならば、$\mathcal{C}$ 上の層を得る。
証明は『位相』補題 \ref{topologies-lemma-characterize-sheaf-big} とまったく同じである。

\medskip\noindent
$\mathcal{C}$ 上の {\it 構造層}とは、規則
$$
\mathcal{O}_{X/S} :
(U, T, \delta)
\longmapsto
\Gamma(T, \mathcal{O}_T)
$$
で定義される層 $\mathcal{O}_{X/S}$ である。$\mathcal{C}$ における被覆の定義から
これは層である。$\mathcal{F}$ を $\mathcal{O}_{X/S}$-加群の層とする。
この場合、比較写像 (\ref{crystalline-equation-comparison}) は $\mathcal{O}_T$-加群の比較写像
\begin{equation}
\label{crystalline-equation-comparison-modules}
c_f : f^*\mathcal{F}_{T'} \longrightarrow \mathcal{F}_T
\end{equation}
を定める。

\medskip\noindent
別種の例は、$(\Sch/X)_{Zar}$ または $X_{Zar}$ 上の層 $\mathcal{G}$ から得られる
（どちらを用いるかは $\mathcal{C} = \text{CRIS}(X/S)$ か
$\mathcal{C} = \text{Cris}(X/S)$ かによる）。規則
$$
\underline{\mathcal{G}} :
(U, T, \delta)
\longmapsto
\mathcal{G}(U)
$$
で定義される $\underline{\mathcal{G}}$ は $\mathcal{C}$ 上の層である。
とくに $\mathcal{G} = \mathbf{G}_a = \mathcal{O}_X$ とすれば
$$
\underline{\mathbf{G}_a} :
(U, T, \delta)
\longmapsto
\Gamma(U, \mathcal{O}_U)
$$
を得る。各対象 $(U, T, \delta)$ に対する標準写像
$\Gamma(T, \mathcal{O}_T) \to \Gamma(U, \mathcal{O}_U)$ により、層の全射
$\mathcal{O}_{X/S} \to \underline{\mathbf{G}_a}$ が定まる。その核を
$\mathcal{J}_{X/S}$ と書くと、短完全列
$$
0 \to
\mathcal{J}_{X/S} \to
\mathcal{O}_{X/S} \to
\underline{\mathbf{G}_a} \to 0
$$
を得る。$\mathcal{J}_{X/S}$ には標準的な分割冪構造が備わることに注意する。
実際、各対象 $(U, T, \delta)$ の第三成分 $\delta$ は、まさに
$\mathcal{O}_T \to \mathcal{O}_U$ の核上の分割冪構造である。
したがって（大）結晶トポスは分割冪トポスである。





\section{加群のクリスタル}
\label{crystalline-section-crystals}

\noindent
クリスタルは非常に一般的な対象である。しかし定義は少し読みにくいため、
まず結晶サイト上の加群の場合に定義を与える。

\begin{definition}
\label{crystalline-definition-modules}
状況 \ref{crystalline-situation-global} において、
$\mathcal{C} = \text{CRIS}(X/S)$ または $\mathcal{C} = \text{Cris}(X/S)$ とし、
$\mathcal{F}$ を $\mathcal{C}$ 上の $\mathcal{O}_{X/S}$-加群の層とする。
\begin{enumerate}
\item $\mathcal{C}$ の各対象 $(U, T, \delta)$ に対して制限 $\mathcal{F}_T$ が
準連接 $\mathcal{O}_T$-加群であるとき、$\mathcal{F}$ は
{\it 局所準連接}であるという。
\item $\mathcal{F}$ が『サイト上の加群』定義
\ref{sites-modules-definition-site-local} の意味で準連接であるとき、
{\it 準連接}であるという。
\item すべての比較写像 (\ref{crystalline-equation-comparison-modules}) が同型であるとき、
$\mathcal{F}$ は {\it $\mathcal{O}_{X/S}$-加群のクリスタル}であるという。
\end{enumerate}
\end{definition}

\noindent
これらの概念は次のように関係する。

\begin{lemma}
\label{crystalline-lemma-crystal-quasi-coherent-modules}
記法 $X/S, \mathcal{I}, \gamma, \mathcal{C}, \mathcal{F}$ を定義
\ref{crystalline-definition-modules} と同じとする。次は同値である。
\begin{enumerate}
\item $\mathcal{F}$ は準連接である。
\item $\mathcal{F}$ は局所準連接であり、$\mathcal{O}_{X/S}$-加群の
クリスタルである。
\end{enumerate}
\end{lemma}

\begin{proof}
(1) を仮定する。$f : (U', T', \delta') \to (U, T, \delta)$ を
$\mathcal{C}$ の対象とする。(a) $\mathcal{F}_T$ が準連接
$\mathcal{O}_T$-加群であり、(b) $c_f : f^*\mathcal{F}_T \to \mathcal{F}_{T'}$
が同型であることを示せばよい。仮定により被覆
$\{(T_i, U_i, \delta_i) \to (T, U, \delta)\}$ を選べ、各 $i$ に対して
$\mathcal{F}$ の $\mathcal{C}/(T_i, U_i, \delta_i)$ への制限は大域表示をもつ。
(a) と (b) は Zariski 局所的に示せば十分なので、
$f : (T', U', \delta') \to (T, U, \delta)$ を
$(T_i, U_i, \delta_i)$ への基底変換で置き換え、$\mathcal{F}$ の
$\mathcal{C}/(T, U, \delta)$ への制限が大域表示
$$
\bigoplus\nolimits_{j \in J}
\mathcal{O}_{X/S}|_{\mathcal{C}/(U, T, \delta)} \longrightarrow
\bigoplus\nolimits_{i \in I}
\mathcal{O}_{X/S}|_{\mathcal{C}/(U, T, \delta)} \longrightarrow
\mathcal{F}|_{\mathcal{C}/(U, T, \delta)}
\longrightarrow 0
$$
をもつと仮定してよい。これが表示
$$
\bigoplus\nolimits_{j \in J} \mathcal{O}_T \longrightarrow
\bigoplus\nolimits_{i \in I} \mathcal{O}_T \longrightarrow
\mathcal{F}_T
\longrightarrow 0
$$
を与えることは明らかなので (a) が成り立つ。さらにこの表示を $T'$ に制限すると
$\mathcal{F}_{T'}$ の同様の表示が得られるので、(b) も成り立つ。

\medskip\noindent
(2) を仮定し、$(U, T, \delta)$ を $\mathcal{C}$ の対象とする。
被覆の各元で $\mathcal{C}$ を局所化した圏に制限すると $\mathcal{F}$ が
大域表示をもつような $(U, T, \delta)$ の被覆を見つければよい。
したがって $T$ はアフィンと仮定してよい。この場合、表示
$$
\bigoplus\nolimits_{j \in J} \mathcal{O}_T \longrightarrow
\bigoplus\nolimits_{i \in I} \mathcal{O}_T \longrightarrow
\mathcal{F}_T
\longrightarrow 0
$$
を選べる。実際、$\mathcal{F}_T$ は準連接 $\mathcal{O}_T$-加群と仮定している。
$\mathcal{F}$ のクリスタル性により、これは $\mathcal{C}$ の任意の射
$f : (U', T', \delta') \to (U, T, \delta)$ に沿って引き戻すと
$\mathcal{F}_{T'}$ の表示になる。したがって
$\mathcal{F}|_{\mathcal{C}/(U, T, \delta)}$ の所望の表示が得られる。
\end{proof}

\begin{definition}
\label{crystalline-definition-crystal-quasi-coherent-modules}
$\mathcal{F}$ が補題 \ref{crystalline-lemma-crystal-quasi-coherent-modules} の同値な条件を
満たすとき、$\mathcal{F}$ は {\it 準連接加群のクリスタル}であるという。
さらに $\mathcal{F}$ が有限局所自由ならば、$\mathcal{F}$ は
{\it 有限局所自由加群のクリスタル}であるという。
\end{definition}

\noindent
もちろん補題 \ref{crystalline-lemma-crystal-quasi-coherent-modules} が示すように、
準連接加群は常にクリスタルなので、この用語はやや冗長である。
しかし文献では標準的な用語である。

\begin{remark}
\label{crystalline-remark-crystal}
クリスタルの一般概念を定式化するため、スタックと強カルテシアン射の言葉を用いる。
『スタック』定義 \ref{stacks-definition-stack} および『圏論』定義
\ref{categories-definition-cartesian-over-C} を参照されたい。
状況 \ref{crystalline-situation-global} において、
$p : \mathcal{C} \to \text{Cris}(X/S)$ をスタックとする。
{\it $S$ に相対的な $X$ 上の $\mathcal{C}$ の対象のクリスタル}とは、
{\it カルテシアン切断} $\sigma : \text{Cris}(X/S) \to \mathcal{C}$、
すなわち $p \circ \sigma = \text{id}$ を満たし、$\text{Cris}(X/S)$ の
すべての射 $f$ に対して $\sigma(f)$ が強カルテシアンであるような関手
$\sigma$ である。大結晶サイトについても同様である。
\end{remark}





\section{微分の層}
\label{crystalline-section-differentials-sheaf}

\noindent
この節では、その方が自然と思われるので（小）結晶サイトに限定する。
定義 \ref{crystalline-definition-derivation} を次のように大域化する。

\begin{definition}
\label{crystalline-definition-global-derivation}
状況 \ref{crystalline-situation-global} において、$\mathcal{F}$ を
$\text{Cris}(X/S)$ 上の $\mathcal{O}_{X/S}$-加群の層とする。
{\it $S$-導分 $D : \mathcal{O}_{X/S} \to \mathcal{F}$}とは、
$\text{Cris}(X/S)$ の各対象 $(U, T, \delta)$ に対して写像
$$
D : \Gamma(T, \mathcal{O}_T) \longrightarrow \Gamma(T, \mathcal{F})
$$
が分割冪 $\Gamma(V, \mathcal{O}_V)$-導分となるような層の写像をいう。
ここで $V \subset S$ は $T \to S$ が $V$ を経由して分解する任意の開部分である。
\end{definition}

\noindent
これは、$D$ が加法的で Leibniz 則を満たし、$S$ から来る関数を零に写し、
分割冪イデアル $\mathcal{J}_{X/S}$ の局所切断 $f$ に対して
$D(f^{[n]}) = f^{[n - 1]}D(f)$ を満たすことを意味する。
これは、これから述べる非常に一般的な概念の特別な場合である。

\medskip\noindent
以下の議論を『サイト上の加群』節
\ref{sites-modules-section-differentials} と比較されたい。
$\mathcal{C}$ をサイト、$\mathcal{A} \to \mathcal{B}$ を
$\mathcal{C}$ 上の環の層の写像、$\mathcal{J} \subset \mathcal{B}$ を
イデアルの層、$\delta$ を $\mathcal{J}$ 上の分割冪構造、
$\mathcal{F}$ を $\mathcal{B}$-加群の層とする。このとき
{\it 分割冪 $\mathcal{A}$-導分}
$D : \mathcal{B} \to \mathcal{F}$ という概念がある。
これは、$D$ が $\mathcal{A}$-線形で Leibniz 則を満たし、
$\mathcal{J}$ の局所切断 $x$ に対して
$D(\delta_n(x)) = \delta_{n - 1}(x)D(x)$ を満たすことを意味する。
この状況では {\it 普遍分割冪 $\mathcal{A}$-導分}
$$
\text{d}_{\mathcal{B}/\mathcal{A}, \delta} :
\mathcal{B}
\longrightarrow
\Omega_{\mathcal{B}/\mathcal{A}, \delta}
$$
が存在する。さらに $\text{d}_{\mathcal{B}/\mathcal{A}, \delta}$ は合成
$$
\mathcal{B}
\longrightarrow
\Omega_{\mathcal{B}/\mathcal{A}}
\longrightarrow
\Omega_{\mathcal{B}/\mathcal{A}, \delta}
$$
である。ここで第一の写像は『サイト上の加群』補題
\ref{sites-modules-lemma-universal-module} の証明で構成された普遍導分であり、
第二の矢印は局所切断
$\text{d}_{\mathcal{B}/\mathcal{A}}(\delta_n(x)) -
\delta_{n - 1}(x)\text{d}_{\mathcal{B}/\mathcal{A}}(x)$
が生成する部分加群による商である。

\medskip\noindent
これを次のような相対的概念に置き換える。環付きトポスの射
$(f, f^\sharp) : (\Sh(\mathcal{C}), \mathcal{O}) \to
(\Sh(\mathcal{C}'), \mathcal{O}')$、イデアルの層
$\mathcal{J} \subset \mathcal{O}$、$\mathcal{J}$ 上の分割冪構造
$\delta$、および $\mathcal{O}$-加群の層 $\mathcal{F}$ が与えられたとする。
この状況で、$D$ が上で定義した分割冪 $f^{-1}\mathcal{O}'$-導分であるとき、
$D : \mathcal{O} \to \mathcal{F}$ を分割冪 $\mathcal{O}'$-導分という。
さらに
$$
\Omega_{\mathcal{O}/\mathcal{O}', \delta} =
\Omega_{\mathcal{O}/f^{-1}\mathcal{O}', \delta}
$$
と書く。これは普遍分割冪 $\mathcal{O}'$-導分の受け皿である。

\medskip\noindent
これを構造射
$$
(X/S)_{\text{Cris}} \longrightarrow \Sh(S_{Zar})
$$
（注意 \ref{crystalline-remark-structure-morphism} を参照）に適用すると、
上の定義 \ref{crystalline-definition-global-derivation} の概念が得られる。
特に普遍分割冪導分
$$
d_{X/S} : \mathcal{O}_{X/S} \to \Omega_{X/S}
$$
が存在する。微分加群が分割冪と両立することを示す装飾は記法から省略する。
結晶サイト上の構造層の通常の微分加群を考えることは、まずないと思われるからである。

\begin{lemma}
\label{crystalline-lemma-module-differentials-divided-power-scheme}
$(T, \mathcal{J}, \delta)$ を分割冪スキームとし、
$T \to S$ をスキームの射とする。上で記述した商
$\Omega_{T/S} \to \Omega_{T/S, \delta}$ は準連接
$\mathcal{O}_T$-加群である。アフィン開集合 $W \subset T$ が
アフィン開集合 $V \subset S$ に写るならば
$$
\Gamma(W, \Omega_{T/S, \delta}) =
\Omega_{\Gamma(W, \mathcal{O}_W)/\Gamma(V, \mathcal{O}_V), \delta}
$$
である。右辺は節 \ref{crystalline-section-differentials} で構成したものである。
\end{lemma}

\begin{proof}
省略する。
\end{proof}

\begin{lemma}
\label{crystalline-lemma-module-of-differentials}
状況 \ref{crystalline-situation-global} において、$\text{Cris}(X/S)$ の
$(U, T, \delta)$ に対し、$T$ への制限 $(\Omega_{X/S})_T$ は
$\Omega_{T/S, \delta}$ であり、制限 $\text{d}_{X/S}|_T$ は
$\text{d}_{T/S, \delta}$ に等しい。
\end{lemma}

\begin{proof}
省略する。
\end{proof}

\begin{lemma}
\label{crystalline-lemma-module-of-differentials-on-affine}
状況 \ref{crystalline-situation-global} において、$\text{Cris}(X/S)$ の任意の
アフィン対象 $(U, T, \delta)$ がアフィン開集合 $V \subset S$ に写るならば
$$
\Gamma((U, T, \delta), \Omega_{X/S}) =
\Omega_{\Gamma(T, \mathcal{O}_T)/\Gamma(V, \mathcal{O}_V), \delta}
$$
である。右辺は節 \ref{crystalline-section-differentials} で構成したものである。
\end{lemma}

\begin{proof}
補題 \ref{crystalline-lemma-module-differentials-divided-power-scheme} と
\ref{crystalline-lemma-module-of-differentials} を組み合わせればよい。
\end{proof}

\begin{lemma}
\label{crystalline-lemma-describe-omega-small}
状況 \ref{crystalline-situation-global} において、$(U, T, \delta)$ を
$\text{Cris}(X/S)$ の対象とする。
$$
(U(1), T(1), \delta(1)) = (U, T, \delta) \times (U, T, \delta)
$$
を $\text{Cris}(X/S)$ における積とする。
$\mathcal{K} \subset \mathcal{O}_{T(1)}$ を閉埋め込み
$\Delta : T \to T(1)$ に対応する準連接イデアルの層とする。
このとき $\mathcal{K} \subset \mathcal{J}_{T(1)}$ は
$\mathcal{J}_{T(1)}$ 上の分割冪構造で保たれ、
$$
(\Omega_{X/S})_T = \mathcal{K}/\mathcal{K}^{[2]}
$$
である。
\end{lemma}

\begin{proof}
$U \to X$ は開埋め込みであり、(\ref{crystalline-equation-forget-small}) は積と可換なので
$U = U(1)$ である。したがって
$\mathcal{K} \subset \mathcal{J}_{T(1)}$ が分かる。
この事実のもとで $T$ 上アフィン局所的に議論し、補題
\ref{crystalline-lemma-module-of-differentials-on-affine} と
\ref{crystalline-lemma-diagonal-and-differentials-affine-site} を用いれば主張が従う。
\end{proof}

\noindent
$\Omega_{X/S}$ は準連接 $\mathcal{O}_{X/S}$-加群のクリスタルではない。
しかし、これと密接に関係する次の二つの性質を満たす
（補題 \ref{crystalline-lemma-crystal-quasi-coherent-modules} と比較されたい）。

\begin{lemma}
\label{crystalline-lemma-omega-locally-quasi-coherent}
状況 \ref{crystalline-situation-global} において、微分の層 $\Omega_{X/S}$ は
次の二つの性質をもつ。
\begin{enumerate}
\item $\Omega_{X/S}$ は局所準連接である。
\item $\text{Cris}(X/S)$ の任意の射
$(U, T, \delta) \to (U', T', \delta')$ で
$f : T \to T'$ が閉埋め込みであるものに対し、写像
$c_f : f^*(\Omega_{X/S})_{T'} \to (\Omega_{X/S})_T$ は全射である。
\end{enumerate}
\end{lemma}

\begin{proof}
部分 (1) は補題 \ref{crystalline-lemma-module-differentials-divided-power-scheme} と
\ref{crystalline-lemma-module-of-differentials} を組み合わせれば従う。
部分 (2) は
$(\Omega_{X/S})_T = \Omega_{T/S, \delta}$
が $\Omega_{T/S}$ の商であり、かつ
$f^*\Omega_{T'/S} \to \Omega_{T/S}$ が全射であることから従う。
\end{proof}







\section{二つの普遍肥厚化}
\label{crystalline-section-universal-thickenings}

\noindent
この節の構成は、結晶サイト上の加群のクリスタルに接続を定義するために用いる。
ここでの構成は、ある意味で節 \ref{crystalline-section-explicit-thickenings} の構成を
「層化し、普遍化した」ものである。

\begin{remark}
\label{crystalline-remark-first-order-thickening}
状況 \ref{crystalline-situation-global} において、$(U, T, \delta)$ を
$\text{Cris}(X/S)$ の対象とする。補題
\ref{crystalline-lemma-module-of-differentials} に従い
$\Omega_{T/S, \delta} = (\Omega_{X/S})_T$ と書く。
$T$ の一次肥厚化 $T'$ を明示的に記述する。すなわち
$$
\mathcal{O}_{T'} = \mathcal{O}_T \oplus \Omega_{T/S, \delta}
$$
とおき、$\Omega_{T/S, \delta}$ が平方零イデアルとなるような代数構造を入れる。
$\mathcal{J} \subset \mathcal{O}_T$ を閉埋め込み $U \to T$ の
イデアルの層とし、
$\mathcal{J}' = \mathcal{J} \oplus \Omega_{T/S, \delta}$ とおく。
$\mathcal{J}'$ 上の分割冪構造を
$$
\delta_n'(f, \omega) = (\delta_n(f), \delta_{n - 1}(f)\omega),
$$
で定義する。補題 \ref{crystalline-lemma-divided-power-first-order-thickening} を参照されたい。
二つの環準同型
$$
p_0, p_1 : \mathcal{O}_T \to \mathcal{O}_{T'}
$$
がある。第一のものは $f \mapsto (f, 0)$、第二のものは
$f \mapsto (f, \text{d}_{T/S, \delta}f)$ で与えられる。
いずれも $\mathcal{J}$ と $\mathcal{J}'$ 上の分割冪構造と両立し、
商写像 $\mathcal{O}_{T'} \to \mathcal{O}_T$ も両立する。
したがって $\text{Cris}(X/S)$ の対象 $(U, T', \delta')$ と可換図式
$$
\xymatrix{
& T \ar[ld]_{\text{id}} \ar[d]^i \ar[rd]^{\text{id}} \\
T & T' \ar[l]_{p_0} \ar[r]^{p_1} & T
}
$$
を $\text{Cris}(X/S)$ において得る。ここで $i$ は、そのイデアルの層が $\Omega_{T/S, \delta}$ と
同一視される一次肥厚化であり、
$p_1 - p_0 : \mathcal{O}_T \to \mathcal{O}_{T'}$ は
普遍導分 $\text{d}_{T/S, \delta}$ と包含
$\Omega_{T/S, \delta} \to \mathcal{O}_{T'}$ の合成に同一視される。
\end{remark}

\begin{remark}
\label{crystalline-remark-second-order-thickening}
状況 \ref{crystalline-situation-global} において、$(U, T, \delta)$ を
$\text{Cris}(X/S)$ の対象とする。補題
\ref{crystalline-lemma-module-of-differentials} に従い
$\Omega_{T/S, \delta} = (\Omega_{X/S})_T$ と書き、その第二外冪を
$\Omega^2_{T/S, \delta}$ とも書く。$T$ の二次肥厚化 $T''$ を
明示的に記述する。すなわち
$$
\mathcal{O}_{T''} =
\mathcal{O}_T \oplus \Omega_{T/S, \delta} \oplus \Omega_{T/S, \delta}
\oplus \Omega^2_{T/S, \delta}
$$
とおき、代数構造を次で定義する。
$$
(f, \omega_1, \omega_2, \eta) \cdot
(f', \omega_1', \omega_2', \eta') =
(ff', f\omega_1' + f'\omega_1, f\omega_2' + f'\omega_2,
f\eta' + f'\eta + \omega_1 \wedge \omega_2' + \omega_1' \wedge \omega_2).
$$
$\mathcal{J} \subset \mathcal{O}_T$ を閉埋め込み $U \to T$ の
イデアルの層とし、$\mathcal{J}''$ を射影
$\mathcal{O}_{T''} \to \mathcal{O}_T$ による $\mathcal{J}$ の逆像とする。
$\mathcal{J}''$ 上の分割冪構造を
$$
\delta_n''(f, \omega_1, \omega_2, \eta) =
(\delta_n(f), \delta_{n - 1}(f)\omega_1, \delta_{n - 1}(f)\omega_2,
\delta_{n - 1}(f)\eta + \delta_{n - 2}(f)\omega_1 \wedge \omega_2)
$$
で定義する。補題 \ref{crystalline-lemma-divided-power-second-order-thickening} を参照されたい。
三つの環準同型
$q_0, q_1, q_2 : \mathcal{O}_T \to \mathcal{O}_{T''}$
を
\begin{align*}
q_0(f) & = (f, 0, 0, 0), \\
q_1(f) & = (f, \text{d}f, 0, 0), \\
q_2(f) & = (f, \text{d}f, \text{d}f, 0)
\end{align*}
で定める。ここで $\text{d} = \text{d}_{T/S, \delta}$ である。
三つとも $\mathcal{J}$ と $\mathcal{J}''$ 上の分割冪構造と両立する。
また三つの環準同型
$q_{01}, q_{12}, q_{02} : \mathcal{O}_{T'} \to \mathcal{O}_{T''}$
がある。ここで $\mathcal{O}_{T'}$ は注意
\ref{crystalline-remark-first-order-thickening} のものである。すなわち
\begin{align*}
q_{01}(f, \omega) & = (f, \omega, 0, 0), \\
q_{12}(f, \omega) & =
(f, \text{d}f, \omega, \text{d}\omega), \\
q_{02}(f, \omega) & = (f, \omega, \omega, 0)
\end{align*}
とおく。これらも与えられた分割冪構造と両立する。
$q_{12}$ について検証しよう。次の計算により $q_{12}$ は環準同型である。
\begin{align*}
q_{12}(f, \omega)q_{12}(g, \eta) & =
(f, \text{d}f, \omega, \text{d}\omega)(g, \text{d}g, \eta, \text{d}\eta) \\
& =
(fg, f\text{d}g + g \text{d}f, f\eta + g\omega,
f\text{d}\eta + g\text{d}\omega + \text{d}f \wedge \eta +
\text{d}g \wedge \omega) \\
& = q_{12}(fg, f\eta + g\omega) = q_{12}((f, \omega)(g, \eta))
\end{align*}
また次の計算により $q_{12}$ は分割冪と両立する。
\begin{align*}
\delta_n''(q_{12}(f, \omega)) & =
\delta_n''((f, \text{d}f, \omega, \text{d}\omega)) \\
& =
(\delta_n(f), \delta_{n - 1}(f)\text{d}f, \delta_{n - 1}(f)\omega,
\delta_{n - 1}(f)\text{d}\omega + \delta_{n - 2}(f)\text{d}(f) \wedge \omega)
\\
& = q_{12}((\delta_n(f), \delta_{n - 1}(f)\omega)) =
q_{12}(\delta'_n(f, \omega))
\end{align*}
$q_{01}$ と $q_{02}$ の検証はより容易である。
$q_0 = q_{01} \circ p_0$、$q_1 = q_{01} \circ p_1$、
$q_1 = q_{12} \circ p_0$, $q_2 = q_{12} \circ p_1$,
$q_0 = q_{02} \circ p_0$、および $q_2 = q_{02} \circ p_1$ である。
したがって $(U, T'', \delta'')$ は $\text{Cris}(X/S)$ の対象であり、
射
$$
\xymatrix{
T''
\ar@<2ex>[r]
\ar@<0ex>[r]
\ar@<-2ex>[r]
&
T'
\ar@<1ex>[r]
\ar@<-1ex>[r]
&
T
}
$$
を得る。これらは $\text{Cris}(X/S)$ の射であり、上で述べた関係を満たす。
応用では、上で記述した環準同型に対応する射を
$q_i : T'' \to T$ および $q_{ij} : T'' \to T'$ と書く。
\end{remark}






\section{de Rham 複体}
\label{crystalline-section-de-Rham}

\noindent
状況 \ref{crystalline-situation-global} において（小）結晶サイト上で考え、
$\Omega^i_{X/S} = \wedge^i_{\mathcal{O}_{X/S}} \Omega_{X/S}$
を $i \geq 0$ に対して定義する。普遍 $S$-導分
$\text{d}_{X/S}$ から、$\text{Cris}(X/S)$ 上の
{\it de Rham 複体}
$$
\mathcal{O}_{X/S} \to \Omega^1_{X/S} \to \Omega^2_{X/S} \to \ldots
$$
が得られる。補題 \ref{crystalline-lemma-module-of-differentials-on-affine} および
注意 \ref{crystalline-remark-divided-powers-de-rham-complex} を参照されたい。


\section{接続}
\label{crystalline-section-connections}

\noindent
状況 \ref{crystalline-situation-global} において、$\text{Cris}(X/S)$ 上の
$\mathcal{O}_{X/S}$-加群 $\mathcal{F}$ が与えられたとする。
{\it 接続}とは、アーベル層の写像
$$
\nabla :
\mathcal{F}
\longrightarrow
\mathcal{F} \otimes_{\mathcal{O}_{X/S}} \Omega_{X/S}
$$
であって、$\mathcal{F}$ と $\mathcal{O}_{X/S}$ の局所切断 $s,f$ に対し
$\nabla(f s) = f\nabla(s) + s \otimes \text{d}f$ を満たすものをいう。
接続が与えられると、標準的な写像
$
\nabla :
\mathcal{F} \otimes_{\mathcal{O}_{X/S}} \Omega^i_{X/S}
\longrightarrow
\mathcal{F} \otimes_{\mathcal{O}_{X/S}} \Omega^{i + 1}_{X/S}
$
が、規則 $\nabla(s \otimes \omega) =
\nabla(s) \wedge \omega + s \otimes \text{d}\omega$
により、注意 \ref{crystalline-remark-connection} と同様に定まる。
$\nabla \circ \nabla = 0$ のとき接続は {\it 可積分}であるという。
$\nabla$ が可積分ならば、$\text{Cris}(X/S)$ 上の
{\it de Rham 複体}
$$
\mathcal{F} \to
\mathcal{F} \otimes_{\mathcal{O}_{X/S}} \Omega^1_{X/S} \to
\mathcal{F} \otimes_{\mathcal{O}_{X/S}} \Omega^2_{X/S} \to \ldots
$$
を得る。任意の $\mathcal{O}_{X/S}$-加群のクリスタルには、
標準的な可積分接続が備わる。

\begin{lemma}
\label{crystalline-lemma-automatic-connection}
状況 \ref{crystalline-situation-global} において、$\mathcal{F}$ を
$\text{Cris}(X/S)$ 上の $\mathcal{O}_{X/S}$-加群のクリスタルとする。
このとき $\mathcal{F}$ には標準的な可積分接続が備わる。
\end{lemma}

\begin{proof}
$(U, T, \delta)$ を $\text{Cris}(X/S)$ の対象とする。
$(U, T', \delta')$ を、注意 \ref{crystalline-remark-first-order-thickening} で
構成した $(\Omega_{X/S})_T = \Omega_{T/S, \delta}$ による
$T$ の無限小肥厚化とする。これには射影
$p_0, p_1 : T' \to T$ と対角射 $i : T \to T'$ が備わる。
仮定により、$\mathcal{O}_{T'}$-加群の同型
$$
p_0^*\mathcal{F}_T \xrightarrow{c_0}
\mathcal{F}_{T'} \xleftarrow{c_1}
p_1^*\mathcal{F}_T
$$
を得る。$c = c_1^{-1} \circ c_0$ を $i$ により $T$ へ引き戻すと、
$\mathcal{F}_T$ の恒等写像が得られる。したがって
$s \in \Gamma(T, \mathcal{F}_T)$ に対し
$\nabla(s) = p_1^*s - c(p_0^*s)$ は $p_1^*\mathcal{F}_T$ の切断であり、
$i$ による引き戻しで零になる。ゆえに $\nabla(s)$ は
$$
\mathcal{F}_T
\otimes_{\mathcal{O}_T}
\Omega_{T/S, \delta}
$$
の切断である。実際、構成により
$\mathcal{O}_{T'} = \mathcal{O}_T \oplus \Omega_{T/S, \delta}$ なので、
これは $p_1^*\mathcal{F}_T \to \mathcal{F}_T$ の核である。
注意 \ref{crystalline-remark-first-order-thickening} の $\text{d}$ の記述を用いれば
$\nabla(fs) = f\nabla(s) + s \otimes \text{d}(f)$ は容易に確かめられる。

\medskip\noindent
こうして得られた写像の族
$$
\nabla : \Gamma(T, \mathcal{F}_T) \to
\Gamma(T, \mathcal{F}_T \otimes_{\mathcal{O}_T} \Omega_{T/S, \delta})
$$
は、$T'$ の構成が $T$ に関して関手的なので $T$ に関して関手的である。
したがって接続を得る。

\medskip\noindent
接続が可積分であることを示すため、注意
\ref{crystalline-remark-second-order-thickening} で構成した対象
$(U, T'', \delta'')$ を考える。$\mathcal{F}$ は層なので
$$
\xymatrix{
q_0^*\mathcal{F}_T \ar[rr]_{q_{01}^*c} \ar[rd]_{q_{02}^*c} & &
q_1^*\mathcal{F}_T \ar[ld]^{q_{12}^*c} \\
& q_2^*\mathcal{F}_T
}
$$
は $\mathcal{O}_{T''}$-加群の可換図式である。
$s \in \Gamma(T, \mathcal{F}_T)$ に対して
$c(p_0^*s) = p_1^*s - \nabla(s)$ である。
$s_i$ を $\mathcal{F}_T$ の局所切断、$\omega_i$ を
$\Omega_{T/S, \delta}$ の局所切断として
$\nabla(s) = \sum p_1^*s_i \cdot \omega_i$ と書く。
$\omega_i$ を $\mathcal{O}_{T'}$ の局所切断とみなすので、
テンソル積の代わりに積を記す。一方で
\begin{align*}
q_{12}^*c \circ q_{01}^*c(q_0^*s) & = 
q_{12}^*c(q_1^*s - \sum q_1^*s_i \cdot q_{01}^*\omega_i) \\
& =
q_2^*s - \sum q_2^*s_i \cdot q_{12}^*\omega_i -
\sum q_2^*s_i \cdot q_{01}^*\omega_i +
\sum q_{12}^*\nabla(s_i) \cdot q_{01}^*\omega_i
\end{align*}
であり、他方で
$$
q_{02}^*c(q_0^*s) = q_2^*s - \sum q_2^*s_i \cdot q_{02}^*\omega_i.
$$
注意 \ref{crystalline-remark-second-order-thickening} の公式から
$q_{01}^*\omega_i + q_{12}^*\omega_i - q_{02}^*\omega_i = \text{d}\omega_i$.
したがって上の二つの式の差は
$$
\sum q_2^*s_i \cdot \text{d}\omega_i -
\sum q_{12}^*\nabla(s_i) \cdot q_{01}^*\omega_i
$$
である。$\mathcal{O}_{T''}$ 上の乗法の定義により
$q_{12}^*\omega \cdot q_{01}^*\omega' = \omega' \wedge \omega =
- \omega \wedge \omega'$ である。したがって上の式は、
$q_2^*\mathcal{F}$ の部分層
$\mathcal{F}_T \otimes \Omega^2_{T/S, \delta}$ の切断とみなした
$\nabla^2(s)$ である。ゆえに可積分条件が得られる。
\end{proof}




\section{余単体代数}
\label{crystalline-section-cosimplicial}

\noindent
この節は別の場所へ移すべきである。
{\it 余単体環}とは、環の圏における余単体対象をいう。
環 $R$ が与えられたとき、{\it 余単体 $R$-代数}とは
$R$-代数の圏における余単体対象をいう。
余単体環 $A_*$ の {\it 余単体イデアル}とは、各 $n$ に対する
イデアル $I_n \subset A_n$ であって、$\Delta$ のすべての
$f : [n] \to [m]$ に対して $A(f)(I_n) \subset I_m$ を満たすものをいう。

\medskip\noindent
$A_*$ を余単体環とする。$\mathcal{C}$ を、環 $A$ と $A$ 上の加群
$M$ の対 $(A, M)$ からなる圏とする。射
$(A, M) \to (A', M')$ は、環準同型 $A \to A'$ と
$A$-加群準同型 $M \to M'$ からなる。ここで $M'$ は
$A \to A'$ と $M'$ 上の $A'$-加群構造を通じて $A$-加群とみなす。
このもとで、{\it $A_*$ 上の余単体加群 $M_*$}とは、
第一成分が $A_*$ に等しい $\mathcal{C}$ の余単体対象
$(A_*, M_*)$ をいう。{\it $A_*$ 上の余単体加群の準同型
$\varphi_* : M_* \to N_*$}とは、第一成分が $1_{A_*}$ である
$\mathcal{C}$ の余単体対象の射
$(A_*, M_*) \to (A_*, N_*)$ をいう。

\medskip\noindent
$A_*$ 上の余単体加群の準同型 $\varphi_*, \psi_* : M_* \to N_*$
の間の {\it ホモトピー}とは、対応する写像
$(A_*, M_*) \to (A_*, N_*)$ の間のホモトピーであって、
その第一成分が自明なホモトピーであるものをいう
（『単体論』例 \ref{simplicial-example-trivial-homotopy} の双対）。
これを具体的に述べる。このようなホモトピーとは、アーベル余単体群の
準同型としての $\varphi_*$ と $\psi_*$ の間のホモトピー
$$
h : M_* \longrightarrow \Hom(\Delta[1], N_*)
$$
であって、各 $n$ に対する写像
$h_n : M_n \to \prod_{\alpha \in \Delta[1]_n} N_n$ が
$A_n$-線形であるものをいう。次の補題は『単体論』補題
\ref{simplicial-lemma-functorial-homotopy} の余単体加群版である。

\begin{lemma}
\label{crystalline-lemma-homotopy-tensor}
$A_*$ を余単体環とし、$\varphi_*, \psi_* : K_* \to M_*$ を
余単体 $A_*$-加群の準同型とする。
\begin{enumerate}
\item
\label{crystalline-item-tensor}
$\varphi_*$ と $\psi_*$ がホモトピックならば、
$$
\varphi_* \otimes 1, \psi_* \otimes 1 :
K_* \otimes_{A_*} L_* \longrightarrow M_* \otimes_{A_*} L_*
$$
は任意の余単体 $A_*$-加群 $L_*$ に対してホモトピックである。
\item
\label{crystalline-item-wedge}
$\varphi_*$ と $\psi_*$ がホモトピックならば、
$$
\wedge^i(\varphi_*), \wedge^i(\psi_*) :
\wedge^i(K_*) \longrightarrow \wedge^i(M_*)
$$
はホモトピックである。
\item
\label{crystalline-item-base-change}
$\varphi_*$ と $\psi_*$ がホモトピックで、
$A_* \to B_*$ が余単体環の準同型ならば、
$$
\varphi_* \otimes 1, \psi_* \otimes 1 :
K_* \otimes_{A_*} B_* \longrightarrow M_* \otimes_{A_*} B_*
$$
は余単体 $B_*$-加群の準同型としてホモトピックである。
\item
\label{crystalline-item-completion}
$I_* \subset A_*$ が余単体イデアルならば、完備化の間に誘導される写像
$$
\varphi^\wedge_*, \psi^\wedge_* :
K_*^\wedge \longrightarrow M_*^\wedge
$$
はホモトピックである。
\item 必要に応じて、例えば対称冪についての項をここに追加する。
\end{enumerate}
\end{lemma}

\begin{proof}
$h : M_* \longrightarrow \Hom(\Delta[1], N_*)$ を与えられた
ホモトピーとする。次数 $n$ では
$$
h_n = (h_{n, \alpha}) :
K_n \longrightarrow
\prod\nolimits_{\alpha \in \Delta[1]_n} K_n
$$
である。『単体論』節
\ref{simplicial-section-homotopy-cosimplicial} を参照されたい。
$h_{n, \alpha}$ の族がホモトピーをなすための必要十分条件は、
すべての $f : [n] \to [m]$ に対して
$$
h_{m, \alpha} \circ M_*(f) = N_*(f) \circ h_{n, \alpha \circ f}
$$
が成り立つことである。『単体論』式
(\ref{simplicial-equation-property-homotopy-cosimplicial}) を参照されたい。
さらに $\psi_n = h_{n, 0 : [n] \to [1]}$ および
$\varphi_n = h_{n, 1 : [n] \to [1]}$ でなければならない。

\medskip\noindent
補題の各場合に対応する写像を構成できる。
場合 (\ref{crystalline-item-tensor})。次数 $n$ で
$$
(h \otimes 1)_{n, \alpha} = h_{n, \alpha} \otimes 1_{L_n} :
K_n \otimes_{A_n} L_n
\longrightarrow
M_n \otimes_{A_n} L_n.
$$
と定めたホモトピー $h \otimes 1$ を用いればよい。
場合 (\ref{crystalline-item-wedge})。次数 $n$ で
$$
\wedge^i(h)_{n, \alpha} = \wedge^i(h_{n, \alpha}) :
\wedge_{A_n}(K_n)
\longrightarrow
\wedge^i_{A_n}(M_n).
$$
と定めたホモトピー $\wedge^ih$ を用いればよい。
場合 (\ref{crystalline-item-base-change})。次数 $n$ で
$$
(h \otimes 1)_{n, \alpha} = h_{n, \alpha} \otimes 1 :
K_n \otimes_{A_n} B_n
\longrightarrow
M_n \otimes_{A_n} B_n.
$$
と定めたホモトピー $h \otimes 1$ を用いればよい。
場合 (\ref{crystalline-item-completion})。次数 $n$ で
$$
(h^\wedge)_{n, \alpha} = h_{n, \alpha}^\wedge :
K_n^\wedge
\longrightarrow
M_n^\wedge.
$$
と定めたホモトピー $h^\wedge$ を用いればよい。
各 $h_{n, \alpha}$ は $A_n$-線形なので、この構成でよい。
\end{proof}






\section{準連接加群のクリスタル}
\label{crystalline-section-quasi-coherent-crystals}

\noindent
状況 \ref{crystalline-situation-affine} において
$X = \Spec(C)$ および $S = \Spec(A)$ とおく。
$\text{Cris}(X/S)$ 上の準連接加群のクリスタルを分類する。
その前に記法を固定する。

\medskip\noindent
$A$ 上の多項式環 $P = A[x_i]$ と、核が
$J = \Ker(P \to C)$ である $A$-代数の全射 $P \to C$ を選ぶ。
\begin{equation}
\label{crystalline-equation-D}
D = \lim_e D_{P, \gamma}(J) / p^eD_{P, \gamma}(J)
\end{equation}
とおき、これを $p$-進完備化した分割冪包絡とする。
この環には分割冪イデアル $\bar J$ と分割冪構造 $\bar \gamma$ が備わる。
補題 \ref{crystalline-lemma-list-properties} を参照されたい。
$D_e = D/p^eD$ とおき、$\bar J$ の $D_e$ における像を
$\bar J_e$ と書く。略記
\begin{equation}
\label{crystalline-equation-omega-D}
\Omega_D = \lim_e \Omega_{D_e/A, \bar\gamma} =
\lim_e \Omega_{D/A, \bar\gamma}/p^e\Omega_{D/A, \bar\gamma}
\end{equation}
を分割冪微分加群の $p$-進完備化に対して用いる。
補題 \ref{crystalline-lemma-differentials-completion} を参照されたい。
これは
$\Omega_{D_{P, \gamma}(J)/A, \bar\gamma}$
の $p$-進完備化でもある。この加群は $\text{d}x_i$ 上自由である。
補題 \ref{crystalline-lemma-module-differentials-divided-power-envelope} を参照されたい。
したがって $\Omega_D$ の任意の元は和
$\sum f_i\text{d}x_i$ と一意的に書け、各 $e$ に対して
$p^eD$ に属さない $f_i$ は有限個しかない。さらに写像
$\text{d}_{D_e/A, \bar\gamma} : D_e \to \Omega_{D_e/A, \bar\gamma}$
は整合して、$p$-進完備化上の分割冪 $A$-導分
\begin{equation}
\label{crystalline-equation-derivation-D}
\text{d} : D \longrightarrow \Omega_D
\end{equation}
を定める。

\medskip\noindent
また「$\Spec(D)$ の積 $\Spec(D(n))$」も必要となる。
説明については命題 \ref{crystalline-proposition-compute-cohomology} とその証明を参照されたい。
形式的には次のように定義する。$n \geq 0$ に対して
$J(n) = \Ker(P \otimes_A \ldots \otimes_A P \to C)$ とする。
ここでテンソル積は $n + 1$ 個の因子をもつ。
\begin{equation}
\label{crystalline-equation-Dn}
D(n) = \lim_e
D_{P \otimes_A \ldots \otimes_A P, \gamma}(J(n))/
p^eD_{P \otimes_A \ldots \otimes_A P, \gamma}(J(n))
\end{equation}
とおき、これを分割冪包絡の $p$-進完備化とする。
その分割冪イデアルを $\bar J(n)$、分割冪を $\bar \gamma(n)$ と書く。
さらに $D(n)_e = D(n)/p^eD(n)$、および $p$-進完備化した微分加群
\begin{equation}
\label{crystalline-equation-omega-Dn}
\Omega_{D(n)} = \lim_e \Omega_{D(n)_e/A, \bar\gamma} =
\lim_e \Omega_{D(n)/A, \bar\gamma}/p^e\Omega_{D(n)/A, \bar\gamma}
\end{equation}
と導分
\begin{equation}
\label{crystalline-equation-derivation-Dn}
\text{d} : D(n) \longrightarrow \Omega_{D(n)}
\end{equation}
を導入する。もちろん $D = D(0)$ である。
環 $D(0), D(1), D(2), \ldots$ は分割冪環の圏における
余単体対象をなすことに注意する。

\begin{lemma}
\label{crystalline-lemma-structure-Dn}
$D$ と $D(n)$ を (\ref{crystalline-equation-D}) および
(\ref{crystalline-equation-Dn}) のものとする。余積への標準射
$P \to P \otimes_A \ldots \otimes_A P$、
$f \mapsto f \otimes 1 \otimes \ldots \otimes 1$
は $D$ 上の代数の同型
\begin{equation}
\label{crystalline-equation-structure-Dn}
D(n) = \lim_e D\langle \xi_i(j) \rangle/p^eD\langle \xi_i(j) \rangle
\end{equation}
を誘導する。ここで
$$
\xi_i(j) = x_i \otimes 1 \otimes \ldots \otimes 1 -
1 \otimes \ldots \otimes 1 \otimes x_i \otimes 1 \otimes \ldots \otimes 1
$$
であり、$j = 1, \ldots, n$ に対して第二の $x_i$ は
$j + 1$ 番目の位置に置く。$D(n)$ は $P$ の $A$ 上の
$n + 1$ 重テンソル積から構成されることを想起されたい。
\end{lemma}

\begin{proof}
等式
$$
P \otimes_A \ldots \otimes_A P = P[\xi_i(j)]
$$
が成り立ち、$J(n)$ は $J$ と元 $\xi_i(j)$ により生成される。
したがって補題 \ref{crystalline-lemma-divided-power-envelope-add-variables} から従う。
\end{proof}

\begin{lemma}
\label{crystalline-lemma-property-Dn}
$D$ と $D(n)$ を (\ref{crystalline-equation-D}) および
(\ref{crystalline-equation-Dn}) のものとする。このとき、第一の対
$(D, \bar J, \bar\gamma)$ と、第二の対
$(D(n), \bar J(n), \bar\gamma(n))$ はいずれも
$\text{Cris}^\wedge(C/A)$ の対象であり
（注意 \ref{crystalline-remark-completed-affine-site} を参照）、
$$
D(n) = \coprod\nolimits_{j = 0, \ldots, n} D
$$
が $\text{Cris}^\wedge(C/A)$ において成り立つ。
\end{lemma}

\begin{proof}
第一の主張は明らかである。第二の主張について、
$(B \to C, \delta)$ を $\text{Cris}^\wedge(C/A)$ の対象とすると
$$
\Mor_{\text{Cris}^\wedge(C/A)}(D, B) = 
\Hom_A((P, J), (B, \Ker(B \to C)))
$$
である。$D(n)$ についても同様であり、$(P, J)$ を
$(P \otimes_A \ldots \otimes_A P, J(n))$ に置き換えればよい。
$P \otimes_A \ldots \otimes_A P$ は余積なので、余積に関する主張が従う。
\end{proof}

\noindent
以下の補題では、次の条件を満たす対 $(M, \nabla)$ を考える。
\begin{enumerate}
\item
\label{crystalline-item-complete}
$M$ は $p$-進完備な $D$-加群である。
\item
\label{crystalline-item-connection}
$\nabla : M \to M \otimes^\wedge_D \Omega_D$ は接続、すなわち
$\nabla(fm) = m \otimes \text{d}f + f\nabla(m)$,
\item
\label{crystalline-item-integrable}
$\nabla$ は可積分である
（注意 \ref{crystalline-remark-connection} を参照）。
\item
\label{crystalline-item-topologically-quasi-nilpotent}
$\nabla$ は {\it 位相的準冪零}である。すなわち、ある作用素
$\theta_i : M \to M$ によって
$\nabla(m) = \sum \theta_i(m)\text{d}x_i$ と書くとき、
任意の $m \in M$ に対して
$\theta_i^k(m) \not \in pM$ となる対 $(i, k)$ は有限個しかない。
\end{enumerate}
作用素 $\theta_i$ は文献では
$\nabla_{\partial/\partial x_i}$ と書かれることもある。
次の補題で、$\text{Cris}(X/S)$ 上の準連接加群のクリスタルから
このような対の圏への関手を構成する。命題
\ref{crystalline-proposition-crystals-on-affine} で、この関手が圏同値であることを示す。

\begin{lemma}
\label{crystalline-lemma-crystals-on-affine}
上の状況において関手
$$
\begin{matrix}
\text{準連接} \\
\mathcal{O}_{X/S}\text{-加群のクリスタル --- サイト }\text{Cris}(X/S)
\end{matrix}
\longrightarrow
\begin{matrix}
\text{条件を満たす対 }(M, \nabla)\text{：} \\
\text{(\ref{crystalline-item-complete}), (\ref{crystalline-item-connection}),
(\ref{crystalline-item-integrable}), および (\ref{crystalline-item-topologically-quasi-nilpotent})}
\end{matrix}
$$
が存在する。
\end{lemma}

\begin{proof}
$\mathcal{F}$ を $X/S$ 上の準連接加群のクリスタルとする。
$e \gg 0$ に対して $(X, T_e, \bar\gamma)$ が
$\text{Cris}(X/S)$ の対象となるように $T_e = \Spec(D_e)$ とおく。
閉埋め込みである射
$$
(X, T_e, \bar\gamma) \to (X, T_{e + 1}, \bar\gamma) \to \ldots
$$
がある。
$$
M =
\lim_e \Gamma((X, T_e, \bar\gamma), \mathcal{F}) =
\lim_e \Gamma(T_e, \mathcal{F}_{T_e}) = \lim_e M_e
$$
とおく。$\mathcal{F}$ は局所準連接なので
$\mathcal{F}_{T_e} = \widetilde{M_e}$ である。
$\mathcal{F}$ はクリスタルなので
$M_e = M_{e + 1}/p^eM_{e + 1}$ である。したがって
$M_e = M/p^eM$ であり、$M$ は $p$-進完備である。
『代数』補題 \ref{algebra-lemma-limit-complete} を参照されたい。

\medskip\noindent
補題 \ref{crystalline-lemma-automatic-connection} により、$\mathcal{F}$ には標準的な
可積分接続
$\nabla : \mathcal{F} \to \mathcal{F} \otimes \Omega_{X/S}$ が備わる。
この接続を上で構成した対象 $T_e$ 上で評価すると、標準的な可積分接続
$$
\nabla : M \longrightarrow M \otimes^\wedge_D \Omega_D
$$
を得る。これが位相的冪零であることを示すため、その意味を具体的に調べる。

\medskip\noindent
次に環 $D(n)$ に対して同じ手続きを行う。
これにより $p$-進完備な $D(n)$-加群 $M(n)$ を得る。
再び $\mathcal{F}$ のクリスタル性を用いると同型
$$
M \otimes^\wedge_{D, p_0} D(1) \rightarrow M(1)
\leftarrow M \otimes^\wedge_{D, p_1} D(1)
$$
を得る。補題 \ref{crystalline-lemma-automatic-connection} の証明と比較されたい。
左から右への合成を $c$ と書き、$m \in M$ を選ぶ。
$\xi_i = x_i \otimes 1 - 1 \otimes x_i$ と書く。
(\ref{crystalline-equation-structure-Dn}) により、一意的に
$$
c(m \otimes 1) = \sum\nolimits_K \theta_K(m) \otimes \prod \xi_i^{[k_i]}
$$
と書ける。ここで $\theta_K(m) \in M$ であり、和は
$k_i \geq 0$ および $\sum k_i < \infty$ を満たす多重指数
$K = (k_i)$ にわたる。$i$ 番目が $1$ で他が零である
$K$ に対して $\theta_i = \theta_K$ とおく。このとき
$$
\nabla(m) = \sum \theta_i(m) \text{d}x_i.
$$
であることは $\nabla$ の定義と比較すれば分かる。実際、補題
\ref{crystalline-lemma-automatic-connection} における定義式は
$p_1^*m = \nabla(m) - c(p_0^*m)$ であるが、Stacks Project では
対角線のイデアルの平方を法として一貫して
$\text{d}f = p_1(f) - p_0(f)$ を用いているため、符号は整合する。
したがって $\xi_i = x_i \otimes 1 - 1 \otimes x_i$ は、
対角線のイデアルの平方を法として $-\text{d}x_i$ に写る。

\medskip\noindent
添字 $i,j$ に対応する余積への標準射を
$q_i : D \to D(2)$ および $q_{ij} : D(1) \to D(2)$ と書く。
補題 \ref{crystalline-lemma-automatic-connection} の証明の最終段落と同様に
$$
q_{02}^*c = q_{12}^*c \circ q_{01}^*c.
$$
である。これは
$$
\sum\nolimits_{K''} \theta_{K''}(m) \otimes \prod {\zeta''_i}^{[k''_i]}
=
\sum\nolimits_{K', K} \theta_{K'}(\theta_K(m))
\otimes \prod {\zeta'_i}^{[k'_i]} \prod \zeta_i^{[k_i]}
$$
が $M \otimes^\wedge_{D, q_2} D(2)$ において成り立つことを意味する。
ここで
\begin{align*}
\zeta_i & = x_i \otimes 1 \otimes 1 - 1 \otimes x_i \otimes 1,\\
\zeta'_i & = 1 \otimes x_i \otimes 1 - 1 \otimes 1 \otimes x_i,\\
\zeta''_i & = x_i \otimes 1 \otimes 1 - 1 \otimes 1 \otimes x_i.
\end{align*}
特に $\zeta''_i = \zeta_i + \zeta'_i$ であり、$D(2)$ は
$q_2(D)$ 上の $\zeta_i, \zeta'_i$ に関する分割冪多項式環の
$p$-進完備化である。補題 \ref{crystalline-lemma-structure-Dn} を参照されたい。
上の式で係数を比較すると直ちに
$\theta_i \circ \theta_j = \theta_j \circ \theta_i$
（これは $\nabla$ の可積分性の別証明を与える）および
$$
\theta_K(m) = (\prod \theta_i^{k_i})(m).
$$
が従う。特に、上で $c(m \otimes 1)$ を表す和は $p$-進収束しなければ
ならないので、各 $i$ と各 $m \in M$ に対して
$\theta_i^k(m)$ が $p$ を法として非零となるものは有限個しかない。
\end{proof}

\begin{proposition}
\label{crystalline-proposition-crystals-on-affine}
関手
$$
\begin{matrix}
\text{準連接} \\
\mathcal{O}_{X/S}\text{-加群のクリスタル --- サイト }\text{Cris}(X/S)
\end{matrix}
\longrightarrow
\begin{matrix}
\text{条件を満たす対 }(M, \nabla)\text{：} \\
\text{(\ref{crystalline-item-complete}), (\ref{crystalline-item-connection}),
(\ref{crystalline-item-integrable}), および (\ref{crystalline-item-topologically-quasi-nilpotent})}
\end{matrix}
$$
は補題 \ref{crystalline-lemma-crystals-on-affine} で構成したものであり、圏同値である。
\end{proposition}

\begin{proof}
$(M, \nabla)$ が与えられたとし、準連接加群のクリスタル
$\mathcal{F}$ を構成する。
$\nabla(m) = \sum \theta_i(m)\text{d}x_i$ と書く。
このとき $\theta_i \circ \theta_j = \theta_j \circ \theta_i$ であり、
$k_i \geq 0$ および $\sum k_i < \infty$ を満たす任意の多重指数
$K = (k_i)$ に対して
$\theta_K(m) = (\prod \theta_i^{k_i})(m)$ とおける。

\medskip\noindent
$(U, T, \delta)$ を、$T$ がアフィンである
$\text{Cris}(X/S)$ の任意の対象とする。
$T = \Spec(B)$ とし、$U \to T$ のイデアルを $J_B \subset B$ とする。
補題 \ref{crystalline-lemma-set-generators} により、ある整数 $e$ と射
$$
f : (U, T, \delta) \longrightarrow (X, T_e, \bar\gamma)
$$
が存在する。ここで $T_e = \Spec(D_e)$ は補題
\ref{crystalline-lemma-crystals-on-affine} の証明と同じである。
このような $e$ と $f$ を選び、対応する分割冪 $A$-代数準同型も
$f : D \to B$ と書く。$\mathcal{F}_T$ を、$B$-加群
$$
M \otimes_{D, f} B.
$$
に付随する準連接 $\mathcal{O}_T$-加群の層とおく。
ただし、これが $f$ の選択に依存しないことを示す必要がある。
$g : D \to B$ をもう一つのそのような射とする。$f$ と $g$ は
$\text{Cris}(X/S)$ における射なので、
$f - g : D \to B$ の像は分割冪イデアル $J_B$ に含まれる。
$\xi_i = f(x_i) - g(x_i) \in J_B$ と書く。
補題 \ref{crystalline-lemma-crystals-on-affine} の証明にならって、同型
$$
c_{f, g} : M \otimes_{D, f} B \longrightarrow M \otimes_{D, g} B
$$
を公式
$$
m \otimes 1 \longmapsto 
\sum\nolimits_K \theta_K(m) \otimes \prod \xi_i^{[k_i]}
$$
で定義する。これは上の議論と、$\nabla$ が位相的準冪零であること
（したがって和は有限である）により意味をもつ。計算により
$$
c_{g, h} \circ c_{f, g} = c_{f, h}
$$
が分かる。ここで
$h : (U, T, \delta) \longrightarrow (X, T_e, \bar\gamma)$ は
第三の射である。また $c_{f, f} = 1$ でもある。
したがってこれらの写像はすべて同型であり、加群 $\mathcal{F}_T$ は
$f$ の選択に依存しない。

\medskip\noindent
$a : (U', T', \delta') \to (U, T, \delta)$ を
$\text{Cris}(X/S)$ のアフィン対象の射とする。
$f' = f \circ a$ を選べば、標準的な同型
$a^*\mathcal{F}_T \to \mathcal{F}_{T'}$ が存在することは明らかである。
この写像が $f$ の選択に依存しないことの検証は省略する。
これらを制限写像として用いると、$\text{Cris}(X/S)$ のアフィン対象からなる
充満部分圏上に準連接加群のクリスタルが得られることは明らかである。
これが $\text{Cris}(X/S)$ 全体のクリスタルへ拡張することの証明は省略する。
また、この手続きが関手をなし、補題
\ref{crystalline-lemma-crystals-on-affine} で構成した関手の準逆となることの証明も省略する。
\end{proof}

\begin{lemma}
\label{crystalline-lemma-crystals-on-affine-smooth}
状況 \ref{crystalline-situation-affine} において、$A \to P' \to C$ を環準同型とし、
$A \to P'$ は滑らか、$P' \to C$ は核 $J'$ をもつ全射とする。
$D'$ を $D_{P', \gamma}(J')$ の $p$-進完備化とする。
このとき上のようなデータ $A \to P \to C$ の選択と、
分割冪 $A$-代数の準同型
$$
a : D \longrightarrow D',\quad b : D' \longrightarrow D
$$
が存在する。これらは $D \to C$ および $D' \to C$ と両立し、
$a \circ b = \text{id}_{D'}$ を満たす。
これらの写像は、$D$ 上で
(\ref{crystalline-item-complete}), (\ref{crystalline-item-connection}),
(\ref{crystalline-item-integrable}), (\ref{crystalline-item-topologically-quasi-nilpotent})
を満たす対 $(M, \nabla)$ の圏と、$D'$ 上で
(\ref{crystalline-item-complete}), (\ref{crystalline-item-connection}),
(\ref{crystalline-item-integrable}), (\ref{crystalline-item-topologically-quasi-nilpotent})
\footnote{$D'$ 上の $(M', \nabla')$ に対してこの条件を定式化するのは
注意を要する。証明を参照されたい。}を満たす対 $(M', \nabla')$ の圏との
圏同値を誘導する。特に、命題
\ref{crystalline-proposition-crystals-on-affine} の圏同値は、$D'$ 上の対への
対応する関手についても成り立つ。
\end{lemma}

\begin{proof}
$A$ 上の多項式代数 $P = A[y_1, \ldots, y_m]$ と全射
$P \to P'$ を選ぶ。$P \to C$ を合成 $P \to P' \to C$ で定義する。
分割冪包絡と完備化の関手性により、全射 $a : D \to D'$ を得る。
$D_e$ が $A$ 上の $C$ の分割冪肥厚化となるよう、十分大きな $e$ を選ぶ。
このとき $D_e \to C$ は局所冪零な核をもつ全射である。
『分割冪代数』補題 \ref{dpa-lemma-nil} を参照されたい。
$D'_e = D'/p^eD'$ とおくと、$D_e \to D'_e$ の核は局所冪零である。
したがって『代数』補題 \ref{algebra-lemma-smooth-strong-lift} により、
写像 $P' \to D'_e$ の持ち上げ $\beta_e : P' \to D_e$ が存在する。
任意の $i \geq 0$ に対して
$p^{e + i}D \to p^{e + i}D'$ は全射なので、
$D_{e + i + 1} \to D_{e + i} \times_{D'_{e + i}} D'_{e + i + 1}$
は平方零な核をもつ全射である。通常の持ち上げ性
（『代数』命題 \ref{algebra-proposition-smooth-formally-smooth}）を
図式
$$
\xymatrix{
P' \ar[r] & D_{e + i} \times_{D'_{e + i}} D'_{e + i + 1} \\
A \ar[u] \ar[r] & D_{e + i + 1} \ar[u]
}
$$
に順次適用すると、$a$ との合成が与えられた写像 $P' \to D'$ となる
$A$-代数準同型 $\beta : P' \to D$ が存在する。
分割冪包絡の普遍性により写像 $D_{P', \gamma}(J') \to D$ を得る。
$D$ は $p$-進完備なので、$a \circ b = \text{id}_{D'}$ を満たす
$b : D' \to D$ を得る。

\medskip\noindent
基底変換関手
$$
F : (M, \nabla) \longmapsto
(M \otimes^\wedge_{D, a} D', \nabla')
\quad\text{および}\quad
G : (M', \nabla') \longmapsto
(M' \otimes^\wedge_{D', b} D, \nabla)
$$
を、(\ref{crystalline-item-complete}), (\ref{crystalline-item-connection}),
(\ref{crystalline-item-integrable}) を満たす接続付き加群上で考える。
注意 \ref{crystalline-remark-base-change-connection} を参照されたい。
$a \circ b = \text{id}_{D'}$ なので、$F \circ G$ は恒等関手である。
$G(M', \nabla')$ が性質 (\ref{crystalline-item-topologically-quasi-nilpotent}) を
もつとき、$(M', \nabla')$ もこの性質をもつということにする。
ここで形式的な議論により、証明を完了するには、
$(M, \nabla)$ が四条件
(\ref{crystalline-item-complete}), (\ref{crystalline-item-connection}),
(\ref{crystalline-item-integrable}), (\ref{crystalline-item-topologically-quasi-nilpotent})
をすべて満たす場合に $G(F(M, \nabla))$ が $(M, \nabla)$ と
同型であることを示せば十分である。このため、命題
\ref{crystalline-proposition-crystals-on-affine} の証明における関手的同型
$$
c_{\text{id}_D, b \circ a} :
M \otimes_{D, \text{id}_D} D
\longrightarrow
M \otimes_{D, b \circ a} D
$$
を用いる。これには $\nabla$ の位相的準冪零性が必要だが、これは仮定している。
残るのは、この写像が水平、すなわち接続と両立することの証明であるが、省略する。

\medskip\noindent
これで最後の主張も従う。
\end{proof}

\begin{remark}
\label{crystalline-remark-equivalence-more-general}
命題 \ref{crystalline-proposition-crystals-on-affine} の圏同値は、$P/A$ が
『代数』補題 \ref{algebra-lemma-smooth-strong-lift} の強持ち上げ性を
満たす全射 $P \to C$ から始めても成り立つ。
これは補題 \ref{crystalline-lemma-crystals-on-affine-smooth} の証明と同様に示せる。
（必要になれば詳細をここに追加する。）
おそらくこの結果には直接証明もあるが、多項式環を用いる利点は、
環 $D(n)$ が分割冪多項式環の $p$-進完備化となり、代数的議論が
簡単になることである。
\end{remark}



\section{コホモロジーに関する一般的注意}
\label{crystalline-section-cohomology-lqc}

\noindent
この節では、アフィンスキームの結晶サイト上の加群のコホモロジーを
代数的問題へ移すための準備を行う。

\begin{lemma}
\label{crystalline-lemma-vanishing-lqc}
状況 \ref{crystalline-situation-global} において、$\mathcal{F}$ を
$\text{Cris}(X/S)$ 上の局所準連接 $\mathcal{O}_{X/S}$-加群とする。
このとき
$$
H^p((U, T, \delta), \mathcal{F}) = 0
$$
が、すべての $p > 0$ と、$T$ または $U$ がアフィンであるすべての
$(U, T, \delta)$ に対して成り立つ。
\end{lemma}

\begin{proof}
$U \to T$ は肥厚化なので、$U$ がアフィンであることと $T$ が
アフィンであることは同値である。『極限』補題
\ref{limits-lemma-affine} を参照されたい。そこで『サイト上のコホモロジー』
補題 \ref{sites-cohomology-lemma-cech-vanish-collection} を、
アフィン対象 $(U, T, \delta)$ の族 $\mathcal{B}$ と、
アフィン開被覆
$\mathcal{U} = \{(U_i, T_i, \delta_i) \to (U, T, \delta)\}$ の族
$\text{Cov}$ に適用する。このような被覆の {\v C}ech 複体
${\check C}^*(\mathcal{U}, \mathcal{F})$ は、アフィンスキーム $T$ の
アフィン開被覆 $\{T_i \to T\}$ に関する準連接
$\mathcal{O}_T$-加群 $\mathcal{F}_T$ の {\v C}ech 複体にほかならない
（ここで $\mathcal{F}$ が局所準連接であるという仮定を用いている）。
したがって『スキームのコホモロジー』補題
\ref{coherent-lemma-cech-cohomology-quasi-coherent} および
\ref{coherent-lemma-quasi-coherent-affine-cohomology-zero} により、
{\v C}ech コホモロジーは零である。ゆえに『サイト上のコホモロジー』補題
\ref{sites-cohomology-lemma-cech-vanish-collection} の仮定が満たされ、
主張を得る。
\end{proof}

\begin{lemma}
\label{crystalline-lemma-compare}
状況 \ref{crystalline-situation-global} において、さらに $X$ と $S$ は
アフィンスキームであると仮定する。分割冪肥厚化 $(X, T, \delta)$ からなる
充満部分圏 $\mathcal{C} \subset \text{Cris}(X/S)$ にカオス位相を入れる
（『サイト』例 \ref{sites-example-indiscrete} を参照）。
任意の局所準連接 $\mathcal{O}_{X/S}$-加群 $\mathcal{F}$ に対して
$$
R\Gamma(\mathcal{C}, \mathcal{F}|_\mathcal{C}) =
R\Gamma(\text{Cris}(X/S), \mathcal{F})
$$
である。
\end{lemma}

\begin{proof}
$(U, T, \delta)$ であって $U$ と $T$ がともにアフィンなものからなる
$\text{Cris}(X/S)$ の充満部分圏を $\text{AffineCris}(X/S)$ と書く。
射の族
$\{(U_i, T_i, \delta_i) \to (U, T, \delta)\}_{i \in I}$ が
$\text{AffineCris}(X/S)$ における被覆であることを、それが
$\text{Cris}(X/S)$ の被覆であることと定め、これをサイトとする。
この定義のもとで包含関手
$$
\text{AffineCris}(X/S) \longrightarrow \text{Cris}(X/S)
$$
は『サイト』定義
\ref{sites-definition-special-cocontinuous-functor} の意味で
特殊余連続関手である。その証明は『位相』補題
\ref{topologies-lemma-affine-big-site-Zariski} の証明とまったく同じである。
したがって表示した包含関手による制限を通じて、
$\text{Cris}(X/S)$ 上の層のトポスと
$\text{AffineCris}(X/S)$ 上の層のトポスは同じである。
よって包含 $\mathcal{C} \subset \text{AffineCris}(X/S)$ について
対応する主張を示せばよい。

\medskip\noindent
$\mathcal{C}$ と $\text{AffineCris}(X/S)$ が積とファイバー積をもつことを、
以下では断りなく用いる（詳細は省略する。補題
\ref{crystalline-lemma-divided-power-thickening-fibre-products} を参照）。
包含関手 $u : \mathcal{C} \to \text{AffineCris}(X/S)$ は充満忠実かつ
連続で、積およびファイバー積と可換である。これは環付きサイトの射
$$
f :
(\text{AffineCris}(X/S), \mathcal{O}_{X/S})
\longrightarrow
(\Sh(\mathcal{C}), \mathcal{O}_{X/S}|_\mathcal{C})
$$
を定めると主張する。これには『サイト』補題
\ref{sites-lemma-directed-morphism} を用いる。
$\mathcal{C}$ はファイバー積をもち、$u$ はそれと可換なので、圏
$\mathcal{I}^u_{(U, T, \delta)}$ は有向圏の非交和である
（『サイト』補題 \ref{sites-lemma-almost-directed} と『圏論』補題
\ref{categories-lemma-split-into-directed} による）。
したがって $\mathcal{I}^u_{(U, T, \delta)}$ が連結であることを
示せば十分である。非空性は補題 \ref{crystalline-lemma-set-generators} から従う。
実際、$U$ と $T$ はアフィンなので、同補題により
$\mathcal{C}$ の対象 $(X, T', \delta')$ と分割冪肥厚化の射
$(U, T, \delta) \to (X, T', \delta')$ が少なくとも一つ存在する。
連結性は、$\mathcal{C}$ が積をもち $u$ がそれと可換であることから従う
（『サイト』補題 \ref{sites-lemma-directed} の証明と比較されたい）。

\medskip\noindent
$f_*\mathcal{F} = \mathcal{F}|_\mathcal{C}$ であることに注意する。
したがって $p > 0$ に対して $R^pf_*\mathcal{F} = 0$ ならば補題が従う。
『サイト上のコホモロジー』補題
\ref{sites-cohomology-lemma-apply-Leray} を参照されたい。
同書補題 \ref{sites-cohomology-lemma-higher-direct-images} により、\newline
$H^p(\text{AffineCris}(X/S)/(X, T, \delta), \mathcal{F}) = 0$
をすべての $(X, T, \delta)$ に対して示せば十分である。
これは補題 \ref{crystalline-lemma-vanishing-lqc} から従う。実際、サイト\newline
$\text{AffineCris}(X/S)/(X, T, \delta)$ のトポスは、同補題で用いた
サイト\newline
$\text{Cris}(X/S)/(X, T, \delta)$ のトポスと同値である。
\end{proof}

\begin{lemma}
\label{crystalline-lemma-complete}
状況 \ref{crystalline-situation-affine} において
$\mathcal{C} = (\text{Cris}(C/A))^{opp}$ および
$\mathcal{C}^\wedge = (\text{Cris}^\wedge(C/A))^{opp}$
とおき、カオス位相を入れる。記法については注意
\ref{crystalline-remark-completed-affine-site} を参照されたい。
トポスの射
$$
g : \Sh(\mathcal{C}) \longrightarrow \Sh(\mathcal{C}^\wedge)
$$
であって、$\mathcal{F}$ が $\mathcal{C}$ 上のアーベル群の層ならば
$$
R^pg_*\mathcal{F}(B \to C, \delta) =
\left\{
\begin{matrix}
\lim_e \mathcal{F}(B_e \to C, \delta) & \text{条件：}p = 0 \\
R^1\lim_e \mathcal{F}(B_e \to C, \delta) & \text{条件：}p = 1 \\
0 & \text{その他の場合}
\end{matrix}
\right.
$$
となるものが存在する。ここで $e \gg 0$ に対して $B_e = B/p^eB$ である。
\end{lemma}

\begin{proof}
圏の間の任意の関手は、同じ向きのカオストポス間の射を定める。
例えば、その関手をサイト間の余連続関手とみなせるからである。
『サイト』節 \ref{sites-section-cocontinuous-morphism-topoi} を参照されたい。
$g_*\mathcal{F}$ の記述の証明は省略する。主張では、十分大きな $e$ に
対してのみ $(B_e \to C, \delta)$ が $\text{Cris}(C/A)$ の対象である。
$\mathcal{I}$ を $\mathcal{C}$ 上の単射的アーベル層とする。
このとき遷移写像
$$
\mathcal{I}(B_e \to C, \delta) \leftarrow
\mathcal{I}(B_{e + 1} \to C, \delta)
$$
は全射である。実際、射
$$
(B_e \to C, \delta)
\longrightarrow
(B_{e + 1} \to C, \delta)
$$
は圏 $\mathcal{C}$ における単射である。したがって単射的アーベル層に
対して補題の表示式の両辺は一致する。$\mathcal{F}$ の単射分解を取れば
結果は容易に得られる（層は前層なので、完全性は対象上の切断群のレベルで
測られる）。
\end{proof}

\begin{lemma}
\label{crystalline-lemma-category-with-covering}
$\mathcal{C}$ をカオス位相を備えた圏とする。
$X$ を $\mathcal{C}$ の対象で、$\mathcal{C}$ のすべての対象から
$X$ への射が存在するものとする。$\mathcal{C}$ は二対象の積をもつと仮定する。
このとき $\mathcal{C}$ 上の任意のアーベル層 $\mathcal{F}$ に対し、
全コホモロジー $R\Gamma(\mathcal{C}, \mathcal{F})$ は複体
$$
\mathcal{F}(X) \to \mathcal{F}(X \times X) \to
\mathcal{F}(X \times X \times X) \to \ldots
$$
で表される。これは余単体アーベル群
$[n] \mapsto \mathcal{F}(X^n)$ に付随する複体である。
\end{lemma}

\begin{proof}
任意の前層が $\mathcal{C}$ 上の層なので、すべての $q > 0$ に対して
$H^q(X^p, \mathcal{F}) = 0$ である。$X$ に関する仮定は
$h_X \to *$ が全射であるということである。
$H^q(X, \mathcal{F}) = H^q(h_X, \mathcal{F})$ および
$H^q(\mathcal{C}, \mathcal{F}) = H^q(*, \mathcal{F})$ を用いると、
主張は『サイト上のコホモロジー』補題
\ref{sites-cohomology-lemma-cech-to-cohomology-sheaf-sets} の特別な場合である。
\end{proof}









\section{余単体的準備}
\label{crystalline-section-cohomology}

\noindent
この節では結晶コホモロジーと de Rham コホモロジーを比較する。
\cite{Bhatt} に従う。

\begin{example}
\label{crystalline-example-cosimplicial-module}
$A_*$ を任意の余単体環とし、規則
$$
M_n = \bigoplus\nolimits_{i = 0, ..., n} A_n e_i
$$
で定義される余単体加群 $M_*$ を考える。
写像 $f : [n] \to [m]$ に対し、$M_*(f) : M_n \to M_m$ を
$e_i$ を $e_{f(i)}$ に写す一意な $A_*(f)$-線形写像と定める。
$M_*$ 上の恒等写像は $0$ とホモトピックであると主張する。
実際、ホモトピーは余単体加群の写像
$$
h : M_* \longrightarrow \Hom(\Delta[1], M_*)
$$
で与えられる。節 \ref{crystalline-section-cosimplicial} を参照されたい。
$j \in \{0, \ldots, n + 1\}$ に対し、
$\alpha^n_j(i) = 0 \Leftrightarrow i < j$ で定まる写像
$\alpha^n_j : [n] \to [1]$ を取る。このとき
$\Delta[1]_n = \{\alpha^n_0, \ldots, \alpha^n_{n + 1}\}$ であり、対応して
$\Hom(\Delta[1], M_*)_n = \prod_{j = 0, \ldots, n + 1} M_n$ である。
『単体論』節 \ref{simplicial-section-homotopy} および
\ref{simplicial-section-homotopy-cosimplicial} を参照されたい。
この積表示を使う代わりに、$\Hom(\Delta[1], M_*)_n$ の元を
関数 $\Delta[1]_n \to M_n$ とみなす。この記法で、次数 $n$ の $h$ を
$$
h_n(e_i)(\alpha^n_j) =
\left\{
\begin{matrix}
e_{i} & \text{ならば} & i < j \\
0 & \text{それ以外} 
\end{matrix}
\right.
$$
で定義する。まず $h$ が余単体加群の射であることを確かめる。
すなわち $f : [n] \to [m]$ に対して
\begin{equation}
\label{crystalline-equation-cosimplicial-morphism}
h_m \circ M_*(f) = \Hom(\Delta[1], M_*)(f) \circ h_n
\end{equation}
を示す。(\ref{crystalline-equation-cosimplicial-morphism}) の左辺を $e_i$ で評価し、
さらに $\alpha^m_j$ で評価すると
$$
h_m(e_{f(i)})(\alpha^m_j) =
\left\{
\begin{matrix}
e_{f(i)} & \text{ならば} & f(i) < j \\
0 & \text{それ以外}
\end{matrix}
\right.
$$
である。$\alpha^m_j \circ f = \alpha^n_{j'}$ であることに注意する。
ここで $0 \leq j' \leq n + 1$ は、
$f(i) < j$ と $i < j'$ が同値となる一意な添字である。
したがって (\ref{crystalline-equation-cosimplicial-morphism}) の右辺を $e_i$ で評価し、
さらに $\alpha^m_j$ で評価すると
$$
M_*(f)(h_n(e_i)(\alpha^m_j \circ f) =
M_*(f)(h_n(e_i)(\alpha^n_{j'})) =
\left\{
\begin{matrix}
e_{f(i)} & \text{ならば} & i < j' \\
0 & \text{それ以外} 
\end{matrix}
\right.
$$
となる。$j'$ の記述から二つの答えは等しい。ゆえに $h$ は
余単体加群の写像である。$0 : \Delta[0] \to \Delta[1]$ と
$1 : \Delta[0] \to \Delta[1]$ を明らかな写像とし、対応する評価写像を
$ev_0, ev_1 : \Hom(\Delta[1], M_*) \to M_*$ と書く。
合成
$$
ev_0 \circ h, ev_1 \circ h : M_* \longrightarrow M_*
$$
がそれぞれ $1$ と $0$ であることは容易に確かめられる。
したがって $h$ は $1$ と $0$ の間の所望のホモトピーである。
\end{example}

\begin{lemma}
\label{crystalline-lemma-vanishing-omega-1}
(\ref{crystalline-equation-omega-Dn}) の記法のもとで、複体
$$
\Omega_{D(0)} \to \Omega_{D(1)} \to \Omega_{D(2)} \to \ldots
$$
は余単体 $D(*)$-加群として零とホモトピックである。
\end{lemma}

\begin{proof}
『単体論』補題 \ref{simplicial-lemma-functorial-homotopy} の原理、
より具体的には補題 \ref{crystalline-lemma-homotopy-tensor} を用いる。
これは（余）単体対象間のホモトピックな写像が、任意の関手によって
ホモトピックな写像へ移されることをいう。補題の複体は、余単体加群
$$
M_* = \left(
\Omega_{P/A} \to
\Omega_{P \otimes_A P/A} \to
\Omega_{P \otimes_A P \otimes_A P/A} \to \ldots
\right)
$$
を余単体環準同型 $P\otimes_A \ldots \otimes_A P \to D(n)$ に沿って
基底変換し、$p$-進完備化したものに等しい。これは補題
\ref{crystalline-lemma-module-differentials-divided-power-envelope} から従う。
(\ref{crystalline-equation-omega-D}) の後の注意も参照されたい。
したがって余単体加群 $M_*$ が零とホモトピックであることを示せば十分である
（基底変換と $p$-進完備化を用いる）。
$\mathbf{Z} \to A$ に沿う基底変換を使えるので、
$A = \mathbf{Z}$ および $P = \mathbf{Z}[\{x_i\}_{i \in I}]$ と
仮定してよい。この場合 $P^{\otimes n + 1}$ は元
$$
x_i(e) = 1 \otimes \ldots \otimes x_i \otimes \ldots \otimes 1
$$
上の多項式代数であり、$x_i$ は $e$ 番目の位置にある。
複体の各加群は生成元 $\text{d}x_i(e)$ 上自由である。
$f : [n] \to [m]$ が写像ならば
$$
M_*(f)(\text{d}x_i(e)) = \text{d}x_i(f(e))
$$
である。したがって $M_*$ は、例 \ref{crystalline-example-cosimplicial-module} で
調べた加群のコピーを $I$ にわたって直和したものであり、主張が従う。
\end{proof}

\begin{lemma}
\label{crystalline-lemma-vanishing}
(\ref{crystalline-equation-Dn}) と (\ref{crystalline-equation-omega-Dn}) の記法のもとで、
$D(*)$ 上の任意の余単体加群 $M_*$ と $i > 0$ に対し、余単体加群
$$
M_0 \otimes^\wedge_{D(0)} \Omega^i_{D(0)} \to
M_1 \otimes^\wedge_{D(1)} \Omega^i_{D(1)} \to
M_2 \otimes^\wedge_{D(2)} \Omega^i_{D(2)} \to \ldots
$$
は零とホモトピックである。ここで $\Omega^i_{D(n)}$ は
$\Omega_{D(n)}$ の第 $i$ 外冪の $p$-進完備化である。
\end{lemma}

\begin{proof}
補題 \ref{crystalline-lemma-vanishing-omega-1} により、$\Omega_{D(*)}$ の自己準同型
$0$ と $1$ はホモトピックである。関手 $\wedge^i$ を適用すると、
余単体加群 $\wedge^i\Omega_{D(*)}$ についても同じことが成り立つ。
補題 \ref{crystalline-lemma-homotopy-tensor} を参照されたい。
同補題をもう一度適用すると、$p$-進完備化
$\Omega^i_{D(*)}$ は零とホモトピー同値である。
$M_*$ とテンソル積を取ると
$M_* \otimes_{D(*)} \Omega^i_{D(*)}$ は零とホモトピックである。
再び補題 \ref{crystalline-lemma-homotopy-tensor} を参照されたい。
最後に $p$-進完備化関手を適用して証明が完了する。
\end{proof}



\section{分割冪 Poincar\'e 補題}
\label{crystalline-section-poincare}

\noindent
可能な限り最も単純な形だけを扱う。

\begin{lemma}
\label{crystalline-lemma-poincare}
$A$ を環とし、$P = A\langle x_i \rangle$ を $A$ 上の
分割冪多項式環とする。任意の $A$-加群 $M$ に対し、複体
$$
0 \to M \to
M \otimes_A P \to
M \otimes_A \Omega^1_{P/A, \delta} \to
M \otimes_A \Omega^2_{P/A, \delta} \to \ldots
$$
は完全である。$D$ を $P$ の $p$-進完備化とし、
$\Omega^i_D$ を $\Omega_{D/A, \delta}$ の第 $i$ 外冪の
$p$-進完備化とする。任意の $p$-進完備な $A$-加群 $M$ に対し、複体
$$
0 \to M \to
M \otimes^\wedge_A D \to
M \otimes^\wedge_A \Omega^1_D \to
M \otimes^\wedge_A \Omega^2_D \to \ldots
$$
は完全である。
\end{lemma}

\begin{proof}
複体
$$
E :
(0 \to A \to P \to \Omega^1_{P/A, \delta} \to
\Omega^2_{P/A, \delta} \to \ldots)
$$
が $A$-加群の複体として零とホモトピー同値であることを示せば十分である。
各多重指数 $K = (k_i)$ に対し、次数 $j$ の項が
$$
\bigoplus\nolimits_{I = \{i_1, \ldots, i_j\} \subset \text{Supp}(K)}
A
\prod\nolimits_{i \not \in I} x_i^{[k_i]}
\prod\nolimits_{i \in I} x_i^{[k_i - 1]}
\text{d}x_{i_1} \wedge \ldots \wedge \text{d}x_{i_j}
$$
からなる部分複体 $E(K)$ を考える。$E = \bigoplus E(K)$ なので、
各複体 $E(K)$ が零とホモトピックであることを示せば十分である。
$K = 0$ ならば $E(K) : (A \to A)$ は零とホモトピックである。
$K$ が空でない有限な台 $S$ をもつならば、複体 $E(K)$ は複体
$$
0 \to A \to
\bigoplus\nolimits_{s \in S} A \to
\wedge^2(\bigoplus\nolimits_{s \in S} A) \to
\ldots \to \wedge^{\# S}(\bigoplus\nolimits_{s \in S} A) \to 0
$$
と同型である。これは例えば『代数詳論』補題
\ref{more-algebra-lemma-homotopy-koszul-abstract} により
零とホモトピックである。
\end{proof}

\noindent
次の補題への別の、より直接的な方法を例 \ref{crystalline-example-integrate} で説明する。

\begin{lemma}
\label{crystalline-lemma-relative-poincare}
$A$ を環とし、$(B, I, \delta)$ を分割冪環で $B$ が
$A$-代数であるものとする。$P = B\langle x_i \rangle$ を $B$ 上の
分割冪多項式環とし、通常どおり分割冪イデアルを
$J = IP + B\langle x_i \rangle_{+}$ とする。
$M$ を可積分接続
$\nabla : M \to M \otimes_B \Omega^1_{B/A, \delta}$ を備えた
$B$-加群とする。このとき de Rham 複体の写像
$$
M \otimes_B \Omega^*_{B/A, \delta}
\longrightarrow
M \otimes_P \Omega^*_{P/A, \delta}
$$
は擬同型である。$D$ と $D'$ をそれぞれ $B$ と $P$ の $p$-進完備化、
$\Omega^i_D$ と $\Omega^i_{D'}$ をそれぞれ
$\Omega^i_{B/A, \delta}$ と $\Omega^i_{P/A, \delta}$ の
$p$-進完備化とする。$M$ を整接続
$\nabla : M \to M \otimes^\wedge_D \Omega^1_D$ を備えた
$p$-進完備な $D$-加群とする。このとき de Rham 複体の写像
$$
M \otimes^\wedge_D \Omega^*_D
\longrightarrow
M \otimes^\wedge_D \Omega^*_{D'}
$$
は擬同型である。
\end{lemma}

\begin{proof}
部分複体
$F^i(\Omega^*_{B/A, \delta}) = \sigma_{\geq i}\Omega^*_{B/A, \delta}$.
が与える $\Omega^*_{B/A, \delta}$ 上の減少フィルトレーション $F^*$ を考える。
『ホモロジー』節 \ref{homology-section-truncations} を参照されたい。
これは
$$
F^i(\Omega^*_{P/A, \delta}) =
F^i(\Omega^*_{B/A, \delta}) \wedge \Omega^*_{P/A, \delta}.
$$
とおくことにより $\Omega^*_{P/A, \delta}$ 上の減少フィルトレーション
$F^*$ を誘導する。分裂短完全列
$$
0 \to \Omega^1_{B/A, \delta} \otimes_B P \to
\Omega^1_{P/A, \delta} \to
\Omega^1_{P/B, \delta} \to 0
$$
があり、最後の加群は $\text{d}x_i$ 上自由である。
したがって $F^i(\Omega^*_{P/A, \delta}) \to
\Omega^*_{P/A, \delta}$ は項ごとに分裂する単射であり、
$$
\text{gr}^i_F(\Omega^*_{P/A, \delta}) =
\Omega^i_{B/A, \delta} \otimes_B \Omega^*_{P/B, \delta}
$$
が複体として成り立つ。そこで
$$
F^i(M \otimes_B \Omega^*_{P/A, \delta}) =
M \otimes_B F^i(\Omega^*_{P/A, \delta})
$$
とおくことで $M \otimes_B \Omega^*_{B/A, \delta}$ 上の
フィルトレーション $F^*$ を定義でき、
$$
\text{gr}^i_F(M \otimes_B \Omega^*_{P/A, \delta}) =
M \otimes_B \Omega^i_{B/A, \delta} \otimes_B \Omega^*_{P/B, \delta}
$$
が複体として成り立つ。補題 \ref{crystalline-lemma-poincare} により、
これらの各複体は次数 $0$ に置いた
$M \otimes_B \Omega^i_{B/A, \delta}$ と擬同型である。
したがって補題の第一の表示写像はフィルター付き複体の射であり、
次数付き部分上で擬同型を誘導する。ゆえにこれは擬同型である。
例えばフィルター付き複体に付随するスペクトル系列を用いればよい。
『ホモロジー』節 \ref{homology-section-filtered-complex} を参照されたい。

\medskip\noindent
第二の擬同型の証明もまったく同じである。
\end{proof}



\section{アフィンの場合のコホモロジー}
\label{crystalline-section-cohomology-affine}

\noindent
節 \ref{crystalline-section-quasi-coherent-crystals} で調べた状況に戻ろう。
$(A, I, \gamma)$ と $A/I \to C$ から出発し、
$X = \Spec(C)$ および $S = \Spec(A)$ とおく。次に、
$A$ 上の多項式環 $P$ と、核が $J$ である全射 $P \to C$ を選ぶ。
(\ref{crystalline-equation-D}) および (\ref{crystalline-equation-Dn}) を参照すれば、
$D$ と $D(n)$ が得られる。
$T(n)_e = \Spec(D(n)/p^eD(n))$ とおくと、
$(X, T(n)_e, \delta(n))$ は $\text{Cris}(X/S)$ の対象である。
$\mathcal{F}$ を $\mathcal{O}_{X/S}$ 加群の層とし、
$$
M(n) = \lim_e \Gamma((X, T(n)_e, \delta(n)), \mathcal{F})
$$
とおく。ただし $n = 0, 1, 2, 3, \ldots$ である。これは余単体環
$D(0), D(1), D(2), \ldots$ 上の余単体加群をなす。

\begin{proposition}
\label{crystalline-proposition-compute-cohomology}
上の記法のもとで、次を仮定する。
\begin{enumerate}
\item $\mathcal{F}$ は局所的準連接である。
\item $f : T \to T'$ が閉埋め込みであるような
$\text{Cris}(X/S)$ の任意の射
$(U, T, \delta) \to (U', T', \delta')$ に対して、写像
$c_f : f^*\mathcal{F}_{T'} \to \mathcal{F}_T$ は全射である。
\end{enumerate}
このとき複体
$$
M(0) \to M(1) \to M(2) \to \ldots
$$
は $R\Gamma(\text{Cris}(X/S), \mathcal{F})$ を計算する。
\end{proposition}

\begin{proof}
仮定 (1) と補題 \ref{crystalline-lemma-compare} により、
$R\Gamma(\text{Cris}(X/S), \mathcal{F})$ は
$R\Gamma(\mathcal{C}, \mathcal{F})$ と同型である。補題
\ref{crystalline-lemma-compare} と \ref{crystalline-lemma-complete} で用いる圏
$\mathcal{C}$ は一致することに注意する。(2) のような閉埋め込み
$f : T \to T'$ をとる。$c_f : f^*\mathcal{F}_{T'} \to \mathcal{F}_T$
が全射であることは、$\mathcal{F}_{T'} \to f_*\mathcal{F}_T$
が全射であることと同値である。したがって $\mathcal{F}$ が
(1) と (2) を満たすなら、$T'$ 上の準連接
$\mathcal{O}_{T'}$ 加群の短完全列
$$
0 \to \mathcal{K} \to \mathcal{F}_{T'} \to f_*\mathcal{F}_T \to 0
$$
を得る。『スキーム』節 \ref{schemes-section-quasi-coherent}、特に
補題 \ref{schemes-lemma-push-forward-quasi-coherent} を参照されたい。
ゆえに $T'$ がアフィンなら、$H^1(T', \mathcal{K})$ の消滅から
制限写像
$\mathcal{F}(U', T', \delta') \to \mathcal{F}(U, T, \delta)$
は全射である。『スキームのコホモロジー』補題
\ref{coherent-lemma-quasi-coherent-affine-cohomology-zero} を参照されたい。
したがって補題 \ref{crystalline-lemma-complete} に現れる逆系の遷移写像は
全射である。補題 \ref{crystalline-lemma-complete} の $g$ に対して、すべての
$p \geq 1$ について
$R^pg_*(\mathcal{F}|_\mathcal{C}) = 0$ と結論できる。
圏 $\mathcal{C}^\wedge$ の対象 $D$ は、補題
\ref{crystalline-lemma-category-with-covering} の仮定を満たす。これは補題
\ref{crystalline-lemma-generator-completion} による。実際、
$$
D \times \ldots \times D = D(n)
$$
が $\mathcal{C}$ で成り立つ。というのも、補題
\ref{crystalline-lemma-property-Dn} により $D(n)$ は
$\text{Cris}^\wedge(C/A)$ における $D$ の $n + 1$ 重余積だからである。
以上で証明された。
\end{proof}

\begin{lemma}
\label{crystalline-lemma-cohomology-is-zero}
仮定と記法を命題 \ref{crystalline-proposition-compute-cohomology} と同じものとする。
このとき
$$
H^j(\text{Cris}(X/S), \mathcal{F} \otimes_{\mathcal{O}_{X/S}} \Omega^i_{X/S})
= 0
$$
がすべての $i > 0$ およびすべての $j \geq 0$ に対して成り立つ。
\end{lemma}

\begin{proof}
補題 \ref{crystalline-lemma-omega-locally-quasi-coherent} により、
$\mathcal{H} = \mathcal{F} \otimes_{\mathcal{O}_{X/S}} \Omega^i_{X/S}$
も命題 \ref{crystalline-proposition-compute-cohomology} の仮定 (1) と (2) を満たす。
$M(n) = \lim_e M(n)_e$ となるように
$M(n)_e = \Gamma((X, T(n)_e, \delta(n)), \mathcal{F})$ と書く。このとき
\begin{align*}
\lim_e \Gamma((X, T(n)_e, \delta(n)), \mathcal{H}) & =
\lim_e M(n)_e \otimes_{D(n)_e} \Omega_{D(n)}/p^e\Omega_{D(n)} \\
& = \lim_e M(n)_e \otimes_{D(n)} \Omega_{D(n)}
\end{align*}
補題 \ref{crystalline-lemma-vanishing} により、余単体加群
$$
M(0)_e \otimes_{D(0)} \Omega^i_{D(0)} \to
M(1)_e \otimes_{D(1)} \Omega^i_{D(1)} \to
M(2)_e \otimes_{D(2)} \Omega^i_{D(2)} \to \ldots
$$
は零にホモトピックである。遷移写像
$M(n)_{e + 1} \to M(n)_e$ は全射なので、付随する複体の逆極限は
非輪状である\footnote{実際、ホモトピーどうしが両立するため、これらは
零にホモトピックですらあるが、ここでは必要ない。この迂回した議論を
用いる理由は、補題の仮定だけでは
$M(n)_e = M(n)_{e + 1}/p^eM(n)_{e + 1}$ であることが分からないため、
極限 $\lim_e M(n)_e \otimes_{D(n)} \Omega^i_{D(n)}$ が
$M(n) \otimes_{D(n)} \Omega^i_{D(n)}$ の $p$ 進完備化とは限らない
からである。$\mathcal{F}$ がクリスタルなら、これは成り立つ。}。
したがって命題 \ref{crystalline-proposition-compute-cohomology} により
$\mathcal{H}$ のコホモロジーは消滅する。
\end{proof}

\begin{proposition}
\label{crystalline-proposition-compute-cohomology-crystal}
仮定を命題 \ref{crystalline-proposition-compute-cohomology} と同じものとするが、
ここでは $\mathcal{F}$ が準連接加群のクリスタルであると仮定する。
命題 \ref{crystalline-proposition-crystals-on-affine} の対応による、$D$ 上の接続付き加群を
$(M, \nabla)$ とする。このとき複体
$$
M \otimes^\wedge_D \Omega^*_D
$$
は $R\Gamma(\text{Cris}(X/S), \mathcal{F})$ を計算する。
\end{proposition}

\begin{proof}
これを、各項が
$$
K^{a, b} = M \otimes_D^\wedge \Omega^a_{D(b)}
$$
である二重複体 $K^{*, *}$ に付随する二つのスペクトル系列を用いて示す。
これまでに何が分かっているだろうか。補題 \ref{crystalline-lemma-vanishing} により、
$a > 0$ である各列 $K^{a, *}$ は非輪状である。命題
\ref{crystalline-proposition-compute-cohomology} により、第一列 $K^{0, *}$ は
$R\Gamma(\text{Cris}(X/S), \mathcal{F})$ と擬同型である。
したがって二重複体に付随する第一のスペクトル系列から、
$R\Gamma(\text{Cris}(X/S), \mathcal{F})$ と
$\text{Tot}(K^{*, *})$ の間に標準的な擬同型があることが分かる。

\medskip\noindent
次に行 $K^{*, b}$ を考えよう。補題 \ref{crystalline-lemma-structure-Dn} により、
$b + 1$ 個の各写像 $D \to D(b)$ は $D(b)$ を $D$ 上の
分割冪多項式環の $p$ 進完備化として表示する。したがって補題
\ref{crystalline-lemma-relative-poincare} により、写像
$$
M \otimes^\wedge_D\Omega^*_D
\longrightarrow
M \otimes^\wedge_{D(b)} \Omega^*_{D(b)} = K^{*, b}
$$
は擬同型である。これらの写像はいずれもコホモロジー上で
{\it 同じ}写像（さらには導来圏で同じ写像）を定めることに注意する。
実際、その逆写像は余対角写像 $D(b) \to D$、すなわち乗法写像
$P \otimes_A \ldots \otimes_A P \to P$ に対応する写像で与えられる。
したがって第二のスペクトル系列の $E_1$ 頁を見ると、
$$
E_1^{a, b} = H^a(M \otimes^\wedge_D\Omega^*_D)
$$
を得、その微分は
$$
E_1^{a, 0} \xrightarrow{0}
E_1^{a, 1} \xrightarrow{1}
E_1^{a, 2} \xrightarrow{0}
E_1^{a, 3} \xrightarrow{1} \ldots
$$
となる。これは、それぞれが与えられた同一視
$H^a(M \otimes^\wedge_D\Omega^*_D) = E_1^{a, 0} = E_1^{a, 1} = \ldots$
の交代和だからである。ゆえに $E_2$ 頁は第一行では
$H^a(M \otimes^\wedge_D\Omega^*_D)$ に等しく、他では零である。
したがって $M \otimes^\wedge_D\Omega^*_D$ と第一行との同一視は、
$M \otimes^\wedge_D\Omega^*_D$ と $\text{Tot}(K^{*, *})$
との擬同型を誘導する。
\end{proof}

\begin{lemma}
\label{crystalline-lemma-compute-cohomology-crystal-smooth}
仮定を命題 \ref{crystalline-proposition-compute-cohomology-crystal} と同じものとする。
$A \to P'$ は滑らかで、$P' \to C$ は核 $J'$ をもつ全射であるような
環準同型 $A \to P' \to C$ をとる。$D'$ を
$D_{P', \gamma}(J')$ の $p$ 進完備化とする。補題
\ref{crystalline-lemma-crystals-on-affine-smooth} の対応による、
$\mathcal{F}$ に対応する $D'$ 上の対を $(M', \nabla')$ とする。
このとき複体
$$
M' \otimes^\wedge_{D'} \Omega^*_{D'}
$$
は $R\Gamma(\text{Cris}(X/S), \mathcal{F})$ を計算する。
\end{lemma}

\begin{proof}
補題 \ref{crystalline-lemma-crystals-on-affine-smooth} のように
$a : D \to D'$ および $b : D' \to D$ を選ぶ。
接続 $\nabla$ を備えた基底変換 $M = M' \otimes_{D', b} D$ は
$\mathcal{F}$ に対応することに注意する。したがって命題
\ref{crystalline-proposition-compute-cohomology-crystal} により、
$M \otimes^\wedge_D \Omega_D^*$ は $\mathcal{F}$ の結晶コホモロジーを
計算する。それゆえ、$a$ と $b$ が誘導する基底変換写像
$$
M' \otimes^\wedge_{D'} \Omega^*_{D'}
\longrightarrow
M \otimes^\wedge_D \Omega^*_D
\quad\text{および}\quad
M \otimes^\wedge_D \Omega^*_D
\longrightarrow
M' \otimes^\wedge_{D'} \Omega^*_{D'}
$$
が擬同型であることを示せば十分である。$a \circ b = \text{id}_{D'}$
なので、一方の向きの合成は複体
$M' \otimes^\wedge_{D'} \Omega^*_{D'}$ 上の恒等写像である。
したがって、$b \circ a : D \to D$ が誘導する写像
$$
M \otimes^\wedge_D \Omega^*_D
\longrightarrow
M \otimes^\wedge_D \Omega^*_D
$$
が擬同型であることを示せばよい。（両辺の複体が同じであることに
注意する。実際 $M = M' \otimes^\wedge_{D', b} D$ だから、
$M \otimes^\wedge_{D, b \circ a} D =
M' \otimes^\wedge_{D', b \circ a \circ b} D =
M' \otimes^\wedge_{D', b} D = M$ である。）実は、$C$ への
増大写像と両立する任意の分割冪 $A$ 代数準同型
$\rho : D \to D$ に対して、誘導写像
$M \otimes^\wedge_D \Omega^*_D \to M \otimes^\wedge_{D, \rho} \Omega^*_D$
は擬同型であると主張する。

\medskip\noindent
$\rho(x_i) = x_i + z_i$ と書く。$\rho$ は $C$ への増大写像と
両立するので、元 $z_i$ は $D$ の分割冪イデアルに属する。
したがって写像 $\rho$ を合成
$$
D \xrightarrow{\sigma} D\langle \xi_i \rangle^\wedge \xrightarrow{\tau} D
$$
として分解できる。ここで第一の写像は $x_i \mapsto x_i + \xi_i$
で与えられ、第二の写像は $\xi_i$ を $z_i$ に写す分割冪
$D$ 代数準同型である。（ここでは多項式環、分割冪多項式環、
分割冪包、および $p$ 進完備化の普遍性を用いた。）
$\alpha(x_i) = x_i - \xi_i$ および
$\alpha(\xi_i) = \xi_i$ を満たす
$D\langle \xi_i \rangle^\wedge$ の{\it 自己同型} $\alpha$ が存在することに
注意する。$x_i$ を $x_i$ に写す $\alpha \circ \sigma$ に補題
\ref{crystalline-lemma-relative-poincare} を適用し、$\alpha$ が同型であることを用いると、
$\sigma$ は $M \otimes^\wedge_D \Omega^*_D$ と
$M \otimes^\wedge_{D, \sigma} \Omega^*_{D\langle x_i \rangle^\wedge}$
との擬同型を誘導すると結論できる。他方、写像 $\tau$ は写像
$D \to D\langle x_i \rangle^\wedge$, $x_i \mapsto x_i$ を左逆にもち、
補題 \ref{crystalline-lemma-relative-poincare} をもう一度用いると、$\tau$ は
$M \otimes^\wedge_{D, \sigma} \Omega^*_{D\langle x_i \rangle^\wedge}$ と
$M \otimes^\wedge_{D, \tau \circ \sigma} \Omega^*_D$ との擬同型を
誘導すると結論できる。これら二つの擬同型を合成すれば、望むとおり
$\rho$ は擬同型
$M \otimes^\wedge_D \Omega^*_D \to M \otimes^\wedge_{D, \rho} \Omega^*_D$
を誘導する。
\end{proof}










\section{二つの反例}
\label{crystalline-section-examples}

\noindent
結晶コホモロジーがうまく働く例に移る前に、考えているスキームが
基底上固有でないか、または特異である場合に、結晶コホモロジーが
あまりうまく働かない理由を説明する二つの例を挙げよう。
第一の例は \cite{BO} にある。

\begin{example}
\label{crystalline-example-torsion}
$A = \mathbf{Z}_p$ とし、分割冪イデアル $(p)$ にはその一意な
分割冪 $\gamma$ を入れる。
$C = \mathbf{F}_p[x, y]/(x^2, xy, y^2)$ とする。表示
$$
C = P/J = \mathbf{Z}_p[x, y]/(x^2, xy, y^2, p)
$$
を選ぶ。節 \ref{crystalline-section-quasi-coherent-crystals} のように、
分割冪イデアル $(\bar J, \bar \gamma)$ をもつ
$D = D_{P, \gamma}(J)^\wedge$ をとる。$D$ における $x$ と $y$ の
像も、それぞれ $x, y$ と書く。元
$$
\tau = \bar\gamma_p(x^2)\bar\gamma_p(y^2) - \bar\gamma_p(xy)^2 \in D
$$
を考える。$D$ において
$$
p! \bar\gamma_p(x^2) \bar\gamma_p(y^2) =
x^{2p} \bar\gamma_p(y^2) = \bar\gamma_p(x^2y^2) =
x^py^p \bar\gamma_p(xy) = p! \bar\gamma_p(xy)^2
$$
なので、$p\tau = 0$ である。また、
\begin{align*}
\text{d}(\bar\gamma_p(x^2) \bar\gamma_p(y^2))
& =
\bar\gamma_{p - 1}(x^2)\bar\gamma_p(y^2)\text{d}x^2 +
\bar\gamma_p(x^2)\bar\gamma_{p - 1}(y^2)\text{d}y^2 \\
& =
2 x \bar\gamma_{p - 1}(x^2)\bar\gamma_p(y^2)\text{d}x +
2 y \bar\gamma_p(x^2)\bar\gamma_{p - 1}(y^2)\text{d}y \\
& =
2/(p - 1)!( x^{2p - 1} \bar\gamma_p(y^2)\text{d}x +
y^{2p - 1} \bar\gamma_p(x^2)\text{d}y ) \\
& =
2/(p - 1)!
(x^{p - 1} \bar\gamma_p(xy^2)\text{d}x +
y^{p - 1} \bar\gamma_p(x^2y)\text{d}y) \\
& =
2/(p - 1)!
(x^{p - 1}y^p \bar\gamma_p(xy)\text{d}x +
x^py^{p - 1} \bar\gamma_p(xy)\text{d}y) \\
& =
2 \bar\gamma_{p - 1}(xy) \bar\gamma_p(xy)(y\text{d}x + x \text{d}y) \\
& = 
\text{d}(\bar\gamma_p(xy)^2)
\end{align*}
であるから、$\Omega_D$ において $\text{d}\tau = 0$ であることにも
注意する。最後に、$D$ において $\tau \not = 0$ であると主張する。
これを見るには、$\tau$ の $B$ における像が零でないような
$\text{Cris}(C/S)$ の対象
$(B \to \mathbf{F}_p[x, y]/(x^2, xy, y^2), \delta)$ を構成すれば十分である。
そのために
$$
B = \mathbf{F}_p[x, y, u, v]/(x^3, x^2y, xy^2, y^3, xu, yu, xv, yv, u^2, v^2)
$$
とし、$C$ への明らかな全射を入れる。$K = \Ker(B \to C)$ とし、写像
$$
\delta_p : K \longrightarrow K,\quad
ax^2 + bxy + cy^2 + du + ev + fuv \longmapsto a^pu + c^pv
$$
を考える。これは『分割冪代数』補題
\ref{dpa-lemma-need-only-gamma-p} の仮定 (1), (2), (3) を満たすことが
確かめられ、したがって分割冪構造を定める。さらに、$\tau$ は
$B$ で零でない $uv$ に写る。$X = \Spec(C)$ および
$S = \Spec(A)$ とおく。次を結論できる。
\begin{enumerate}
\item $H^0(\text{Cris}(X/S), \mathcal{O}_{X/S})$ は $p$ ねじれをもつ。
\item Frobenius による引き戻し
$F^* : H^0(\text{Cris}(X/S), \mathcal{O}_{X/S})
\to H^0(\text{Cris}(X/S), \mathcal{O}_{X/S})$ は単射でない。
\end{enumerate}
実際、命題 \ref{crystalline-proposition-compute-cohomology-crystal} により、
$\tau$ は $H^0(\text{Cris}(X/S), \mathcal{O}_{X/S})$ の零でない
ねじれ元を定める。同様に、$\sigma : D \to D$ を $P$ 上の
Frobenius の任意の持ち上げが誘導する写像とすると、
$F^*(\tau) = \sigma(\tau)$ である。$\sigma(x) = x^p$ および
$\sigma(y) = y^p$ と選べば、簡単な計算から $F^*(\tau) = 0$ が分かる。
\end{example}

\noindent
次の例は、アフィン $n$ 空間についてさえ結晶コホモロジーが
期待されるものを与えないことを示す。

\begin{example}
\label{crystalline-example-affine-n-space}
$A = \mathbf{Z}_p$ とし、分割冪イデアル $(p)$ にはその一意な
分割冪 $\gamma$ を入れる。
$C = \mathbf{F}_p[x_1, \ldots, x_r]$ とする。表示
$$
C = P/J = P/pP\quad\text{ただし}\quad P = \mathbf{Z}_p[x_1, \ldots, x_r]
$$
を選ぶ。『分割冪代数』補題 \ref{dpa-lemma-gamma-extends} により、
$pP$ は分割冪をもつことに注意する。したがって、分割冪イデアル
$(p)$ をもつ $D = P^\wedge$ とおけば、節
\ref{crystalline-section-quasi-coherent-crystals} の状況が得られる。
$R\Gamma(\text{Cris}(X/S), \mathcal{O}_{X/S})$ は複体
$$
D \to \Omega^1_D \to \Omega^2_D \to \ldots \to \Omega^r_D
$$
によって表される。命題 \ref{crystalline-proposition-compute-cohomology-crystal} を
参照されたい。$r > 0$ と仮定すると、次を結論できる。
\begin{enumerate}
\item $S = \Spec(\mathbf{Z}_p)$ 上の
$X = \mathbf{A}^r_{\mathbf{F}_p}$ の結晶構造層の結晶コホモロジーは、
次数 $0, \ldots, r$ 以外では零である。
\item $H^0(\text{Cris}(X/S), \mathcal{O}_{X/S}) = \mathbf{Z}_p$ である。
\item コホモロジー群 $H^r(\text{Cris}(X/S), \mathcal{O}_{X/S})$ は
無限であり、ねじれアーベル群ではない。
\item コホモロジー群 $H^r(\text{Cris}(X/S), \mathcal{O}_{X/S})$ は
$p$ 進位相について分離的でない。
\end{enumerate}
最初の二つの主張は妥当だが、(3) と (4) は驚くべきものである。
これらの主張は、上に表示した複体を具体的に書けば直ちに従う。
$r = 1$ の場合についてこれを行おう。このとき考えるのは
$p$ 進完備加群の二項複体
$$
\text{d} :
D = \left(
\bigoplus\nolimits_{n \geq 0} \mathbf{Z}_p x^n
\right)^\wedge
\longrightarrow
\Omega^1_D = \left(
\bigoplus\nolimits_{n \geq 1} \mathbf{Z}_p x^{n - 1}\text{d}x
\right)^\wedge
$$
だけである。この写像は $\text{diag}(0, 1, 2, 3, 4, \ldots)$ で与えられるが、
右辺には第一直和因子がない。ここから
$\bigoplus_{n > 0} \mathbf{Z}_p/n\mathbf{Z}_p$ が余核の部分群であることは
明らかであり、したがって余核は無限である。実際、元
$$
\omega = \sum\nolimits_{e > 0} p^e x^{p^{2e} - 1}\text{d}x
$$
は明らかに余核のねじれ元ではない。しかし事情はさらに悪い。元
$$
\eta = \sum\nolimits_{e > 0} p^e x^{p^e - 1}\text{d}x
$$
を考える。すべての $t > 0$ に対して、$\eta$ は $\text{d}$ の像を法として
$\sum_{e > t} p^e x^{p^e - 1}\text{d}x$ と合同であり、この元は
$p^t$ で割り切れる。したがって
$H^1(\text{Cris}(X/S), \mathcal{O}_{X/S})$ における $\eta$ の
コホモロジー類は
$\bigcap p^tH^1(\text{Cris}(X/S), \mathcal{O}_{X/S})$ に含まれる。
しかし $\eta$ は $\text{d}$ の像には属さない。像に属するとすれば、
ある $a \in \mathbf{Z}_p$ に対する
$a + \sum_{e > 0} x^{p^e}$ の像でなければならないが、これは左辺の
加群の元ではないからである。ゆえに $\eta$ のコホモロジー類は零でなく、
(4) が成り立つ。（実際、$\eta$ のコホモロジー類は非ねじれである。）
肯定的な面として、コホモロジー群
$H^i(\text{Cris}(X/S), \mathcal{O}_{X/S})$ は $p$ 進完備加群の
複体のコホモロジーであるため、導来的完備な $\mathbf{Z}_p$ 加群である。
『代数詳論』節 \ref{more-algebra-section-derived-completion} を参照されたい。
\end{example}








\section{応用}
\label{crystalline-section-applications}

\noindent
この節では、前節までの内容のいくつかの応用をまとめる。

\begin{proposition}
\label{crystalline-proposition-compare-with-de-Rham}
状況 \ref{crystalline-situation-global} において、$\mathcal{F}$ を
$\text{Cris}(X/S)$ 上の準連接加群のクリスタルとする。複体の切断写像
$$
(\mathcal{F} \to
\mathcal{F} \otimes_{\mathcal{O}_{X/S}} \Omega^1_{X/S} \to
\mathcal{F} \otimes_{\mathcal{O}_{X/S}} \Omega^2_{X/S} \to \ldots)
\longrightarrow \mathcal{F}[0],
$$
は擬同型ではないが、$Ru_{X/S, *}$ を適用すると擬同型になる。
実際、任意の $i > 0$ に対して
$$
Ru_{X/S, *}(\mathcal{F} \otimes_{\mathcal{O}_{X/S}} \Omega^i_{X/S}) = 0.
$$
が成り立つ。
\end{proposition}

\begin{proof}
補題 \ref{crystalline-lemma-automatic-connection} により、命題に示した de Rham 複体が
得られる。$\mathcal{H} = \mathcal{F} \otimes \Omega^i_{X/S}$ と略記する。
$X' \subset X$ を、あるアフィン開部分スキーム $S' \subset S$ に写る
アフィン開部分スキームとする。このとき
$$
(Ru_{X/S, *}\mathcal{H})|_{X'_{Zar}} =
Ru_{X'/S', *}(\mathcal{H}|_{\text{Cris}(X'/S')}),
$$
である。補題 \ref{crystalline-lemma-localize} を参照されたい。したがって補題
\ref{crystalline-lemma-cohomology-is-zero} により、$Ru_{X/S, *}\mathcal{H}$ は
$X_{Zar}$ 上の層の複体であって、その任意のアフィン開集合上の
コホモロジーは自明である。$X$ の位相はアフィン開集合からなる基底を
もつので、$Ru_{X/S, *}\mathcal{H}$ は零と擬同型である。
\end{proof}

\begin{remark}
\label{crystalline-remark-vanishing}
命題 \ref{crystalline-proposition-compare-with-de-Rham} の証明から、結論
$$
Ru_{X/S, *}(\mathcal{F} \otimes_{\mathcal{O}_{X/S}} \Omega^i_{X/S}) = 0
$$
が、命題 \ref{crystalline-proposition-compute-cohomology} の条件 (1) と (2) を満たす
任意の $\mathcal{O}_{X/S}$ 加群 $\mathcal{F}$ について、$i > 0$ なら
成り立つことが分かる。これは次のクリスタルでない層にも適用できる。
すなわち、すべての $i$ に対する $\Omega^i_{X/S}$、および
$\mathcal{F}$ が準連接 $\mathcal{O}_X$ 加群であるときの
$\underline{\mathcal{F}}$ という形の任意の層である。特に層
$\underline{\mathcal{O}_X} = \underline{\mathbf{G}_a}$ に適用できる。
しかし de Rham 複体を構成するには、$\mathcal{F}$ がクリスタルであることを
必要とする補題 \ref{crystalline-lemma-automatic-connection} のような結果が要ることに
注意する。したがって命題 \ref{crystalline-proposition-compare-with-de-Rham} の完全な
主張が成り立つ加群の層の集まりは、（現在のところ）ちょうど
準連接加群のクリスタルの圏である。
\end{remark}

\noindent
状況 \ref{crystalline-situation-global} において、$\mathcal{F}$ を
$\text{Cris}(X/S)$ 上の準連接加群のクリスタルとする。
$(U, T, \delta)$ を $\text{Cris}(X/S)$ の対象とする。命題
\ref{crystalline-proposition-compare-with-de-Rham} により、標準的な写像
\begin{equation}
\label{crystalline-equation-restriction}
R\Gamma(\text{Cris}(X/S), \mathcal{F})
\longrightarrow
R\Gamma(T, \mathcal{F}_T \otimes_{\mathcal{O}_T} \Omega^*_{T/S, \delta})
\end{equation}
を構成できる。実際、
$R\Gamma(\text{Cris}(X/S), \mathcal{F}) =
R\Gamma(\text{Cris}(X/S), \mathcal{F} \otimes \Omega^*_{X/S})$ であり、
大域コホモロジー類を $T$ に制限でき、さらに補題
\ref{crystalline-lemma-module-of-differentials} により $\Omega_{X/S}$ は
$\Omega_{T/S, \delta}$ に制限される。

















\section{いくつかの追加結果}
\label{crystalline-section-missing}

\noindent
この節では、証明がまだ欠けているいくつかの結果に言及する。
これらを一連の注意として述べ、詳しい証明を加えた時点で初めて
実際の補題や命題に改める。

\begin{remark}[高次順像]
\label{crystalline-remark-compute-direct-image}
$p$ を素数とする。
$(S, \mathcal{I}, \gamma) \to (S', \mathcal{I}', \gamma')$ を
$\mathbf{Z}_{(p)}$ 上の分割冪スキームの射とする。スキームの射の可換図式
$$
\xymatrix{
X \ar[r]_f \ar[d] & X' \ar[d] \\
S_0 \ar[r] & S'_0
}
$$
をとり、$X$ と $X'$ 上で $p$ は局所的に冪零であると仮定する。
$\mathcal{F}$ を $\text{Cris}(X/S)$ 上の $\mathcal{O}_{X/S}$ 加群とする。
このとき $Rf_{\text{cris}, *}\mathcal{F}$ は次のように計算できる。

\medskip\noindent
$\text{Cris}(X'/S')$ の対象 $(U', T', \delta')$ が与えられたとき、
$U = X \times_{X'} U' = f^{-1}(U')$ とおく。これは $X$ の開部分スキームである。
$(T_0, T, \delta)$ を、次の図式が分割冪スキームの圏で Cartesian となる
$S$ 上の分割冪スキームとする。
$$
\xymatrix{
T \ar[r] \ar[d] & T' \ar[d] \\
S \ar[r] & S'
}
$$
補題 \ref{crystalline-lemma-fibre-product} を参照されたい。誘導される射
$U \to T_0$ があり、注意 \ref{crystalline-remark-functoriality-cris} により射
$(U/T)_{\text{cris}} \to (X/S)_{\text{cris}}$ が得られる。
$\mathcal{F}_U$ を $\mathcal{F}$ の引き戻しとする。
$\tau_{U/T} : (U/T)_{\text{cris}} \to T_{Zar}$ を構造射とする。このとき
\begin{equation}
\label{crystalline-equation-identify-pushforward}
\left(Rf_{\text{cris}, *}\mathcal{F}\right)_{T'} =
R(T \to T')_*\left(R\tau_{U/T, *} \mathcal{F}_U \right)
\end{equation}
である。ここで左辺は制限である（節 \ref{crystalline-section-sheaves} を参照）。

\medskip\noindent
ヒント：まず $\text{Cris}(U/T)$ は、集合の層
$f_{\text{cris}}^{-1}h_{(U', T', \delta')}$ における
$\text{Cris}(X/S)$ の局所化であることを示す。ここで局所化は
『サイト』補題 \ref{sites-lemma-localize-topos-site} の意味である。
次に、単射的 $\mathcal{O}_{X/S}$ 加群の引き戻しが
$\text{Cris}(U/T)$ 上の単射的 $\mathcal{O}_{U/T}$ 加群であることを用い、
$\mathcal{F}$ が単射的加群である場合と加群の順像の場合に帰着する。
最後に、通常の順像について結果が成り立つことを確かめる。
\end{remark}

\begin{remark}[Mayer-Vietoris]
\label{crystalline-remark-mayer-vietoris}
注意 \ref{crystalline-remark-compute-direct-image} の状況で、開被覆
$X = X' \cup X''$ があるとする。$X''' = X' \cap X''$ とおく。
$f'$, $f''$, $f''$ を、それぞれ $f$ の $X'$, $X''$, $X'''$ への制限とする。
また $\mathcal{F}'$, $\mathcal{F}''$, $\mathcal{F}'''$ を、それぞれ
$\mathcal{F}$ の $X'$, $X''$, $X'''$ の結晶サイトへの制限とする。
このとき $D(\mathcal{O}_{X'/S'})$ における識別三角
$$
Rf_{\text{cris}, *}\mathcal{F}
\longrightarrow
Rf'_{\text{cris}, *}\mathcal{F}' \oplus Rf''_{\text{cris}, *}\mathcal{F}''
\longrightarrow
Rf'''_{\text{cris}, *}\mathcal{F}'''
\longrightarrow
Rf_{\text{cris}, *}\mathcal{F}[1]
$$
が存在する。

\medskip\noindent
ヒント：部分圏 $\text{Cris}(X'/S)$, $\text{Cris}(X''/S)$,
$\text{Cris}(X'''/S)$ が $\text{Cris}(X/S)$ 上の終対象層の開部分対象に
対応し、最後のものが最初の二つの共通部分であることから形式的に従う。
\end{remark}

\begin{remark}[{\v C}ech 複体]
\label{crystalline-remark-cech-complex}
$p$ を素数とする。$(A, I, \gamma)$ を分割冪環とし、$A$ は
$\mathbf{Z}_{(p)}$ 代数であるとする。$S = \Spec(A)$ および
$S_0 = \Spec(A/I)$ とおく。$X$ を $S_0$ 上の分離的スキーム\footnote{この仮定は
厳密には必要ない。超被覆を用いれば、この注意の構成を一般の場合へ
拡張できる。}とし、$X$ 上で $p$ は局所的に冪零であるとする。
$\mathcal{F}$ を準連接 $\mathcal{O}_{X/S}$ 加群のクリスタルとする。

\medskip\noindent
$X$ のアフィン開被覆 $X = \bigcup_{\lambda \in \Lambda} U_\lambda$ を選ぶ。
$U_\lambda = \Spec(C_\lambda)$ と書く。$A$ 上の多項式代数
$P_\lambda$ と全射 $P_\lambda \to C_\lambda$ を選ぶ。これらの選択を
固定すると、$R\Gamma(\text{Cris}(X/S), \mathcal{F})$ を計算する
{\v C}ech 複体を構成できる。

\medskip\noindent
$n \geq 0$ および $\lambda_0, \ldots, \lambda_n \in \Lambda$ に対して
$U_{\lambda_0 \ldots \lambda_n} = U_{\lambda_0} \cap \ldots
\cap U_{\lambda_n}$ と書く。これは仮定によりアフィンスキームである。
$U_{\lambda_0 \ldots \lambda_n} = \Spec(C_{\lambda_0 \ldots \lambda_n})$
と書く。また
$$
P_{\lambda_0 \ldots \lambda_n} =
P_{\lambda_0} \otimes_A \ldots \otimes_A P_{\lambda_n}
$$
とおく。これには $C_{\lambda_0 \ldots \lambda_n}$ への標準的な全射がある。
その核を $J_{\lambda_0 \ldots \lambda_n}$ と書き、
$D_{\lambda_0 \ldots \lambda_n}$ を、$\gamma$ に関する
$P_{\lambda_0 \ldots \lambda_n}$ 内の
$J_{\lambda_0 \ldots \lambda_n}$ の分割冪包の $p$ 進完備化とする。
$M_{\lambda_0 \ldots \lambda_n}$ を、命題
\ref{crystalline-proposition-crystals-on-affine} を介して $\mathcal{F}$ の
$\text{Cris}(U_{\lambda_0 \ldots \lambda_n}/S)$ への制限に対応する
$P_{\lambda_0 \ldots \lambda_n}$ 加群とする。構成により、次数 $n$ の環が
$$
D(n) =
\prod\nolimits_{\lambda_0 \ldots \lambda_n}
D_{\lambda_0 \ldots \lambda_n}
$$
である余単体分割冪環 $D(*)$ を得る。ここでは分割冪包の関手性と、
同様に定義した環 $P(*)$ 上の自明な余単体構造を用いる。
$M_{\lambda_0 \ldots \lambda_n}$ は対象
$\Spec(D_{\lambda_0 \ldots \lambda_n})$ における $\mathcal{F}$ の「値」なので、
規則
$$
M(n) = \prod\nolimits_{\lambda_0 \ldots \lambda_n}
M_{\lambda_0 \ldots \lambda_n}
$$
で定める $M(*)$ は余単体 $D(*)$ 加群をなす。ここで
$$
R\Gamma(\text{Cris}(X/S), \mathcal{F}) = s(M(*))
$$
と主張する。ここで $s(-)$ は余単体加群に付随する余鎖複体を表す。
『単体論』節 \ref{simplicial-section-dold-kan-cosimplicial} を参照されたい。

\medskip\noindent
ヒント：これは命題 \ref{crystalline-proposition-compute-cohomology} の証明と同様である。
特に、その命題の仮定を満たす任意の加群について結果は成り立つ。
\end{remark}

\begin{remark}[交代 {\v C}ech 複体]
\label{crystalline-remark-alternating-cech-complex}
$p$ を素数とする。$(A, I, \gamma)$ を分割冪環とし、$A$ は
$\mathbf{Z}_{(p)}$ 代数であるとする。$S = \Spec(A)$ および
$S_0 = \Spec(A/I)$ とおく。$X$ を $S_0$ 上の分離的かつ準コンパクトな
スキームとし、$X$ 上で $p$ は局所的に冪零であるとする。
$\mathcal{F}$ を準連接 $\mathcal{O}_{X/S}$ 加群のクリスタルとする。

\medskip\noindent
$X$ の有限アフィン開被覆
$X = \bigcup_{\lambda \in \Lambda} U_\lambda$ と、$\Lambda$ 上の全順序を選ぶ。
$U_\lambda = \Spec(C_\lambda)$ と書く。$A$ 上の多項式代数
$P_\lambda$ と全射 $P_\lambda \to C_\lambda$ を選ぶ。これらの選択を
固定すると、$R\Gamma(\text{Cris}(X/S), \mathcal{F})$ を計算する
交代 {\v C}ech 複体を構成できる。

\medskip\noindent
注意 \ref{crystalline-remark-cech-complex} で導入した記法を用いる。
$\Omega_{\lambda_0 \ldots \lambda_n}$ を、分割冪構造と両立する
$D_{\lambda_0 \ldots \lambda_n}$ の $A$ 上の微分加群の
$p$ 進完備化とする。$\nabla$ を、命題
\ref{crystalline-proposition-crystals-on-affine} から得られる
$M_{\lambda_0 \ldots \lambda_n}$ 上の可積分接続とする。
各項が
$$
M^{n, m} =
\bigoplus\nolimits_{\lambda_0 < \ldots < \lambda_n}
M_{\lambda_0 \ldots \lambda_n}
\otimes^\wedge_{D_{\lambda_0 \ldots \lambda_n}}
\Omega^m_{D_{\lambda_0 \ldots \lambda_n}}.
$$
である二重複体 $M^{\bullet, \bullet}$ を考える。$n$ を増加させる微分
$d_1$ には通常の {\v C}ech 微分を用い、微分 $d_2$ には接続、すなわち
de Rham 複体の微分を用いる。次を主張する。
$$
R\Gamma(\text{Cris}(X/S), \mathcal{F}) = \text{Tot}(M^{\bullet, \bullet})
$$
ここで $\text{Tot}(-)$ は二重複体に付随する全複体を表す。
『ホモロジー』定義 \ref{homology-definition-associated-simple-complex} を
参照されたい。

\medskip\noindent
ヒント：命題 \ref{crystalline-proposition-compare-with-de-Rham} により
$$
R\Gamma(\text{Cris}(X/S), \mathcal{F}) = R\Gamma(\text{Cris}(X/S),
\mathcal{F} \otimes_{\mathcal{O}_{X/S}} \Omega_{X/S}^\bullet)
$$
である。式の右辺は、被覆
$X = \bigcup_{\lambda \in \Lambda} U_\lambda$ と複体
$\mathcal{F} \otimes_{\mathcal{O}_{X/S}} \Omega_{X/S}^\bullet$ に対する
交代 {\v C}ech 複体にほかならない。この被覆は $\text{Cris}(X/S)$ の
終対象層の開被覆を誘導する。命題
\ref{crystalline-proposition-compute-cohomology-crystal} を参照されたい。
ここで、すべての部分上で正しい答えを与えるなら交代 {\v C}ech 複体が
コホモロジーを計算するという、サイト上のコホモロジーの一般結果から
結論が従う（将来の参照をここに挿入）。
\end{remark}

\begin{remark}[準連接性]
\label{crystalline-remark-quasi-coherent}
注意 \ref{crystalline-remark-compute-direct-image} の状況で、$S \to S'$ および
$X \to S_0$ は準コンパクトかつ準分離的であると仮定する。このとき
準連接 $\mathcal{O}_{X/S}$ 加群のクリスタル $\mathcal{F}$ に対して、層
$R^if_{\text{cris}, *}\mathcal{F}$ は局所的準連接である。

\medskip\noindent
ヒント：$(U', T', \delta')$ を $\text{Cris}(X'/S')$ の任意の対象としたとき、
$T'$ への制限が準連接 $\mathcal{O}_{T'}$ 加群であることを示す必要がある。
$T'$ がアフィンの場合に示せば十分である。式
(\ref{crystalline-equation-identify-pushforward})、$T \to T'$ が準コンパクトかつ
準分離的であること（$T$ は $S \to S'$ による $T'$ の基底変換上
アフィンである）、および『スキームのコホモロジー』補題
\ref{coherent-lemma-quasi-coherence-higher-direct-images} を用いると、層
$R^i\tau_{U/T, *}\mathcal{F}_U$ が準連接であることを示せば十分だと分かる。
$U \to T_0$ も準コンパクトかつ準分離的であることに注意する。
『スキーム』補題 \ref{schemes-lemma-quasi-compact-permanence} および
\ref{schemes-lemma-quasi-compact-permanence} を参照されたい。

\medskip\noindent
これにより、$S$ 上で $p$ が局所的に冪零である場合に、
$R^i\tau_{X/S, *}\mathcal{F}$ が $S$ 上準連接であることを示す問題に
帰着する。ここで $\tau_{X/S}$ は構造射である。注意
\ref{crystalline-remark-structure-morphism} を参照されたい。$S$ 上局所的に
議論できるので、$S$ はアフィンと仮定してよい（補題
\ref{crystalline-lemma-localize} を参照）。$X$ を被覆するアフィンの個数に関する帰納法と
Mayer--Vietoris（注意 \ref{crystalline-remark-mayer-vietoris}）により、$X$ も
アフィンである場合に帰着する。これは『スキームのコホモロジー』補題
\ref{coherent-lemma-quasi-coherence-higher-direct-images} の証明と同様である。
$X = \Spec(C)$ および $S = \Spec(A)$ とし、$(A, I, \gamma)$ と
$A \to C$ は状況 \ref{crystalline-situation-affine} のとおりであるとする。
節 \ref{crystalline-section-quasi-coherent-crystals} のように、$A$ 上の多項式代数
$P$ と全射 $P \to C$ を選ぶ。命題
\ref{crystalline-proposition-crystals-on-affine} により、$\mathcal{F}$ に対応する加群を
$(M, \nabla)$ とする。命題
\ref{crystalline-proposition-compute-cohomology-crystal} を適用すると、
$R\Gamma(\text{Cris}(X/S), \mathcal{F})$ は
$M \otimes_D \Omega_D^*$ によって表される。$A$ で $p$ は冪零なので、
完備化は不要であることに注意する。これが
$S = \Spec(A)$ の主開集合をとることと両立することを示す必要がある。
$g \in A$ とする。同様に
$R\Gamma(\text{Cris}(X_g/S_g), \mathcal{F})$ は
$M_g \otimes_{D_g} \Omega_{D_g}^*$ によって計算される。ここでも
$p$ 進完備化が不要であることを用いる。局所化は $A$ 加群上の完全関手なので、
これで結論が従う。
\end{remark}

\begin{remark}[有界性]
\label{crystalline-remark-bounded-cohomology}
注意 \ref{crystalline-remark-compute-direct-image} の状況で、$S \to S'$ は
準コンパクトかつ準分離的であり、$X \to S_0$ は有限型かつ準分離的であると
仮定する。このとき整数 $i_0$ が存在し、任意の準連接
$\mathcal{O}_{X/S}$ 加群のクリスタル $\mathcal{F}$ に対して、すべての
$i > i_0$ について $R^if_{\text{cris}, *}\mathcal{F} = 0$ である。

\medskip\noindent
ヒント：注意 \ref{crystalline-remark-quasi-coherent} と同様に議論し、
『スキームのコホモロジー』補題
\ref{coherent-lemma-quasi-coherence-higher-direct-images} を用いると、命題
\ref{crystalline-proposition-compute-cohomology-crystal} の状況で $C$ が $A$ 上
有限型の代数であるとき、$i \gg 0$ に対して
$H^i(\text{Cris}(X/S), \mathcal{F}) = 0$ であることを示す問題に帰着する。
有限変数多項式代数を選べば $i \gg 0$ に対して
$\Omega^i_D = 0$ となるので、これは明らかである。
\end{remark}

\begin{remark}[具体的な有界性]
\label{crystalline-remark-bounded-cohomology-over-point}
状況 \ref{crystalline-situation-global} において、$\mathcal{F}$ を準連接
$\mathcal{O}_{X/S}$ 加群のクリスタルとする。$S_0$ はただ一つの点をもち、
$X \to S_0$ は有限表示であると仮定する。
\begin{enumerate}
\item $\dim X = d$ であり、$X/S_0$ の埋め込み次元が $e$ なら、
$i > d + e$ に対して $H^i(\text{Cris}(X/S), \mathcal{F}) = 0$ である。
\item $X$ が分離的で $q$ 個のアフィンによって被覆でき、
$X/S_0$ の埋め込み次元が $e$ なら、$i > q + e$ に対して
$H^i(\text{Cris}(X/S), \mathcal{F}) = 0$ である。
\end{enumerate}
ヒント：(1) の場合には
$$
H^i(\text{Cris}(X/S), \mathcal{F}) = H^i(X_{Zar}, Ru_{X/S, *}\mathcal{F})
$$
であることと、$Ru_{X/S, *}\mathcal{F}$ が局所的には $X$ を $S$ 上
次元 $e$ の滑らかなスキームへ埋め込んで構成する de Rham 複体により
計算されることを用いる（補題
\ref{crystalline-lemma-compute-cohomology-crystal-smooth} を参照）。これらの
de Rham 複体はすべての次数 $> e$ で零である。したがって (1) は
『コホモロジー』命題 \ref{cohomology-proposition-vanishing-Noetherian}
から従う。(2) の場合には交代 {\v C}ech 複体（注意
\ref{crystalline-remark-alternating-cech-complex} を参照）を用いて、$X$ がアフィンの
場合に帰着する。アフィンの場合には、$X$ を $S$ 上次元 $e$ の
滑らかなスキームへ埋め込むことに付随する de Rham 複体を用いて示す。
そのような埋め込みの構成には多少の作業が必要である。
\end{remark}

\begin{remark}[基底変換写像]
\label{crystalline-remark-base-change}
注意 \ref{crystalline-remark-compute-direct-image} の状況で、
$S = \Spec(A)$ および $S' = \Spec(A')$ はアフィンであると仮定する。
$\mathcal{F}'$ を $\mathcal{O}_{X'/S'}$ 加群とし、$\mathcal{F}$ を
$\mathcal{F}'$ の引き戻しとする。このとき標準的な基底変換写像
$$
L(S' \to S)^*R\tau_{X'/S', *}\mathcal{F}'
\longrightarrow
R\tau_{X/S, *}\mathcal{F}
$$
がある。ここで $\tau_{X/S}$ と $\tau_{X'/S'}$ は構造射である。
注意 \ref{crystalline-remark-structure-morphism} を参照されたい。大域切断をとると、
$D(A)$ における基底変換写像
\begin{equation}
\label{crystalline-equation-base-change-map}
R\Gamma(\text{Cris}(X'/S'), \mathcal{F}') \otimes^\mathbf{L}_{A'} A
\longrightarrow
R\Gamma(\text{Cris}(X/S), \mathcal{F})
\end{equation}
が得られる。

\medskip\noindent
ヒント：『サイト上のコホモロジー』注意
\ref{sites-cohomology-remark-base-change} の非常に一般的な基底変換写像と、
標準写像 $Lf_{\text{cris}}^*\mathcal{F}' \to
f_{\text{cris}}^*\mathcal{F}' = \mathcal{F}$ を合成する。
\end{remark}

\begin{remark}[基底変換同型]
\label{crystalline-remark-base-change-isomorphism}
次の条件がすべて満たされるなら、写像
(\ref{crystalline-equation-base-change-map}) は同型である。
\begin{enumerate}
\item $p$ は $A'$ で冪零である。
\item $\mathcal{F}'$ は準連接 $\mathcal{O}_{X'/S'}$ 加群のクリスタルである。
\item $X' \to S'_0$ は準コンパクトかつ準分離的な射である。
\item $X = X' \times_{S'_0} S_0$,
\item $\mathcal{F}'$ は平坦な $\mathcal{O}_{X'/S'}$ 加群である。
\item $X' \to S'_0$ は局所完全交叉射である（『射の詳論』定義
\ref{more-morphisms-definition-lci} を参照）。例えば
$X' \to S'_0$ がシントミックまたは滑らかなら、これは成り立つ。
\item $X'$ と $S_0$ は $S'_0$ 上 Tor 独立である（『代数詳論』定義
\ref{more-algebra-definition-tor-independent} を参照）。例えば
$S_0 \to S'_0$ または $X' \to S'_0$ のいずれかが平坦なら、これは成り立つ。
\end{enumerate}
ヒント：条件 (1) により、以下の議論では $p$ 進完備化は何も変えず、
無視できる。条件 (3) と Mayer--Vietoris（注意
\ref{crystalline-remark-mayer-vietoris} を参照）を用いると、$X'$ がアフィンの
場合に帰着する。実際、条件 (6) によりさらに縮小すれば、
$X' = \Spec(C')$ であり、表示
$C' = A'/I'[x_1, \ldots, x_n]/(\bar f'_1, \ldots, \bar f'_c)$
が与えられ、$\bar f'_1, \ldots, \bar f'_c$ は $A'/I'$ における
Koszul 正則列であると仮定できる。（これは、滑らか位相局所的に
$\bar f'_1, \ldots, \bar f'_c$ が正則列をなすという意味である。
『代数詳論』補題
\ref{more-algebra-lemma-Koszul-regular-flat-locally-regular} を参照。）
$\bar f'_i$ の持ち上げ $f'_i \in A'[x_1, \ldots, x_n]$ を選ぶ。
(4) により $X = \Spec(C)$ であり、
$C = A/I[x_1, \ldots, x_n]/(\bar f_1, \ldots, \bar f_c)$ となる。
ここで $f_i \in A[x_1, \ldots, x_n]$ は $f'_i$ の像である。
性質 (7) により $\bar f_1, \ldots, \bar f_c$ は
$A/I[x_1, \ldots, x_n]$ における Koszul 正則列である。
$\gamma'$ に関する、$A'[x_1, \ldots, x_n]$ 内の
$I'A'[x_1, \ldots, x_n] + (f'_1, \ldots, f'_c)$ の分割冪包は
$$
D' = A'[x_1, \ldots, x_n]\langle \xi_1, \ldots, \xi_c \rangle/(\xi_i - f'_i)
$$
である。補題 \ref{crystalline-lemma-describe-divided-power-envelope} を参照されたい。
次に $\xi_1 - f'_1, \ldots, \xi_n - f'_n$ が環
$A'[x_1, \ldots, x_n]\langle \xi_1, \ldots, \xi_c\rangle$ における
Koszul 正則列であることを確かめる。同様に、$\gamma$ に関する
$A[x_1, \ldots, x_n]$ 内の
$IA[x_1, \ldots, x_n] + (f_1, \ldots, f_c)$ の分割冪包は
$$
D = A[x_1, \ldots, x_n]\langle \xi_1, \ldots, \xi_c\rangle/(\xi_i - f_i)
$$
であり、$\xi_1 - f_1, \ldots, \xi_n - f_n$ は環
$A[x_1, \ldots, x_n]\langle \xi_1, \ldots, \xi_c\rangle$ における
Koszul 正則列である。したがって $D' \otimes_{A'}^\mathbf{L} A = D$ である。
条件 (2) により、$\mathcal{F}'$ は接続付き $D'$ 加群からなる対
$(M', \nabla)$ に対応する。命題 \ref{crystalline-proposition-crystals-on-affine} を
参照されたい。このとき $M = M' \otimes_{D'} D$ は引き戻し
$\mathcal{F}$ に対応する。仮定 (5) により $M'$ は平坦な $D'$ 加群なので、
$$
M = M' \otimes_{D'} D = M' \otimes_{D'} D' \otimes_{A'}^\mathbf{L} A
= M' \otimes_{A'}^\mathbf{L} A
$$
である。節 \ref{crystalline-section-quasi-coherent-crystals} で定義した微分加群
$\Omega_{D'}$ と $\Omega_D$ は同じ生成元上の自由 $D'$ 加群なので、
$$
M \otimes_D \Omega^\bullet_D =
M' \otimes_{D'} \Omega^\bullet_{D'} \otimes_{D'} D =
M' \otimes_{D'} \Omega^\bullet_{D'} \otimes_{A'}^\mathbf{L} A
$$
を得る。命題 \ref{crystalline-proposition-compute-cohomology-crystal} により、
これは求める主張を証明する。
\end{remark}

\begin{remark}[導来逆極限]
\label{crystalline-remark-rlim}
$p$ を素数とする。$(A, I, \gamma)$ を分割冪環とし、$A$ は
$\mathbf{Z}_{(p)}$ 上の代数で、$p$ は $A/I$ で冪零であるとする。
$S = \Spec(A)$ および $S_0 = \Spec(A/I)$ とおく。$X$ を $S_0$ 上の
スキームとし、$X$ 上で $p$ は局所的に冪零であるとする。
$\mathcal{F}$ を任意の $\mathcal{O}_{X/S}$ 加群とする。$e \gg 0$ に対して
$(p^e) \subset I$ は $\gamma$ によって保たれる。『分割冪代数』補題
\ref{dpa-lemma-extend-to-completion} を参照されたい。
$e \gg 0$ に対して $S_e = \Spec(A/p^eA)$ とおく。このとき
$\text{Cris}(X/S_e)$ は $\text{Cris}(X/S)$ の充満部分圏であり、
$\mathcal{F}$ の $\text{Cris}(X/S_e)$ への制限を $\mathcal{F}_e$ と書く。
このとき
$$
R\Gamma(\text{Cris}(X/S), \mathcal{F}) =
R\lim_e R\Gamma(\text{Cris}(X/S_e), \mathcal{F}_e)
$$

\medskip\noindent
ヒント：$\mathcal{F}$ が単射的である場合に示せば十分である。この場合、層
$\mathcal{F}_e$ も単射的加群であり、遷移写像
$\Gamma(\mathcal{F}_{e + 1}) \to \Gamma(\mathcal{F}_e)$ は全射である。
また $\text{Cris}(X/S)$ の任意の対象は、$\text{Cris}(X/S)$ の定義により
局所的には圏 $\text{Cris}(X/S_e)$ のいずれか一つの対象なので、
$\Gamma(\mathcal{F}) = \lim_e \Gamma(\mathcal{F}_e)$ である。
\end{remark}

\begin{remark}[比較]
\label{crystalline-remark-comparison}
$p$ を素数とする。$(A, I, \gamma)$ を分割冪環とし、$p$ は $A$ で
冪零であるとする。$S = \Spec(A)$ および $S_0 = \Spec(A/I)$ とおく。
$Y$ を $S$ 上の滑らかなスキームとし、$X = Y \times_S S_0$ とおく。
$\mathcal{F}$ を準連接 $\mathcal{O}_{X/S}$ 加群のクリスタルとする。
このとき次が成り立つ。
\begin{enumerate}
\item $\gamma$ は $Y$ における $X$ のイデアル上の分割冪構造へ延長され、
$(X, Y, \gamma)$ は $\text{Cris}(X/S)$ の対象となる。
\item 制限 $\mathcal{F}_Y$（節 \ref{crystalline-section-sheaves} を参照）は、標準的な
可積分接続
$\nabla : \mathcal{F}_Y \to
\mathcal{F}_Y \otimes_{\mathcal{O}_Y} \Omega_{Y/S}$ を備える。
\item $D(A)$ において
$$
R\Gamma(\text{Cris}(X/S), \mathcal{F}) =
R\Gamma(Y, \mathcal{F}_Y \otimes_{\mathcal{O}_Y} \Omega^\bullet_{Y/S})
$$
である。
\end{enumerate}
ヒント：(1) については『分割冪代数』補題
\ref{dpa-lemma-gamma-extends} を参照されたい。(2) については補題
\ref{crystalline-lemma-automatic-connection} を参照されたい。(3) については、
(\ref{crystalline-equation-restriction}) の写像があることに注意する。この写像は
$X$ がアフィンなら同型である。補題
\ref{crystalline-lemma-compute-cohomology-crystal-smooth} を参照されたい。
これにより $Ru_{X/S, *}\mathcal{F}$ と
$\mathcal{F}_Y \otimes \Omega^\bullet_{Y/S}$ は
$Y_{Zar} = X_{Zar}$ 上の複体として擬同型である。
$R\Gamma(\text{Cris}(X/S), \mathcal{F}) =
R\Gamma(X_{Zar}, Ru_{X/S, *}\mathcal{F})$ なので、結果が従う。
\end{remark}

\begin{remark}[完全性]
\label{crystalline-remark-perfect}
$p$ を素数とする。$(A, I, \gamma)$ を分割冪環とし、$p$ は $A$ で
冪零であるとする。$S = \Spec(A)$ および $S_0 = \Spec(A/I)$ とおく。
$X$ を $S_0$ 上の固有かつ滑らかなスキームとする。$\mathcal{F}$ を
有限局所自由な準連接 $\mathcal{O}_{X/S}$ 加群のクリスタルとする。
このとき $R\Gamma(\text{Cris}(X/S), \mathcal{F})$ は
$D(A)$ の完全対象である。

\medskip\noindent
ヒント：注意 \ref{crystalline-remark-base-change-isomorphism} により
$$
R\Gamma(\text{Cris}(X/S), \mathcal{F}) \otimes_A^\mathbf{L} A/I
\cong
R\Gamma(\text{Cris}(X/S_0), \mathcal{F}|_{\text{Cris}(X/S_0)})
$$
である。注意 \ref{crystalline-remark-comparison} により
$$
R\Gamma(\text{Cris}(X/S_0), \mathcal{F}|_{\text{Cris}(X/S_0)}) =
R\Gamma(X, \mathcal{F}_X \otimes \Omega^\bullet_{X/S_0})
$$
である。de Rham 複体上の愚直フィルトレーションを用いると、複体
$$
R\Gamma(X, \mathcal{F}_X \otimes \Omega^q_{X/S_0})
$$
が $D(A/I)$ の完全複体なら、直前に表示した複体も $D(A/I)$ で
完全であることが分かる。『代数詳論』補題
\ref{more-algebra-lemma-two-out-of-three-perfect} を参照されたい。
$\mathcal{F}_X \otimes \Omega^q_{X/S_0}$ が有限局所自由層であり、
$X \to S_0$ が固有かつ平坦であることを用いる連接コホモロジーの
標準的な議論により、これは成り立つ（将来の参照をここに挿入）。
『代数詳論』補題
\ref{more-algebra-lemma-perfect-modulo-nilpotent-ideal} を適用すると、
$$
R\Gamma(\text{Cris}(X/S), \mathcal{F}) \otimes_A^\mathbf{L} A/I^n
$$
はすべての $n$ に対して $D(A/I^n)$ の完全対象である。しかし
$A$ が Noether 環でなければ、これだけでは十分でない。実際、$p$ が
冪零であるという仮定により $I$ は局所的に冪零であるにもかかわらず、
ある $n$ に対して $I^n = 0$ と結論することはできない。
『分割冪代数』補題 \ref{dpa-lemma-nil} を参照されたい。
反例は $\mathbf{F}_p\langle x \rangle$ である。
$\mathcal{F} = \mathcal{O}_{X/S}$ の場合に一般にこれを証明するには、
\url{https://math.columbia.edu/~dejong/wordpress/?p=2227} の議論が使える。
係数 $\mathcal{F}$ が自明でない場合、\cite{Faltings-very} の議論は
次のようであると思われる。『代数詳論』補題
\ref{more-algebra-lemma-perfect-modulo-nilpotent-ideal} により
$pA = 0$ の場合に帰着する。この場合 Frobenius 写像
$A \to A$, $a \mapsto a^p$ は $A \to A/I \xrightarrow{\varphi} A$ と
分解する。実際 $x \in I$ に対して $x^p = 0$ である。
$X^{(1)} = X \otimes_{A/I, \varphi} A$ とおく。$X$ の絶対 Frobenius 射は
射 $F_X : X \to X^{(1)}$ を経由して分解する。これは一種の相対
Frobenius である。アフィン局所的に $X = \Spec(C)$ なら、
$X^{(1)} = \Spec( C \otimes_{A/I, \varphi} A)$ であり、$F_X$ は
$C \otimes_{A/I, \varphi} A \to C$, $c \otimes a \mapsto c^pa$ に
対応する。これにより環付きトポスの射
$$
(X/S)_{\text{cris}}
\xrightarrow{(F_X)_{\text{cris}}}
(X^{(1)}/S)_{\text{cris}}
\xrightarrow{u_{X^{(1)}/S}}
\Sh(X^{(1)}_{Zar})
$$
が定まる。その合成を $\text{Frob}_X$ と書く。次に局所計算によって、
$R\text{Frob}_{X, *}\mathcal{F}$ が $\mathcal{O}_{X^{(1)}}$ 加群の
完全複体によって表されることを示す。
\end{remark}

\begin{remark}[完備な場合の完全性]
\label{crystalline-remark-complete-perfect}
$p$ を素数とする。$(A, I, \gamma)$ を分割冪環とし、$A$ は
$p$ 進完備環で、$p$ は $A/I$ で冪零であるとする。
$S = \Spec(A)$ および $S_0 = \Spec(A/I)$ とおく。$X$ を $S_0$ 上の
固有かつ滑らかなスキームとする。$\mathcal{F}$ を有限局所自由な
準連接 $\mathcal{O}_{X/S}$ 加群のクリスタルとする。このとき
$R\Gamma(\text{Cris}(X/S), \mathcal{F})$ は $D(A)$ の完全対象である。

\medskip\noindent
ヒント：注意 \ref{crystalline-remark-rlim} により、
$K = R\Gamma(\text{Cris}(X/S), \mathcal{F})$ は $A/p^eA$ 上の
コホモロジーの導来極限 $K = R\lim K_e$ である。注意
\ref{crystalline-remark-perfect} により、各 $K_e$ は $D(A/p^eA)$ の完全複体である。
$A$ は $p$ 進完備なので、『代数詳論』補題
\ref{more-algebra-lemma-Rlim-perfect-gives-complete} から結果が従う。
\end{remark}

\begin{remark}[完備な場合の比較]
\label{crystalline-remark-complete-comparison}
$p$ を素数とする。$(A, I, \gamma)$ を分割冪環とし、$A$ は Noether な
$p$ 進完備環で、$p$ は $A/I$ で冪零であるとする。
$S = \Spec(A)$ および $S_0 = \Spec(A/I)$ とおく。$Y$ を $S$ 上の
固有かつ滑らかなスキームとし、$X = Y \times_S S_0$ とおく。
$\mathcal{F}$ を準連接 $\mathcal{O}_{X/S}$ 加群の有限型クリスタルとする。
このとき次が成り立つ。
\begin{enumerate}
\item 可積分接続
$$
\nabla :
\mathcal{F}_Y
\longrightarrow
\mathcal{F}_Y \otimes_{\mathcal{O}_Y} \Omega_{Y/S}
$$
を備えた連接 $\mathcal{O}_Y$ 加群 $\mathcal{F}_Y$ が存在し、
$\mathcal{F}_Y/p^e\mathcal{F}_Y$ は注意 \ref{crystalline-remark-comparison} で得られた
$A/p^eA$ 上の接続付き加群である。
\item $D(A)$ において
$$
R\Gamma(\text{Cris}(X/S), \mathcal{F}) =
R\Gamma(Y, \mathcal{F}_Y \otimes_{\mathcal{O}_Y} \Omega^\bullet_{Y/S})
$$
である。
\end{enumerate}
ヒント：$\mathcal{F}_Y$ の存在は Grothendieck の存在定理である
（将来の参照をここに挿入）。両辺は法 $p^e$ の各版の $R\lim$ として
計算されるので、コホモロジーの同型が従う。左辺については注意
\ref{crystalline-remark-rlim} を参照し、右辺については形式関数定理、すなわち
『スキームのコホモロジー』定理
\ref{coherent-theorem-formal-functions} を用いる。法 $p^e$ の各版は
注意 \ref{crystalline-remark-comparison} により同型である。
\end{remark}




\section{純非分離写像に沿う引き戻し}
\label{crystalline-section-pull-back-along-pth-root}

\noindent
$\alpha_p$ 被覆とは、形
$$
X' = \Spec(C[z]/(z^p - c)) \longrightarrow \Spec(C) = X
$$
の射をいう。ここで $C$ は $\mathbf{F}_p$ 代数で、$c \in C$ である。
同値な言い方をすれば、$X'$ は $X$ 上の $\alpha_p$ トーサーである。
{\it 反復 $\alpha_p$ 被覆}\footnote{これは標準的な用語ではない。}とは、
標数 $p$ のスキームの射であって、標的上局所的に有限個の
$\alpha_p$ 被覆の合成となるものをいう。この節では、そのような射に沿う
引き戻しが、素数 $p$ を可逆にした後で結晶コホモロジー上の擬同型を
誘導することを証明する。実際には、この結果の精密な形を証明する。
まず、定式化にいくつかの記法を要する準備補題から始める。

\medskip\noindent
注意 \ref{crystalline-remark-base-change-connection} の仮定を満たす環準同型
$B \to B'$ と商 $\Omega_B \to \Omega$ および
$\Omega_{B'} \to \Omega'$ があると仮定する。したがって
(\ref{crystalline-equation-base-change-map-complexes}) は、可積分接続
$\nabla : M \to M \otimes_B \Omega_B$ を備えたすべての $B$ 加群
$M$ に対して標準的な複体の写像
$$
c_M^\bullet :
M \otimes_B \Omega^\bullet
\longrightarrow
M \otimes_B (\Omega')^\bullet
$$
を与える。

\medskip\noindent
$a \in B$, $z \in B'$、および次の仮定を満たす写像
$\theta : B' \to B'$ があるとする。
\begin{enumerate}
\item
\label{crystalline-item-d-a-zero}
$\text{d}(a) = 0$,
\item
\label{crystalline-item-direct-sum}
$\Omega' = B' \otimes_B \Omega \oplus B'\text{d}z$ である。すべての
$f \in B'$ に対して
$\text{d}(f) = \text{d}_1(f) + \partial_z(f) \text{d}z$
と書く。ここで $\text{d}_1(f) \in B' \otimes \Omega$ および
$\partial_z(f) \in B'$ である。
\item
\label{crystalline-item-theta-linear}
$\theta : B' \to B'$ は $B$ 線形である。
\item
\label{crystalline-item-integrate}
$\partial_z \circ \theta = a$,
\item
\label{crystalline-item-injective}
$B \to B'$ は普遍的単射である。したがって $\Omega \to \Omega'$ は単射である。
\item
\label{crystalline-item-factor}
すべての $f \in B'$ に対して $af - \theta(\partial_z(f)) \in B$ である。
\item
\label{crystalline-item-horizontal}
$(\theta \otimes 1)(\text{d}_1(f)) - \text{d}_1(\theta(f)) \in \Omega$
がすべての $f \in B'$ に対して成り立つ。ここで
$\theta \otimes 1 : B' \otimes \Omega \to B' \otimes \Omega$
\end{enumerate}
これらの条件は論理的に独立ではない。例えば仮定
(\ref{crystalline-item-integrate}) から
$\partial_z(af - \theta(\partial_z(f))) = 0$ が従う。したがって
$B \to B'$ の像が $\partial_z$ によって消される元全体なら、
(\ref{crystalline-item-factor}) が従う。条件 (\ref{crystalline-item-horizontal}) についても
同様の議論ができる。

\begin{lemma}
\label{crystalline-lemma-find-homotopy}
上の状況で複体の写像
$$
e_M^\bullet :
M \otimes_B (\Omega')^\bullet
\longrightarrow
M \otimes_B \Omega^\bullet
$$
が存在し、$c_M^\bullet \circ e_M^\bullet$ と
$e_M^\bullet \circ c_M^\bullet$ は $a$ 倍写像とホモトピックである。
\end{lemma}

\begin{proof}
この証明ではすべてのテンソル積を $B$ 上でとる。仮定
(\ref{crystalline-item-direct-sum}) により、すべての $i \geq 0$ に対して
$$
M \otimes (\Omega')^i =
(B' \otimes M \otimes \Omega^i)
\oplus
(B' \text{d}z \otimes M \otimes \Omega^{i - 1})
$$
である。$M \otimes (\Omega')^i$ の加法的生成元の一組は、形
$f \omega$ の元と形 $f \text{d}z \wedge \eta$ の元からなる。
ここで $f \in B'$, $\omega \in M \otimes \Omega^i$, および
$\eta \in M \otimes \Omega^{i - 1}$ である。

\medskip\noindent
$f \in B'$ に対して
$$
\epsilon(f) = af - \theta(\partial_z(f))
\quad\text{および}\quad
\epsilon'(f) = (\theta \otimes 1)(\text{d}_1(f)) - \text{d}_1(\theta(f))
$$
と書く。仮定 (\ref{crystalline-item-factor}) および (\ref{crystalline-item-horizontal}) により、
$\epsilon(f) \in B$ および $\epsilon'(f) \in \Omega$ である。
規則 $e^i_M(f\omega) = \epsilon(f) \omega$ および
$e^i_M(f \text{d}z \wedge \eta) = \epsilon'(f) \wedge \eta$ によって
$e_M^\bullet$ を定義する。写像族 $e^i_M$ が複体の写像であることは
以下で分かる。

\medskip\noindent
規則 $h^i(f \omega) = 0$ および
$h^i(f \text{d}z \wedge \eta) = \theta(f) \eta$ により
$$
h^i : M \otimes_B (\Omega')^i \longrightarrow M \otimes_B (\Omega')^{i - 1}
$$
を定義する。ここで元は上のとおりである。次を主張する。
$$
\text{d} \circ h + h \circ \text{d} = a - c_M^\bullet \circ e_M^\bullet
$$
(\ref{crystalline-item-d-a-zero}) により $a$ 倍写像は複体の写像であることに注意する。
したがって仮定 (\ref{crystalline-item-injective}) により $c_M^\bullet$ は単射的な
複体の写像なので、$e_M^\bullet$ は複体の写像であると結論できる。
主張を証明するために計算すると、
\begin{align*}
(\text{d} \circ h + h \circ \text{d})(f \omega)
& =
h\left(\text{d}(f) \wedge \omega + f \nabla(\omega)\right)
\\
& =
\theta(\partial_z(f)) \omega
\\
& =
a f\omega - \epsilon(f)\omega 
\\
& =
a f \omega - c^i_M(e^i_M(f\omega))
\end{align*}
である。第二の等号は $\nabla(\omega)$ に $\text{d}z$ が現れないことにより、
第三の等号は仮定 (6) による。同様に
\begin{align*}
(\text{d} \circ h + h \circ \text{d})(f \text{d}z \wedge \eta)
& =
\text{d}(\theta(f) \eta) +
h\left(\text{d}(f) \wedge \text{d}z \wedge \eta -
f \text{d}z \wedge \nabla(\eta)\right)
\\
& =
\text{d}(\theta(f)) \wedge \eta + \theta(f) \nabla(\eta)
- (\theta \otimes 1)(\text{d}_1(f)) \wedge \eta
- \theta(f) \nabla(\eta)
\\
& =
\text{d}_1(\theta(f)) \wedge \eta +
\partial_z(\theta(f)) \text{d}z \wedge \eta -
(\theta \otimes 1)(\text{d}_1(f)) \wedge \eta
\\
& =
a f \text{d}z \wedge \eta - \epsilon'(f) \wedge \eta \\
& = a f \text{d}z \wedge \eta - c^i_M(e^i_M(f \text{d}z \wedge \eta))
\end{align*}
である。第二の等号は
$\text{d}(f) \wedge \text{d}z \wedge \eta =
- \text{d}z \wedge \text{d}_1(f) \wedge \eta$ により、第四の等号は仮定
(\ref{crystalline-item-integrate}) による。他方、定義から直ちに
$e^i_M(c^i_M(\omega)) = \epsilon(1) \omega = a \omega$ である。
これで補題が証明された。
\end{proof}

\begin{example}
\label{crystalline-example-integrate}
上の状況の標準的な例は、$B' = B\langle z \rangle$ が分割冪環
$(B, J, \delta)$ 上の分割冪多項式環で、$J' = B'_{+} + JB' \subset B'$
上に分割冪 $\delta'$ をもつ場合に得られる。すなわち、
$\Omega = \Omega_{B, \delta}$ および $\Omega' = \Omega_{B', \delta'}$
とする。この場合 $a = 1$ ととることができ、
$$
\theta( \sum b_m z^{[m]} ) = \sum b_m z^{[m + 1]}
$$
とする。ここで
$$
f - \theta(\partial_z(f)) = f(0)
$$
は定数項に等しいことに注意する。したがってこの場合、補題
\ref{crystalline-lemma-find-homotopy} は結晶 Poincar\'e 補題、すなわち補題
\ref{crystalline-lemma-relative-poincare} を再現する。
\end{example}

\begin{lemma}
\label{crystalline-lemma-computation}
状況 \ref{crystalline-situation-affine} において、$D$ と $\Omega_D$ は
(\ref{crystalline-equation-D}) および (\ref{crystalline-equation-omega-D}) のとおりであるとする。
$\lambda \in D$ とする。$D'$ を
$$
D[z]\langle \xi \rangle/(\xi - (z^p - \lambda))
$$
の $p$ 進完備化とし、$\Omega_{D'}$ を $D'$ の $A$ 上の
分割冪微分加群の $p$ 進完備化とする。
(\ref{crystalline-item-complete}), (\ref{crystalline-item-connection}), (\ref{crystalline-item-integrable}),
および (\ref{crystalline-item-topologically-quasi-nilpotent}) を満たす $D$ 上の
任意の対 $(M, \nabla)$ に対して、標準的な複体の写像
(\ref{crystalline-equation-base-change-map-complexes})
$$
c_M^\bullet : M \otimes_D^\wedge \Omega^\bullet_D
\longrightarrow
M \otimes_D^\wedge \Omega^\bullet_{D'}
$$
は次の性質をもつ。逆向きの写像 $e_M^\bullet$ が存在し、
$c_M^\bullet \circ e_M^\bullet$ と $e_M^\bullet \circ c_M^\bullet$ は
ともに $p$ 倍写像とホモトピックである。
\end{lemma}

\begin{proof}
補題 \ref{crystalline-lemma-find-homotopy} を $a = p$ として用いて証明する。
したがって $\theta : D' \to D'$ を構成し、
(\ref{crystalline-item-d-a-zero}), (\ref{crystalline-item-direct-sum}), (\ref{crystalline-item-theta-linear}),
(\ref{crystalline-item-integrate}), (\ref{crystalline-item-injective}), (\ref{crystalline-item-factor}),
(\ref{crystalline-item-horizontal}) を証明する必要がある。まず環 $D$, $D'$ と
加群 $\Omega_D$, $\Omega_{D'}$ に関する情報をまとめる。

\medskip\noindent
等式
$$
D[z]\langle \xi \rangle/(\xi - (z^p - \lambda))
=
D\langle \xi \rangle[z]/(z^p - \xi - \lambda)
$$
から、$D'$ は自由 $D$ 加群
$$
\bigoplus\nolimits_{i = 0, \ldots, p - 1}
\bigoplus\nolimits_{n \geq 0}
z^i \xi^{[n]} D
$$
の $p$ 進完備化であることが分かる。ここで $\xi^{[0]} = 1$ である。
したがって $D \to D'$ は連続な $D$ 線形切断をもち、特に
$D \to D'$ は普遍的単射である。すなわち (\ref{crystalline-item-injective}) が成り立つ。
$D'$ を、分割冪イデアル $\overline{J}' = \overline{J}D' + (\xi)$ をもつ
$A$ 上の分割冪代数とみなす。このとき $D'$ はまた、$D$ において
$z^p - \lambda$ が生成するイデアルの分割冪包の $p$ 進完備化である。
補題 \ref{crystalline-lemma-describe-divided-power-envelope} を参照されたい。したがって
$$
\Omega_{D'} = \Omega_D \otimes_D^\wedge D' \oplus D'\text{d}z
$$
である。補題 \ref{crystalline-lemma-module-differentials-divided-power-envelope} を
参照されたい。これで (\ref{crystalline-item-direct-sum}) が証明された。
(\ref{crystalline-item-d-a-zero}) は明らかである。

\medskip\noindent
ここで $\theta$ を構成する。（この証明のいくつかの公式を検証する
PARI/gp スクリプト theta.gp を作成した。これは Stacks project の
scripts サブディレクトリにある。）その前に、元 $z^i \xi^{[n]}$ の
微分を計算する。$\text{d}z^i = i z^{i - 1} \text{d}z$ である。
$n \geq 1$ に対して
$$
\text{d}\xi^{[n]} =
\xi^{[n - 1]} \text{d}\xi =
- \xi^{[n - 1]}\text{d}\lambda + p z^{p - 1} \xi^{[n - 1]}\text{d}z
$$
である。これは $\xi = z^p - \lambda$ による。
$0 < i < p$ および $n \geq 1$ に対して
\begin{align*}
\text{d}(z^i\xi^{[n]})
& =
iz^{i - 1}\xi^{[n]}\text{d}z + z^i\xi^{[n - 1]}\text{d}\xi \\
& =
iz^{i - 1}\xi^{[n]}\text{d}z + z^i\xi^{[n - 1]}\text{d}(z^p - \lambda) \\
& =
- z^i\xi^{[n - 1]}\text{d}\lambda +
(iz^{i - 1}\xi^{[n]} + pz^{i + p - 1}\xi^{[n - 1]})\text{d}z \\
& =
- z^i\xi^{[n - 1]}\text{d}\lambda +
(iz^{i - 1}\xi^{[n]} + pz^{i - 1}(\xi + \lambda)\xi^{[n - 1]})\text{d}z \\
& =
- z^i\xi^{[n - 1]}\text{d}\lambda +
((i + pn)z^{i - 1}\xi^{[n]} + p\lambda z^{i - 1}\xi^{[n - 1]})\text{d}z
\end{align*}
である。最後の等号は $\xi \xi^{[n - 1]} = n\xi^{[n]}$ による。
したがって
\begin{align*}
\partial_z(z^i) & = i z^{i - 1} \\
\partial_z(\xi^{[n]}) & = p z^{p - 1} \xi^{[n - 1]} \\
\partial_z(z^i\xi^{[n]}) & =
(i + pn) z^{i - 1} \xi^{[n]} + p \lambda z^{i - 1}\xi^{[n - 1]}
\end{align*}
である。これらの公式に従い、規則
$$
\begin{matrix}
\theta(z^j)
& = & p\frac{z^{j + 1}}{j + 1}
& j = 0, \ldots p - 1, \\
\theta(z^{p - 1}\xi^{[m]})
& = & \xi^{[m + 1]}
& m \geq 1, \\
\theta(z^j \xi^{[m]})
& = &
\frac{p z^{j + 1} \xi^{[m]} - \theta(p\lambda z^j \xi^{[m - 1]})}{(j + 1 + pm)}
& 0 \leq j < p - 1, m \geq 1
\end{matrix}
$$
により $\theta$ を定義する。最後の行では $m$ に関する帰納法を用いて
$\theta$ の値を定める。これを展開すると、$0 \leq j < p - 1$ および
$1 \leq m$ に対して
$$
\theta(z^j \xi^{[m]}) =
\textstyle{\frac{p z^{j + 1} \xi^{[m]}}{(j + 1 + pm)} -
\frac{p^2 \lambda z^{j + 1} \xi^{[m - 1]}}{(j + 1 + pm)(j + 1 + p(m - 1))} +
\ldots +
\frac{(-1)^m p^{m + 1} \lambda^m z^{j + 1}}
{(j + 1 + pm) \ldots (j + 1)}}
$$
を得るが、この式は以下では使わない。$\theta$ が $D'$ 上の
$p$ 進連続な $D$ 線形写像へ一意に延長されることは明らかである。
構成により (\ref{crystalline-item-theta-linear}) と (\ref{crystalline-item-integrate}) が成り立つ。
残るのは (\ref{crystalline-item-factor}) と (\ref{crystalline-item-horizontal}) の証明である。

\medskip\noindent
(\ref{crystalline-item-factor}) と (\ref{crystalline-item-horizontal}) の証明。
$\theta$ は $D$ 線形かつ連続なので、
$p - \theta \circ \partial_z$,
および $(\theta \otimes 1) \circ \text{d}_1 - \text{d}_1 \circ \theta$ を
元 $z^i\xi^{[n]}$ に作用させた値が、それぞれ $D$ および $\Omega_D$ の
元であることを示せば十分である\footnote{これは直接計算できる。
$p - \theta \circ \partial_z$ を $z^i\xi^{[n]}$ に作用させると、
$p$ に写る $1$ と $-p\lambda$ に写る $\xi$ を除いて零になる。
$(\theta \otimes 1) \circ \text{d}_1 - \text{d}_1 \circ \theta$ を
$z^i\xi^{[n]}$ に作用させると、$-\lambda$ に写る
$z^{p - 1}\xi$ を除いて零になる。}。
$D_0 = \mathbf{Z}_{(p)}[\lambda]$ および
$D_0' = \mathbf{Z}_{(p)}[z, \lambda]\langle \xi \rangle/(\xi - z^p + \lambda)$
とおく。上の各式は $D_0'$ または $\Omega_{D_0'}$ の元であることに
注意する。したがって $D_0 \to D_0'$ の場合に結果を示せば十分である。
$D_0$ と $D_0'$ はねじれのない環であり、
$D_0 \otimes \mathbf{Q} = \mathbf{Q}[\lambda]$ および
$D'_0 \otimes \mathbf{Q} = \mathbf{Q}[z, \lambda]$ であることに注意する。
ゆえに $D_0 \subset D'_0$ は $\partial_z$ によって消される元からなる
部分環であり、(\ref{crystalline-item-factor}) は (\ref{crystalline-item-integrate}) から従う。
補題 \ref{crystalline-lemma-find-homotopy} の直前の議論を参照されたい。同様に
$\text{d}_1(f) = \partial_\lambda(f)\text{d}\lambda$ なので、
$$
\left((\theta \otimes 1) \circ \text{d}_1 - \text{d}_1 \circ \theta\right)(f)
=
\left(\theta(\partial_\lambda(f)) - \partial_\lambda(\theta(f))\right)
\text{d}\lambda
$$
である。係数に $\partial_z$ を適用すると
\begin{align*}
\partial_z\left(
\theta(\partial_\lambda(f)) - \partial_\lambda(\theta(f))
\right)
& =
p \partial_\lambda(f) - \partial_z(\partial_\lambda(\theta(f))) \\
& =
p \partial_\lambda(f) - \partial_\lambda(\partial_z(\theta(f))) \\
& =
p \partial_\lambda(f) - \partial_\lambda(p f) = 0
\end{align*}
を得る。したがって望むとおり係数は $z$ に依存しない。
これで補題の証明が完了する。
\end{proof}

\noindent
反復 $\alpha_p$ 被覆 $X' \to X$ は、この節の導入で定義したとおり
有限局所自由であることに注意する。したがって $X$ が連結なら、
$X' \to X$ の次数は一定であり、$p$ の冪である。

\begin{lemma}
\label{crystalline-lemma-pullback-along-p-power-cover}
$p$ を素数とする。$(S, \mathcal{I}, \gamma)$ を
$\mathbf{Z}_{(p)}$ 上の分割冪スキームとし、$p \in \mathcal{I}$ とする。
$S_0 = V(\mathcal{I}) \subset S$ とおく。$f : X' \to X$ を、次数が
一定値 $q$ である $S_0$ 上のスキームの反復 $\alpha_p$ 被覆とする。
$\mathcal{F}$ を $X$ 上の準連接層の任意のクリスタルとし、
$\mathcal{F}' = f_{\text{cris}}^*\mathcal{F}$ とおく。識別三角
$$
Ru_{X/S, *}\mathcal{F}
\longrightarrow
f_*Ru_{X'/S, *}\mathcal{F}'
\longrightarrow
E
\longrightarrow
Ru_{X/S, *}\mathcal{F}[1]
$$
において、対象 $E$ のコホモロジー層は $q$ によって消される。
\end{lemma}

\begin{proof}
$X' \to X$ は同相写像なので、$X$ と $X'$ の台位相空間を同一視できる
ことに注意する。問題は明らかに $X$ 上局所的である。したがって
$X$, $X'$, $S$ はアフィンであり、$X' \to X$ は合成
$$
X' = X_n \to X_{n - 1} \to X_{n - 2} \to \ldots \to X_0 = X
$$
で与えられると仮定してよい。ここで各射 $X_{i + 1} \to X_i$ は
$\alpha_p$ 被覆である。$\mathcal{F}$ の $X_i$ への引き戻しを
$\mathcal{F}_i$ と書く。各写像
$$
R\Gamma(\text{Cris}(X_i/S), \mathcal{F}_i)
\longrightarrow
R\Gamma(\text{Cris}(X_{i + 1}/S), \mathcal{F}_{i + 1})
$$
が、第三項のコホモロジー群が $p$ によって消される三角に入ることを
示せば十分である。（ここでは三角圏 $D(X)$ の公理 TR4 を用いる。
詳細は省略する。）

\medskip\noindent
したがって $S = \Spec(A)$, $X = \Spec(C)$, $X' = \Spec(C')$ であり、
ある $c \in C$ に対して $C' = C[z]/(z^p - c)$ であると仮定してよい。
$A$ 上の多項式代数 $P$ と全射 $P \to C$ を選ぶ。
$D$ を、(\ref{crystalline-equation-D}) のように $P$ における
$\Ker(P \to C)$ の分割冪包の $p$ 進完備化とする。
$P' = P[z]$ とおき、$P' \to C'$ を $z$ が $C'$ における $z$ の類へ
写る全射とする。$c \in C$ の持ち上げ $\lambda \in D$ を選ぶ。
このとき、$P'$ における $\Ker(P' \to C')$ の分割冪包の
$p$ 進完備化 $D'$ は
$D[z]\langle \xi \rangle/(\xi - (z^p - \lambda))$ の $p$ 進完備化と
同型である。補題 \ref{crystalline-lemma-computation} とその証明を参照されたい。
したがって命題 \ref{crystalline-proposition-compute-cohomology-crystal} による
準連接加群のクリスタルのコホモロジー計算とこの補題から、結果が従う。
\end{proof}

\noindent
次の補題の評価はおそらく最良ではない。

\begin{lemma}
\label{crystalline-lemma-pullback-along-p-power-cover-cohomology}
記法と仮定を補題 \ref{crystalline-lemma-pullback-along-p-power-cover} と同じものとする。
写像
$$
f^* :
H^i(\text{Cris}(X/S), \mathcal{F})
\longrightarrow
H^i(\text{Cris}(X'/S), \mathcal{F}')
$$
の核と余核は $q^{i + 1}$ によって消される。
\end{lemma}

\begin{proof}
$E$ の零でないコホモロジー層は次数 $-1$ 以上にあり、したがって
スペクトル系列
$H^a(\mathcal{H}^b(E)) \Rightarrow H^{a + b}(E)$ は収束する。
これと識別三角に付随するコホモロジー長完全列を合わせれば、
この評価が得られる。
\end{proof}

\noindent
状況 \ref{crystalline-situation-global} において $p \in \mathcal{I}$ と仮定する。
$$ 
X^{(1)} = X \times_{S_0, F_{S_0}} S_0.
$$
とおく。相対 Frobenius 射を $F_{X/S_0} : X \to X^{(1)}$ と書く。

\begin{lemma}
\label{crystalline-lemma-pullback-relative-frobenius}
上の状況で、$X \to S_0$ は相対次元 $d$ の滑らかな射であると仮定する。
このとき $F_{X/S_0}$ は次数 $p^d$ の反復 $\alpha_p$ 被覆である。
したがって補題 \ref{crystalline-lemma-pullback-along-p-power-cover} および
\ref{crystalline-lemma-pullback-along-p-power-cover-cohomology} をこの状況に適用できる。
特に、$\text{Cris}(X^{(1)}/S)$ 上の準連接加群の任意のクリスタル
$\mathcal{G}$ に対して、写像
$$
F_{X/S_0}^* : H^i(\text{Cris}(X^{(1)}/S), \mathcal{G})
\longrightarrow
H^i(\text{Cris}(X/S), F_{X/S_0, \text{cris}}^*\mathcal{G})
$$
の核と余核は $p^{d(i + 1)}$ によって消される。
\end{lemma}

\begin{proof}
第一の主張を証明すれば十分である。そのために、$X$ は
$\mathbf{A}^d_{S_0}$ 上 \'etale であると仮定してよい。
『射』補題 \ref{morphisms-lemma-smooth-etale-over-affine-space} を
参照されたい。この \'etale 射を
$\varphi : X \to \mathbf{A}^d_{S_0}$ と書く。この場合、
$X/S_0$ の相対 Frobenius は図式
$$
\xymatrix{
X \ar[d] \ar[r] & X^{(1)} \ar[d] \\
\mathbf{A}^d_{S_0} \ar[r] & \mathbf{A}^d_{S_0}
}
$$
に入る。ここで下の水平射は
$\mathbf{A}^d_{S_0}$ の $S_0$ 上の相対 Frobenius 射である。
これはすべての座標を $p$ 乗する射なので、反復 $\alpha_p$ 被覆である。
この図式が Cartesian 正方形であることに注意すれば証明は完了する。
『\'Etale 射』補題 \ref{etale-lemma-relative-frobenius-etale} を
参照されたい。
\end{proof}













\section{結晶コホモロジー上の Frobenius 作用}
\label{crystalline-section-frobenius}

\noindent
この節では、Frobenius 引き戻しが、素数 $p$ を可逆にした後で
結晶コホモロジー上の擬同型を誘導することを証明する。ただし、この主張を
定式化するためにも、特別な状況で議論する必要がある。

\begin{situation}
\label{crystalline-situation-F-crystal}
状況 \ref{crystalline-situation-global} において、次を仮定する。
\begin{enumerate}
\item $p \in I$ を満たすある分割冪環 $(A, I, \gamma)$ に対して
$S = \Spec(A)$ である。
\item すべての $x \in A$ に対して $\sigma(x) = x^p \bmod pA$ を満たす
分割冪環の準同型 $\sigma : A \to A$ が与えられている。
\end{enumerate}
\end{situation}

\noindent
状況 \ref{crystalline-situation-F-crystal} では、射
$\Spec(\sigma) : S \to S$ は絶対 Frobenius
$F_{S_0} : S_0 \to S_0$ の持ち上げである。また図式
$$
\xymatrix{
X \ar[d] \ar[r]_{F_X} & X \ar[d] \\
S_0 \ar[r]^{F_{S_0}} & S_0
}
$$
は可換である。ここで $F_X : X \to X$ は $X$ の絶対 Frobenius 射である。
したがって結晶トポスの射
$$
(F_X)_{\text{cris}} :
(X/S)_{\text{cris}}
\longrightarrow
(X/S)_{\text{cris}}
$$
を得る。注意 \ref{crystalline-remark-functoriality-cris} を参照されたい。
以下では Saavedra の記法に従って $F$-クリスタルに関する用語を定める。
\cite{Saavedra} を参照されたい。

\begin{definition}
\label{crystalline-definition-F-crystal}
状況 \ref{crystalline-situation-F-crystal} における
{\it $X/S$ 上の（$\sigma$ に関する）$F$-クリスタル}とは、
有限局所自由な $\mathcal{O}_{X/S}$ 加群のクリスタル
$\mathcal{E}$ と写像
$$
F_\mathcal{E} : (F_X)_{\text{cris}}^*\mathcal{E} \longrightarrow \mathcal{E}
$$
からなる対 $(\mathcal{E}, F_\mathcal{E})$ をいう。整数 $i \geq 0$ と写像
$V : \mathcal{E} \to (F_X)_{\text{cris}}^*\mathcal{E}$ が存在して
$V \circ F_{\mathcal{E}} = p^i \text{id}$ となるとき、
$F$-クリスタルは{\it 非退化}であるという。
\end{definition}

\begin{remark}
\label{crystalline-remark-F-crystal-variants}
$(\mathcal{E}, F)$ を定義 \ref{crystalline-definition-F-crystal} のような
$F$-クリスタルとする。文献では、非退化条件を $F$-クリスタルの定義に
含めることが多い。さらに $F \circ V = p^n\text{id}$ も仮定することが
多い。以下の結果に必要なのは、ある整数 $j \geq 0$ が存在し、
$\Ker(F)$ と $\Coker(F)$ が $p^j$ によって消されることである。
$\mathcal{E}$ の階数が有界なら、例えば $X$ が準コンパクトなら、
これら二つの条件はいずれも定義で述べた非退化条件から従う。
実際、$R$ を環、$r \geq 1$ を整数とし、
$K, L \in \text{Mat}(r \times r, R)$ は
$K L = p^i 1_{r \times r}$ を満たす行列であるとする。このとき
$\det(K)\det(L) = p^{ri}$ である。$L'$ を $L$ の余因子行列、すなわち
$L' L = L L' = \det(L)$ を満たす行列とする。
$K' = p^{ri} K$ および $j = ri + i$ とおく。
$K L = p^i$ なので $K' L = p^j 1_{r \times r}$ であり、
$$
L K' = L K \det(L) \det(M) = L K L L' \det(M) = L p^i L' \det(M) =
p^j 1_{r \times r}
$$
でもある。したがって $V$ が定義 \ref{crystalline-definition-F-crystal} のとおりなら、
$N > i \cdot \text{rank}(\mathcal{E})$ を満たす整数に対して
$V' = p^N V$ とおくことで、
$V' \circ F = p^{N + i}$ および $F \circ V' = p^{N + i}$ を得る。
\end{remark}

\begin{theorem}
\label{crystalline-theorem-cohomology-F-crystal}
状況 \ref{crystalline-situation-F-crystal} において、$(\mathcal{E}, F_\mathcal{E})$ を
非退化 $F$-クリスタルとする。$A$ は $p$ 進完備 Noether 環であり、
$X \to S_0$ は固有かつ滑らかであると仮定する。このとき標準写像
$$
F_\mathcal{E} \circ (F_X)_{\text{cris}}^* :
R\Gamma(\text{Cris}(X/S), \mathcal{E}) \otimes^\mathbf{L}_{A, \sigma} A
\longrightarrow
R\Gamma(\text{Cris}(X/S), \mathcal{E})
$$
は $p$ を可逆にすると同型になる。
\end{theorem}

\begin{proof}
まずこの射を三つの射の合成として書く。すなわち
$$
X^{(1)} = X \times_{S_0, F_{S_0}} S_0
$$
とおき、相対 Frobenius 射を $F_{X/S_0} : X \to X^{(1)}$ と書く。
$\mathcal{E}^{(1)}$ を $\Spec(\sigma)$ による $\mathcal{E}$ の基底変換、
言い換えれば可換図式
$$
\xymatrix{
X^{(1)} \ar[r] \ar[d] & X \ar[d] \\
S \ar[r]^{\Spec(\sigma)} & S
}
$$
に付随する結晶トポスの射による $\mathcal{E}$ の
$\text{Cris}(X^{(1)}/S)$ への引き戻しとする。このとき基底変換写像
\begin{equation}
\label{crystalline-equation-base-change-sigma}
R\Gamma(\text{Cris}(X/S), \mathcal{E}) \otimes^\mathbf{L}_{A, \sigma} A
\longrightarrow
R\Gamma(\text{Cris}(X^{(1)}/S), \mathcal{E}^{(1)})
\end{equation}
がある。注意 \ref{crystalline-remark-base-change} を参照されたい。
$F_{X/S_0} : X \to X^{(1)}$ と射影 $X^{(1)} \to X$ の合成は
絶対 Frobenius 射 $F_X$ であることに注意する。したがって
$F_{X/S_0}^*\mathcal{E}^{(1)} = (F_X)_{\text{cris}}^*\mathcal{E}$ である。
ゆえに $F_{X/S_0}$ による引き戻しは写像
\begin{equation}
\label{crystalline-equation-to-prove}
F_{X/S_0}^* :
R\Gamma(\text{Cris}(X^{(1)}/S), \mathcal{E}^{(1)})
\longrightarrow
R\Gamma(\text{Cris}(X/S), (F_X)^*_{\text{cris}}\mathcal{E})
\end{equation}
である。最後に $F_\mathcal{E}$ を用いて写像
\begin{equation}
\label{crystalline-equation-F-E}
R\Gamma(\text{Cris}(X/S), (F_X)^*_{\text{cris}}\mathcal{E})
\longrightarrow
R\Gamma(\text{Cris}(X/S), \mathcal{E})
\end{equation}
を得る。定理の写像は、上の三つの写像
(\ref{crystalline-equation-base-change-sigma}), (\ref{crystalline-equation-to-prove}), および
(\ref{crystalline-equation-F-E}) の合成である。注意
\ref{crystalline-remark-base-change-isomorphism} により、第一の写像は $p$ の
すべての冪を法として擬同型である。関係する複体は $D(A)$ で完全なので、
これは擬同型である。注意 \ref{crystalline-remark-complete-perfect} を参照されたい。
$F_\mathcal{E}$ は $p$ の冪を除けば逆写像をもつので、第三の写像は
$p$ を可逆にすると擬同型である。注意
\ref{crystalline-remark-F-crystal-variants} を参照されたい。最後に補題
\ref{crystalline-lemma-pullback-relative-frobenius} により、第二の写像は
$p$ を可逆にすると同型である。
\end{proof}
